\documentclass[11pt]{article}

\usepackage[T1]{fontenc}
\usepackage[utf8]{inputenc}
\usepackage{lmodern}
\usepackage[margin=1in]{geometry}
\usepackage{amsmath,amssymb,amsthm,mathtools}
\usepackage{enumitem}
\usepackage{booktabs}
\usepackage{array}
\usepackage{longtable}
\usepackage{xcolor}
\usepackage{microtype}
\usepackage{etoolbox}
\usepackage{url}
\usepackage[colorlinks=true,linkcolor=blue!50!black,citecolor=green!35!black,urlcolor=blue!60!black]{hyperref}
\usepackage[capitalize,nameinlink]{cleveref}

% ---------------------------------------------------------------------------
% Theorem environments
% ---------------------------------------------------------------------------
\theoremstyle{plain}
\newtheorem{theorem}{Theorem}[section]
\newtheorem{lemma}[theorem]{Lemma}
\newtheorem{corollary}[theorem]{Corollary}
\newtheorem{proposition}[theorem]{Proposition}
\newtheorem{conjecture}[theorem]{Conjecture}
\newtheorem{question}[theorem]{Question}
\newtheorem{problem}[theorem]{Problem}
\theoremstyle{definition}
\newtheorem{definition}[theorem]{Definition}
\newtheorem{remark}[theorem]{Remark}
\newtheorem{example}[theorem]{Example}

% The Lean proof is the proof; the text gives a faithful outline of it.
\newenvironment{outline}[1][Outline of the formal proof]{\begin{proof}[#1]}{\end{proof}}
% Results that are not formalized get complete proofs, labelled as such.
\newenvironment{informalproof}{\begin{proof}[Proof (not formalized)]}{\end{proof}}

% ---------------------------------------------------------------------------
% Lean references
%   \lean{name}            a declaration, relative to the namespace MotionDepthCodes
%   \formal{file}{names}   the formal counterpart of a numbered result: a file relative to
%                          lean/MotionDepthCodes/ and a comma-separated list of declarations
%   \partformal{file}{names}{gap}
%                          a numbered result whose formal counterpart covers all but the
%                          part described in {gap}; the gap is proved or cited in the text
%   \notformalized         a numbered result that has no formal counterpart
%   \cited{where}          a result of an earlier paper used without proof: where it is proved,
%                          and whether it is Lean-verified there
% ---------------------------------------------------------------------------
% Lean names are typeset like paths, with line breaks allowed after "_", "." and "/".
\DeclareUrlCommand\leanname{\def\UrlFont{\ttfamily\upshape}}
\renewcommand{\UrlBreaks}{\do\_\do\.\do\/\do\-}
\newcommand{\lean}[1]{\leanname{#1}}
\newcommand{\leanfile}[1]{\leanname{#1}}
\newcommand{\leanlist}[1]{\begingroup\def\leansep{}%
  \renewcommand*{\do}[1]{\leansep\leanname{##1}\def\leansep{, }}\docsvlist{#1}\endgroup}
\newcommand{\formal}[2]{\par\noindent{\small\upshape\raggedright\textsf{Lean:} \leanlist{#2}\unskip
  \nobreak\hfil\penalty50\hskip1em\hbox{}\nobreak\hfill[\leanfile{#1}]\parfillskip=0pt\par}%
  \smallskip\noindent\ignorespaces}
\newcommand{\partformal}[3]{\par\noindent{\small\upshape\raggedright\textsf{Lean (partial):} \leanlist{#2}\unskip
  \nobreak\hfil\penalty50\hskip1em\hbox{}\nobreak\hfill[\leanfile{#1}]\parfillskip=0pt\par}%
  \par\noindent{\small\upshape\textsf{Not formalized:} #3}\par\smallskip\noindent\ignorespaces}
\newcommand{\notformalized}{\par\noindent{\small\upshape\textsf{Not formalized;} a complete proof is given
  below.}\par\smallskip\noindent\ignorespaces}
\newcommand{\cited}[1]{\par\noindent{\small\upshape\textsf{Cited:} #1}\par\smallskip\noindent\ignorespaces}

% ---------------------------------------------------------------------------
% Notation
% ---------------------------------------------------------------------------
\newcommand{\N}{\mathbb{N}}
\newcommand{\Z}{\mathbb{Z}}
\newcommand{\Q}{\mathbb{Q}}
\newcommand{\R}{\mathbb{R}}
\newcommand{\C}{\mathbb{C}}
\newcommand{\F}{\mathbb{F}}
\newcommand{\T}{\mathsf{T}}
\newcommand{\clog}[1]{\lceil \log_2 #1 \rceil}
\newcommand{\abs}[1]{\lvert #1 \rvert}
\newcommand{\card}[1]{\lvert #1 \rvert}
\newcommand{\eps}{\varepsilon}
\newcommand{\prev}{\mathrm{prev}}
\newcommand{\cost}{\mathrm{cost}}
\newcommand{\gn}{\mathrm{gn}}
\newcommand{\Sh}{\mathrm{Sh}}
\newcommand{\cl}{\mathrm{cl}}
\newcommand{\Aut}{\mathrm{Aut}}
\newcommand{\minrk}{\mathrm{minrk}}
\newcommand{\MAIS}{\mathrm{MAIS}}
\newcommand{\Crux}{\mathrm{Crux}}
\newcommand{\TCP}{\mathrm{TCP}}
\DeclareMathOperator{\rank}{rank}
\DeclareMathOperator{\dist}{dist}
\DeclareMathOperator{\poly}{poly}


\title{Curved Motion, Depth Reduction, and Storage Codes\\ for Matrix Transposition on Tapes}
\author{Hans Brende \and Claude Opus 5.5 (Anthropic)}
\date{October 2026}

\begin{document}

\maketitle

\begin{abstract}
Can a deterministic multitape Turing machine transpose a $d\times d$ matrix with $\lceil\log_2 d\rceil$-bit entries in time
linear in the input length? The question is open, and a negative answer would give the first superlinear lower bound for
off-line integer multiplication on such machines. We ask what a linear-time transposer would need, and we find three
ingredients: curved head motion; depth-robust dependencies, or extreme coding on shallow ones; and storage codes
that beat rate~$\tfrac12$.
We prove that no machine, oblivious or adaptive, whose work heads are all home-based with legs made of boundedly
many constant-rate runs transposes in linear time, and that the macroscopic head motion of every linear-time transposer
contains a linear amount of curved, not countably-piecewise affine, travel (Theorems~C and~Cq). We show that two statements together would
rule out oblivious linear-time transposers: depth reduction of the time--space dependency DAG by a cheap genie, and an
index-coding lower bound on the remaining shallow cones (Theorem~D). We prove depth reduction for every coupling of
home-based walks (Theorem~B), show that packing certificates cannot establish depth-robustness, and give an exact
accounting of label costs on the sharpest open instance, the deep-nesting pair (Lemma~Zq, Theorems~E and~E$_1$). For the index-coding side we study the dual notion, storage codes, on sparse graphs. Rate bounds pull back along
homomorphisms (Theorem~H), so every storage code, nonlinear included, on a graph that maps to $C_5$, to $K_{14/5}$ or to
the Petersen graph has rate at most $\tfrac12$, whatever its degree. The Clebsch graph has guessing number exactly~$10$, so bounded degree alone does not force rate~$\tfrac12$.
On triangle-free graphs of maximum degree~$3$, no linear storage code over any field beats rate~$\tfrac12$ (Theorem~S).
On locally bipartite graphs of bounded degree, linear codes with bounded complex coefficients have rate at most
$\tfrac12+o(1)$, so gains of finite-field linear codes with a fixed number of symbols per vertex need
$p$-torsion (Theorem~A). The entropy-based results on storage codes (Theorem~H, symmetrization, the good targets and the Clebsch graph),
Theorem~S and the exact label accounting (Lemma~Zq, the exact cases of the chain and transfer identities, and
Theorem~E) are formalized in Lean~4, apart from the classical Tutte--Berge formula
used in the matching form of Theorem~S; all other results are proved in full but not formalized. We end with a roadmap: positive answers to depth reduction for interacting walks and to the index-coding
question on shallow cones, which is open for finite-field and nonlinear codes even in its undirected form, would
rule out oblivious linear-time transposers,
and our restricted results say what a construction would have to combine.

\end{abstract}

\begin{quote}\small
\textbf{Contributions.} Hans Brende chose the problem, fixed the verification standard, and directed the research.
Claude Opus 5.5, a large language model developed by Anthropic, developed the mathematics, wrote the Lean
formalization, which uses a port of J.~P. Traverso Gianini's formalization of the Tutte--Berge
formula~\cite{TraversoGianini2026}, and wrote this paper under Hans Brende's direction.
\end{quote}

\tableofcontents
\clearpage

\section{Introduction}\label{sec:intro}

\subsection{The problem}

Can a deterministic multitape Turing machine transpose a $d\times d$ matrix with $b(d)=\max(1,\clog d)$-bit entries, written row
by row, in time linear in the input length $L=d^2b(d)$? This is the statement~(T) of the first paper of this
series~\cite{FirstPaper2026}. Transposition is the data-movement step of fast integer multiplication on tapes, and by a
reduction of Harvey and van der Hoeven~\cite[Cor.~6.9]{HvdH2025} a linear-time multiplication machine would give a
linear-time transposer. So a proof of $\neg$(T) would give the first superlinear lower bound for off-line integer
multiplication on multitape Turing machines, and a proof of~(T) would remove the data-movement obstacle to linear-time
multiplication.

The first paper proved, in Lean, that every transposer needs $\Omega(\log d)$ monotone passes and that the macroscopic
head motion of every linear-time transposer must fold back (Theorems~P, P$'$ and~M*), and it reduced $\neg$(T) to a network-coding statement on tape networks,
the Tape Coding Conjecture. The second paper~\cite{SecondPaper2026} excluded, in its models, local coding schemes with
low-treewidth interaction graphs and regular (piecewise affine) limiting motion, and the third~\cite{ThirdPaper2026} related bit and word
transposition. The question remains open, and this paper does not decide it.

\subsection{Three ingredients}

This paper asks what a linear-time transposer would need, and finds three ingredients that an oblivious one would have to
combine; the first is needed by adaptive machines as well.
\begin{itemize}[leftmargin=2em]
\item \emph{Curved, head-dependent motion.} Its work heads cannot all move along home-based legs made of boundedly many
constant-rate runs, and a linear amount of their macroscopic motion must be curved, in a way that no common change of
clock removes.
\item \emph{Depth-robust dependencies, or extreme coding on shallow ones.} Cut its time--space dependency DAG with a cheap
genie, which reveals a small set of intermediate values. Either the genie leaves large depth, or the shallow cones that remain must carry index codes saving all but an
$\eps$ fraction of the input.
\item \emph{Codes that beat the pairing.} On sparse side information that looks locally like a tree, such codes must beat
rate~$\tfrac12$, the rate of the code that pairs up receivers.
\end{itemize}
Our theorems prove restricted versions of each of these statements, and reduce the oblivious case of (T) to two
precisely stated open problems (Theorem~D): depth reduction for interacting walks, whose sharpest open instance has two
walks, and an index-coding lower bound on the side information of shallow cones, whose undirected form is open for
finite-field and nonlinear codes.

\subsection{Results}

\paragraph{Curved motion (\cref{sec:motion}).} \emph{Theorem~C.} No machine, oblivious or adaptive, whose work heads are
all home-based on every input, with legs made of boundedly many constant-rate runs, transposes in linear time
(\cref{thm:C}). The homes may depend on the input, and the number of legs, their scales and their interleaving are
arbitrary. \emph{Theorem~Cq}: in every macroscopic limit of a linear-time transposer, the work heads travel a total
distance of at least half a block unit, a linear number of cells, along motion that is not
countably-piecewise affine (\cref{thm:Cq}). Theorem~C follows from it, and it contains the second paper's affine-time
obstruction (\cref{cor:mot-lambda}). These results are unconditional; they use an extension of Theorem~M* from interval partitions to Borel partitions, and the
hyperfiniteness of amenable equivalence relations.

\paragraph{The decomposition (\cref{sec:decomposition}).} \emph{Theorem~D.} If the dependency DAG of every oblivious
linear-time run can be cut to bounded depth by revealing $\eps L$ bits (\emph{depth reduction}), and if index codes on the
side information of bounded-depth cones need a broadcast of $\rho L$ bits with $\rho>\eps$ (\emph{the crux}), then no
oblivious linear-time transposer exists (\cref{thm:D,cor:dec-obl}). The composition is elementary; its content is the formulation and
the bookkeeping, including the forced trade-off $\eps\ge\beta/(D+1)-\beta/L$ between the genie and the depth~$D$,
where $\beta$ is the number of bits revealed per step (\cref{prop:tradeoff}). \emph{Theorem~B} proves depth reduction for
every coupling of $k$ home-based walks of length~$T$: $\eps T$ cut edges leave depth $2^{O(k^3/\eps)}$, whatever the
walks do (\cref{thm:B}). So an oblivious transposer with home-based work heads
exists only if the crux fails, even in a weak form (Corollary~B, \cref{cor:B}).

\paragraph{Depth reduction (\cref{sec:depth}).} For walks that turn far from home, depth reduction is open (\cref{prob:tcp});
for one walk it holds by planarity. We prove that no packing certificate, whatever structures it packs, can
certify depth-robustness
(\cref{prop:packing} and \cref{lem:dep-euler}(c)), and we sketch why per-head band bookkeeping, a natural strategy for depth reduction, fails on
deeply nested walks (\cref{rem:dep-N,rem:dep-Nprime}; proof sketches only). On what we regard as the sharpest open instance, the \emph{deep-nesting pair}, we develop an exact accounting of label
costs: \emph{Lemma~Zq} expresses the cost of any label, a map from times to $\{0,\dots,q-1\}$ that describes a cut of depth
$q-1$, through residues modulo~$q$, and the
chain-decomposition and transfer identities describe labels that are exact on one head. \emph{Theorem~E}: every
hop-disjoint family of negative-defect cycles of head-$2$ hops forces such labels to pay about one violation per cycle
(\cref{thm:E}). \emph{Theorem~E$_1$}: on the golden deep-nesting pair with one level of nesting, labels exact on one head pay a fixed
price per hop on the cover at every depth, at least $67/4096$ at all but the smallest scale (\cref{thm:E1}). Whether the pairs are depth-robust as $m$ grows is open; we state one form of the
missing step, for labels on the cover and a suitable odd shift, as \cref{conj:Aprime}.

\paragraph{Storage codes (\cref{sec:codes}).} The crux is open in general. Its undirected form asks whether storage codes on
locally bipartite graphs of bounded degree can beat rate~$\tfrac12$ (\cref{prob:undirected-crux}).
\begin{itemize}[leftmargin=2em]
\item \emph{Theorem~H}: rate bounds pull back along homomorphisms of graphs and digraphs, for all codes
(\cref{thm:H}). With exact certificates for $C_5$, the Petersen graph and $K_{14/5}$, every storage code, nonlinear
included, on a graph that maps to one of them has rate at most $\tfrac12$, with no degree assumption
(\emph{Corollary~H2}, \cref{cor:H2}). On cubic graphs of large girth the nonlinear case follows from a positive answer to Nešetřil's
Pentagon Problem (\cref{cor:H3}).
\item \emph{Corollary~H4}: the Clebsch graph has guessing number exactly~$10$, and Shannon bound exactly~$5/8$
(\cref{cor:H4,prop:clebsch-exact}). This answers, for the Clebsch graph, a question of Cameron, Dang and Riis~\cite{CameronDangRiis2014}, and
shows that bounded degree alone does not force rate~$\tfrac12$.
\item \emph{Theorem~S}: on triangle-free graphs~$G$ of maximum degree~$3$ with $N$ vertices, every linear storage code over
every field has dimension at most $\lfloor N/2\rfloor$, and with the Tutte--Berge formula the minrank is $N-\nu(G)$, where
$\nu(G)$ is the matching number (\cref{thm:S,thm:S-barrier}).
\item \emph{Theorem~A}: on locally bipartite graphs of bounded degree, complex matrices with unit diagonal, graph support
and bounded entries have rank at least $(\tfrac12-o(1))N$; so over finite fields a gain of a linear code with a fixed
number of symbols per vertex must come from $p$-torsion (\cref{thm:A,cor:torsion}).
\end{itemize}
We also show limits of local and spectral arguments (\cref{sec:ana-limits}), and that Mycielski's construction does
not preserve goodness, the Shannon rate bound~$\tfrac12$ of \cref{def:polymatroid} (\cref{sec:ana-targets}).

\paragraph{Roadmap (\cref{sec:construction,sec:open}).} We list what any linear-time transposer must do, describe an
architecture compatible with these conditions, and state the open problems whose positive solution would rule out
oblivious linear-time transposers.

\subsection{How to read this paper: the formal proofs}

The results on storage codes of \cref{sec:entropic,sec:thmH,sec:symmetrization,sec:targets,sec:clebsch,sec:subcubic}
(Theorem~H, Corollaries H2--H4, Theorem~S) and the exact label accounting of \cref{sec:depth} (Lemma~Zq, the exact cases of
the chain and transfer identities, Theorem~E) are formalized in Lean. The only exception is the classical Tutte--Berge
formula, which the matching form of Theorem~S cites. All other results are proved in full but not
formalized, among them Theorems~D, B, C, Cq, E$_1$ and~A, \cref{prop:tradeoff,prop:packing}, the general forms of the
chain and transfer identities, and the results of \cref{sec:ana-A,sec:ana-fields,sec:ana-limits,sec:ana-targets}. Each of them comes with a
complete proof, and each numbered result states its status as described in \cref{sec:prelim-lean}. Claims with
only a proof sketch, among them the nesting barrier for per-head band rules (\cref{rem:dep-N,rem:dep-Nprime}), are stated
in remarks and are not used in proofs.

\subsection{Related work}

\paragraph{Depth reduction and index coding.} The route from a computation of small depth to index coding goes back to
Valiant's depth reduction for circuits~\cite{Valiant1977} and Riis's guessing games~\cite{Riis2007}, and was used, in the language of
min-rank, by Jukna and Schnitger~\cite{JuknaSchnitger2011}; a related network-coding view of oblivious external-memory
algorithms, and through them of oblivious two-tape machines, is due to Adler, Harvey, Jain, Kleinberg and
Lehman~\cite[Section~6, Theorems~6.2 and~6.3]{AHJKL2006}. On multitape machines, Paul, Pippenger, Szemerédi and Trotter separated deterministic
from nondeterministic linear time by a segregator argument on the dependency graph of a run~\cite{PPST1983}; depth
reduction at linear cost may delete far more, a small constant fraction of the edges, but must then leave constant
depth (\cref{sec:dep-problem}).

\paragraph{Guessing games and storage codes.} Guessing numbers were introduced by Riis~\cite{Riis2007} and studied for
undirected graphs by Christofides and Markström~\cite{ChristofidesMarkstrom2011}, Baber, Christofides, Dang, Vaughan and
Riis~\cite{Baber2017}, and Cameron, Dang and Riis~\cite{CameronDangRiis2016}. Storage codes on graphs and their duality
with index coding were studied by Mazumdar~\cite{Mazumdar2014} and by Arbabjolfaei and Kim~\cite{ArbabjolfaeiKim2015};
Mazumdar, McGregor and Vorotnikova bounded the storage capacity above by vertex covers and averages of vertex covers,
and below by fractional matchings~\cite{MMV2019}, and Barg and Zémor constructed linear storage codes of rate
approaching~$3/4$ on triangle-free graphs of growing degree~\cite{BargZemor2022}. Index coding was introduced by Birk and
Kol~\cite{BirkKol2006} and formalized, with minrank, by Bar-Yossef, Birk, Jayram and Kol~\cite{BBJK2011}; the Shannon
linear-programming bound for it is due to Blasiak, Kleinberg and Lubetzky~\cite{BKL2013}, and minrank goes back to
Haemers~\cite{Haemers1979}.

\paragraph{Motion.} Theorems~C and~Cq rest on the first paper's Theorem~M* and on the theorem of Connes, Feldman and
Weiss that amenable nonsingular equivalence relations are hyperfinite~\cite{CFW1981}; the Lebesgue-measure form of
Theorem~Cq for relatively open affine pieces (\cref{cor:mot-lambda}) slightly generalizes the second paper's
affine-time obstruction~\cite{SecondPaper2026}. With one work tape, much stronger lower bounds are classical:
$\Omega(L^{3/2}/(\log L)^{1/2})$ steps without an output tape~\cite{DMS1991}, and $\Theta(L^{5/4})$ for Boolean matrices
with a write-only output tape~\cite{DietzfelbingerMaass1993}. Our motion results hold for any number of tapes, and they
leave room for curved motion such as the monotone expanders of Bourgain and Yehudayoff~\cite{BourgainYehudayoff2013}, on
which the first paper's architecture is built.

\subsection{Organization}

\Cref{sec:prelim} fixes notation. \Cref{sec:motion} proves the motion theorems, \cref{sec:decomposition} the
decomposition and depth reduction for home-based walks, \cref{sec:depth} the results on depth reduction for interacting
walks, and \cref{sec:codes} the results on storage codes. \Cref{sec:construction} collects the consequences for constructions, and
\cref{sec:open} the roadmap and open problems. The appendices index the Lean declarations and explain how to build and
check the library.

\section{Preliminaries}\label{sec:prelim}

\subsection{Machines and the statement (T)}\label{sec:prelim-machines}

We use the machine model of the first paper~\cite{FirstPaper2026}: a deterministic Turing machine with a finite state
set, a read-only two-way input tape, $k\ge1$ two-way work tapes over a finite alphabet, and a write-only one-way output
tape that receives at most one bit per step (the first paper also allows $k=0$; an idle work tape can always be added). All heads start at cell~$0$; head moves are $-1$, $0$ or $+1$; every step
writes the cell under each work head. For $d\ge1$ let $b(d)=\max(1,\clog d)$. The input is a $d\times d$ matrix with
$b(d)$-bit entries written row by row, of length $L=d^2b(d)$, and the machine must write its transpose in the same
encoding. Statement~(T) says that some machine does this within $c(L+1)$ steps for a constant~$c$ and all~$d$. Its
formal version, \lean{Transpose.LinearTimeTransposable}, is fixed in a frozen file of the first paper's library.

A machine is \emph{oblivious} if, for each input length, the positions of all heads at every time, the times at which it
appends output bits, and its halting time do not depend on the input. An oblivious machine has a single head schedule
per input length; an adaptive machine has one per input.

The Lean library of this paper does not model machines. Its theorems are about polymatroids, storage codes, matrices, and
labels on finite sets with hop bijections; the machine-level statements of \cref{sec:decomposition,sec:motion} are not formalized, and each comes
with a complete proof.

\subsection{Walks, hops and legs}\label{sec:prelim-walks}

A \emph{walk} of \emph{length}~$T$ is a map $x\colon\{0,\dots,T-1\}\to\Z$ with $\abs{x(t+1)-x(t)}\le1$. The work heads of a machine trace
walks $x_1,\dots,x_k$, one per work tape, and in a run these walks may be coupled arbitrarily.

For a head $h$ and a time $u$, let $\prev_h(u)=\max\{t<u : x_h(t)=x_h(u)\}$ when this set is nonempty. The pair
$(\prev_h(u),u)$ is the \emph{hop} of head $h$ ending at $u$: the symbol that head $h$ reads at step $u$ was written at
step $\prev_h(u)$, because every step writes the cell under each work head. A hop is \emph{long} if
$u-\prev_h(u)\ge2$; the other hops are pauses. The \emph{legs} of a walk partition its time interval greedily: the first
leg starts at time~$0$, each leg is the longest interval beginning at its first time on which the walk is weakly
monotone, and the next leg starts at the following time. So a leg ends just before the first step that reverses its direction.
A walk is \emph{home-based}, with home cell~$0$,
if every leg begins or ends within distance~$1$ of~$0$. Theorems about head motion are invariant under shifting a
tape, so the home cell can be any cell.

\subsection{Storage codes and guessing numbers}\label{sec:prelim-codes}

\emph{Side information} on a finite vertex set $V$ is a map $N\colon V\to 2^V$. For a simple graph $G$ on $V$ it is the
neighbourhood map $N(v)=\{u:uv\in E(G)\}$ (\lean{nbhd}); for a digraph it is the out-neighbourhood map
(\lean{outNbhd}). Let $\alpha$ be a finite alphabet.

\begin{definition}[Storage codes]\label{def:storage-code}
\formal{StorageCode.lean}{IsStorageCode, IsRecoverable}
A set $C\subseteq\alpha^V$ is a \emph{storage code} for $N$ if for all $x,y\in C$ and every $v\in V$, $x$ and $y$ agree at
$v$ whenever they agree on $N(v)$. A random word $X\colon\Omega\to\alpha^V$ is \emph{recoverable} for $N$ if the same holds
for all pairs of its values.
\end{definition}

So in a storage code each symbol is a function of the symbols at its neighbours. The \emph{rate} of $C$ is
$\log\card C/(\card V\log\card\alpha)$. The \emph{guessing number} $\gn(G,q)$ is the largest value of
$\log_q\card C$ over storage codes on $G$ with $\card\alpha=q$, and the asymptotic guessing number is
$\gn(G)=\sup_q\gn(G,q)$~\cite[Def.~3.7 and Thm.~3.6]{ChristofidesMarkstrom2011}; guessing numbers for a fixed alphabet go
back to Riis~\cite{Riis2007}. Storage codes are dual to index codes: for the
same side information, the optimal rates of the two problems are complementary asymptotically, and exactly for linear
codes~\cite{Mazumdar2014,ArbabjolfaeiKim2015}.
We work with storage codes throughout.
A \emph{linear} storage code over a field $F$ is a subspace of $F^V$ with the defining property; equivalently
(\cref{lem:linear-pattern}), a subspace of the kernel of a matrix $M\in F^{V\times V}$ with unit diagonal whose off-diagonal
entries vanish outside the edges of~$G$. The largest dimension of such a kernel is $\card V$ minus the
\emph{minrank} of $G$ over $F$~\cite{Haemers1979,BBJK2011}.

\subsection{Polymatroids and rate bounds}\label{sec:prelim-polymatroids}

\begin{definition}[Polymatroids, validity, rate bounds]\label{def:polymatroid}
\formal{Polymatroid.lean}{IsPolymatroid, IsValid, RateBound, IsGood}
A \emph{polymatroid} on $V$ is a function $h\colon 2^V\to\R$ with $h(\emptyset)=0$ that is monotone and submodular:
$h(A\cup B)+h(A\cap B)\le h(A)+h(B)$. It is \emph{valid} for side information $N$ if $h(N(v)\cup\{v\})=h(N(v))$ for every
$v$. Side information $N$ satisfies the \emph{rate bound} $c$ if every valid polymatroid satisfies
\[
  h(V)\;\le\;c\sum_{v\in V}h(\{v\}).
\]
A graph $G$ is \emph{good} if its neighbourhood map satisfies the rate bound~$\tfrac12$, that is, $2h(V)\le\sum_vh(\{v\})$ for
every valid polymatroid.
\end{definition}

The rate bound quantifies over all polymatroids, so it is a statement about the \emph{Shannon} information
inequalities; the smallest admissible $c$ is the Shannon bound $\lambda_{\Sh}$ of the side information.
For a random word $X$ on a finite probability space, $S\mapsto H(X_S)$ is a polymatroid (\cref{lem:entropic}),
and it is valid whenever $X$ is recoverable (the converse holds when every outcome has positive probability). Hence a rate bound $c$ bounds every storage code, linear or not, over
every alphabet: $H(X_V)\le c\sum_vH(X_v)$ for every recoverable $X$, and $\log\card C\le c\,\card V\log\card\alpha$
for every storage code $C$ (\cref{thm:entropic-rate}). Entropy is the finite Shannon entropy of the second
paper~\cite{SecondPaper2026}, whose Lean definitions we import.

\subsection{Formal status of the results}\label{sec:prelim-lean}

Every numbered result names its formal status in one of four ways. A line
\begin{center}
{\small\textsf{Lean:} \texttt{declaration, \dots}\hfill[\texttt{file}]}
\end{center}
gives the Lean declarations of a fully formalized result. The declarations live in the namespace
\texttt{MotionDepthCodes} and the files in the directory \leanfile{lean/MotionDepthCodes/}. The text after such a statement, labelled \emph{Outline of
the formal proof}, is a guide to the formal proof; it is a convenience, not a proof. A line beginning \textsf{Lean
(partial)} marks a result whose formal counterpart covers all but an explicitly named part, which the text proves or
cites. The words \textsf{Not formalized} mark a result with no formal counterpart; it comes with a complete proof,
labelled \emph{Proof (not formalized)}. A line beginning \textsf{Cited} marks a result of an earlier paper that is used
without proof, and says where it is proved and whether it is Lean-verified there. Where the English statement
simplifies, the Lean statement is authoritative.
Results from the literature are cited; a few short ones are reproved in our notation, and say so. Numerical
observations are labelled as such, claims that have only a proof sketch are stated in remarks beginning \emph{Proof sketch
only}, and neither is used in proofs. Numbered definitions carry a status line only when they
are formalized.

The library depends on Lean~4 and Mathlib (both v4.34.1) and on the Lean libraries of the first two papers of this
series, imported through relative paths. It contains no \texttt{sorry}, no \texttt{native\_decide} and no unsafe
evaluation, and every declaration cited here depends only on the axioms \texttt{propext}, \texttt{Classical.choice} and
\texttt{Quot.sound} (\cref{app:verify}). \Cref{app:index} lists every declaration of this library cited in the paper.

\section{Curved motion}\label{sec:motion}

This section constrains the head motion of every linear-time transposer, oblivious or adaptive; the decomposition of
the oblivious case follows in \cref{sec:decomposition}. The first paper~\cite{FirstPaper2026} rescales the work-head trajectories of a machine, passes to uniform limits
along subsequences of sizes, and proves that the limits of a linear-time transposer satisfy an isoperimetric
inequality $(\star)$. This is Theorem~M*, restated as \cref{thm:Mstar}. In this section we turn $(\star)$ into a
constraint on the shape of the motion. We combine it with the second paper's labels for finitely many affine pieces or
affine maps, built from F\o lner sets of finitely many affine maps and Lean-verified there~\cite{SecondPaper2026}, and
with the area formula, which charges the remaining returns of the heads by their speed (\cref{sec:mot-tools}). The outcome is Theorem~Cq (\cref{thm:Cq}): in every limit of a linear-time transposer, the work heads travel a total
distance of at least $\frac12$, in block units, along motion that is not countably-piecewise affine. Theorem~C
(\cref{thm:C}) applies this to machines whose work heads are all home-based, with legs made of boundedly many runs at
constant rate: no such machine, oblivious or adaptive, transposes in linear time.
\Cref{sec:mot-consequences} draws consequences. One work tape never suffices. Theorem~Cr (\cref{thm:Cr}) passes from
affine pieces to affine returns: in every common clock, that is, every reparametrization of time by a homeomorphism
whose inverse maps null sets to null sets, and for every countable set of affine maps, the times at which a head's next
return to its level is not given by one of these maps carry at least one block unit of travel in total; so a common
change of clock and fixed changes of scale do not help. Machines whose work heads all follow stroke trees are excluded, and a construction
must use head-dependent curved motion on at least two interacting tapes (\cref{cor:mot-design}).

These are restricted impossibility results, and they are unconditional: they use neither depth reduction
(\cref{prob:tcp}) nor any statement about codes. On the construction side they say which motion a linear-time
transposer must have. No numbered result of this section is formalized in this paper's Lean library. Theorem~M* and the finite affine
labels, on which they rest, are Lean-verified in the first two papers; the remaining cited ingredient, the area formula
for Lipschitz functions, is classical, as are the other facts of measure theory used below. Amenability enters only
through F\o lner sets of finitely many affine maps, so no hyperfiniteness theorem is needed.

\subsection{Rescaled limits and Theorem M*}\label{sec:mot-setting}

In this section the letter $t$ indexes the work heads, from $0$ as in the first paper, and $h$ denotes a home cell.
Let $M$ be a machine with $k\ge1$ work tapes whose work alphabet has $g+1$ symbols, and let $G=\max(1,\clog{(g+1)})$
be the number of bits per work cell~\cite[Definition~6.7]{FirstPaper2026}. For $d\ge1$ put $r_d=\lceil d\,b(d)/G\rceil$
and $R_d=d\,r_d$, so that $L(d)/G\le R_d\le L(d)/G+d$. Time is cut into \emph{slots} of $r_d$ steps and each work tape
into \emph{blocks} of $r_d$ cells; a block holds about one row of the matrix. On an input $w$ of size $d$ let
$x_{t,w}(s)$ be the cell of work head~$t$ after $s$ steps, frozen once $M$ halts. The \emph{rescaled trajectory} of
head~$t$ is
\[
  P^{(d)}_{t,w}(u)=\frac1d\Bigl\lfloor\frac{x_{t,w}(\lfloor ud\rfloor r_d)}{r_d}\Bigr\rfloor\qquad(u\ge0),
\]
the block of the head at the first step of slot $\lfloor ud\rfloor$, in units of $d$ blocks~\cite[Section~7]{FirstPaper2026}.
Macroscopic time and space are thus both measured in units of $R_d$ steps or cells. We fix a time constant $c>0$ and
put $c_0=cG+1$; a run of at most $cL(d)$ steps ends by macroscopic time $cL(d)/R_d\le cG<c_0$.

A \emph{trace class} of $M$ at size $d$~\cite[Definition~6.1]{FirstPaper2026} is a set $F$ of matrix inputs of size $d$
with a \emph{representative} $w\in F$ such that all members of $F$ halt after the same number of steps and have the
same block-level trace as $w$: slot by slot, the same block residencies of the work heads, the same input rows touched
and the same output length. Its \emph{co-entropy} is the least $D$ with $2^{L(d)}\le2^D\card F$. If $M$ is oblivious,
all matrix inputs of size $d$ form one trace class, of co-entropy~$0$~\cite[Lemma~6.2]{FirstPaper2026}.

\begin{definition}[Limits of a machine]\label{def:mot-limit}
A family $(P_t)_{t<k}$ of functions $[0,c_0]\to\R$ is a \emph{limit of $M$ at time constant $c$} if there are sizes
$d_i\to\infty$ and, for all large $i$, trace classes $F_i$ of $M$ at size $d_i$ with representatives $w_i$, common
halting times $T_i\le cL(d_i)$ and co-entropies $D_i$ with $D_i/L(d_i)\to0$, such that $M$ outputs the transpose on
every input in $F_i$ and $P^{(d_i)}_{t,w_i}\to P_t$ uniformly on $[0,c_0]$ for every $t<k$.
\end{definition}

For a limit, and more generally for any family of $1$-Lipschitz functions $P_t$ on $[0,c_0]$, we write $\lambda$ for
Lebesgue measure and
\[
  \nu_t=\abs{P_t'(u)}\,\mathrm du\ \text{ on }[0,c_0],\qquad
  \Lambda_t(p)=\{u\in[0,c_0]:P_t(u)=p\}.
\]
The derivative $P_t'$ exists almost everywhere and $\abs{P_t'}\le1$, so $\nu_t\le\lambda$. We call $\nu_t$ the
\emph{speed measure} of head~$t$: $\nu_t(A)$ is the distance travelled by the limit path during the times in $A$.
Distance $1$, which we call a \emph{block unit}, is $d$ blocks, that is, $R_d\approx L(d)/G$ cells.

The inequality $(\star)$ of the first paper uses labellings by finitely many intervals. An \emph{excised set} is a
finite union $S$ of closed intervals. A \emph{labelling} is a map $\beta\colon\R\to\N$ that is \emph{piecewise constant
off $S$}: for some finite set $Z$, $\beta$ is constant on every interval $[u,v]\subseteq[0,c_0]\setminus(S\cup Z)$. Its
\emph{classes off $S$} are the sets $\{u\in[0,c_0]\setminus S:\beta(u)=\ell\}$, $\ell\in\N$. Two points $s<s'$ of
$\Lambda_t(p)$ are \emph{consecutive} if no point of $\Lambda_t(p)$ lies strictly between them. Let $\kappa_t(p)$ be the
number of consecutive pairs $s<s'$ of $\Lambda_t(p)$ with $s,s'\notin S$ and $\beta(s)\ne\beta(s')$ if $\Lambda_t(p)$ is
finite, and $\kappa_t(p)=0$ if $\Lambda_t(p)$ is infinite. For $\theta\in[0,\infty]$ the inequality reads
\[
  (\star)\qquad 1\;\le\;\sum_{t<k}\int\kappa_t(p)\,\mathrm dp\;+\;2\theta\;+\;2k\,\lambda(S\cap[0,c_0]),
\]
where $\int$ denotes the lower Lebesgue integral, which is defined and monotone for every function with values in
$[0,\infty]$, measurable or not. We say that $(P_t)_{t<k}$ \emph{satisfies $(\star)_\emptyset$} if $(\star)$ holds with
$S=\emptyset$ for every $\theta\in(0,\frac12)$ and every labelling whose classes have measure at most $\theta$.

\begin{theorem}[Theorem M*~\cite{FirstPaper2026}]\label{thm:Mstar}
\cited{proved and Lean-verified in~\cite{FirstPaper2026}; not reproved here.}
Let $M$ be a machine with $k\ge1$ work tapes and let $c>0$.
\begin{enumerate}[label=(\alph*)]
\item Every limit $(P_t)_{t<k}$ of $M$ at time constant $c$ consists of $1$-Lipschitz functions on $[0,c_0]$ and
satisfies $(\star)$ for every $\theta\in[0,\infty]$, every excised set $S$ and every labelling, piecewise
constant off $S$, whose classes off $S$ have measure at most $\theta$. In particular it satisfies $(\star)_\emptyset$.
\item Let $d_i\to\infty$, and suppose that for all large $i$ the machine $M$ transposes every matrix of size $d_i$ and
halts within $cL(d_i)$ steps on it. Then $M$ has a limit at time constant $c$ along a subsequence of $(d_i)$, with trace
classes of co-entropy at most $C\,d_i\,b(d_i)$ for a constant $C$ that depends only on $M$ and $c$.
\end{enumerate}
\end{theorem}

This combines Lemma~7.1(a), Theorems~7.8 and~7.9 and the popular trace theorem (Theorem~6.8) of~\cite{FirstPaper2026};
the Lean declarations there are \leanname{Transpose.Tape.theoremMStar_level_class_on} for~(a) (with
\leanname{Transpose.Tape.blockTraj_limit_lipschitz} for the Lipschitz property) and
\leanname{Transpose.Tape.theoremMStar_level_adaptive} for~(b). The classes in~(b) are the \emph{popular} trace classes;
for an oblivious machine one can take all matrix inputs. A machine that witnesses~(T) with constant $c$ transposes
within $2cL(d)$ steps on every matrix input of every size, because $c(L+1)\le2cL$ (see the proof of Corollary~7.11
of~\cite{FirstPaper2026}; Lean: \leanname{Transpose.Tape.adaptive_hyps_of_linearTime}), so by~(b) it has limits at time
constant $2c$ along a subsequence of every sequence of sizes. We use only the case $S=\emptyset$ of~(a).

\subsection{Levels, returns and finite affine labels}\label{sec:mot-tools}

Statements about $1$-Lipschitz functions, here and below, hold for every $c_0>0$; for walks and limits of machines,
$c_0=cG+1$. In this subsection $P$ is a $1$-Lipschitz function on $[0,c_0]$, $\nu=\abs{P'}\,\mathrm du$ and
$\Lambda(p)=\{u\in[0,c_0]:P(u)=p\}$; for a family $(P_t)$ we add the index $t$ to everything.

\begin{lemma}[Levels]\label{lem:mot-levels}
\notformalized
\begin{enumerate}[label=(\alph*)]
\item For every Borel set $A\subseteq[0,c_0]$ the function $p\mapsto\#(A\cap\Lambda(p))\in\{0,1,\dots,\infty\}$ is
Lebesgue measurable, and $\int_\R\#(A\cap\Lambda(p))\,\mathrm dp=\nu(A)$.
\item $\Lambda(p)$ is finite for almost every $p$, and $\nu(\Lambda(p))=0$ for every $p$.
\item $\nu(A)\le\lambda(A)$ for every Borel set $A\subseteq[0,c_0]$; in particular $\lambda$-null sets are $\nu$-null.
\end{enumerate}
\end{lemma}

\begin{informalproof}
(a) Extend $P$ to a $1$-Lipschitz function on $\R$ that is constant on each side of $[0,c_0]$. The area formula for
Lipschitz maps from $\R$ to $\R$~\cite[Theorem~3.2.3]{Federer1969} states that for every Lebesgue measurable set
$A\subseteq\R$ the multiplicity function $p\mapsto\#(A\cap P^{-1}(p))$ is measurable and has integral $\int_A\abs{P'}$.
For $A\subseteq[0,c_0]$ the multiplicity is $\#(A\cap\Lambda(p))$. (For intervals $A$ this is Banach's indicatrix
theorem~\cite{Banach1925}.)
(b) For $A=[0,c_0]$, (a) gives $\int\#\Lambda(p)\,\mathrm dp=\nu([0,c_0])\le c_0<\infty$, so $\#\Lambda(p)<\infty$ for
almost every $p$. For $A=\Lambda(p_0)$, which is closed, the integrand in (a) vanishes for $p\ne p_0$, so
$\nu(\Lambda(p_0))=0$.
(c) $\abs{P'}\le1$ almost everywhere.
\end{informalproof}

All labellings below are piecewise constant off $S=\emptyset$, and $\kappa$ is the crossing count of $(\star)$ with
$S=\emptyset$. For a set $R\subseteq[0,c_0]$ of \emph{retained} times, let $\kappa^R(p)$ be the number of consecutive
pairs $s<s'$ of $\Lambda(p)$ with $s,s'\in R$ and $\beta(s)\ne\beta(s')$ if $\Lambda(p)$ is finite, and $\kappa^R(p)=0$
otherwise; so $\kappa=\kappa^{[0,c_0]}$. The points of $\Lambda(p)$ outside $R$ are the \emph{omitted} visits.

\begin{lemma}[Deleting visits]\label{lem:mot-delete}
\notformalized
For every labelling $\beta$, every set $R\subseteq[0,c_0]$ and every level $p$,
$\kappa(p)\le\kappa^R(p)+2\,\#(\Lambda(p)\setminus R)$. If $R$ is Borel, then
\[
  \int\kappa\,\mathrm dp\;\le\;\int\kappa^R\,\mathrm dp+2\,\nu([0,c_0]\setminus R).
\]
\end{lemma}

\begin{informalproof}
If $\Lambda(p)$ is infinite, then $\kappa(p)=0$. Otherwise every consecutive pair counted by $\kappa(p)$ but not by
$\kappa^R(p)$ has a point in $\Lambda(p)\setminus R$, and each point of $\Lambda(p)$ lies in at most two consecutive
pairs. For Borel $R$ the function $p\mapsto\#(\Lambda(p)\setminus R)$ is measurable with integral
$\nu([0,c_0]\setminus R)$, by \cref{lem:mot-levels}(a), and the lower integral of a sum is the sum of the lower integrals
when one summand, here $2\,\#(\Lambda(p)\setminus R)$, is measurable.
\end{informalproof}

The pointwise inequality is \leanname{TransposeFollowup.AffineMotion.levelCross_delete_le} in the second paper's Lean
library. The next lemma is the finite step of that paper's proof of its affine-time obstruction. For finitely many affine
pieces, the visits of a level in two pieces on which the path is not constant are related by an affine map of time; the
finitely many levels of the constant pieces form a null set of levels. The labels are built from F\o lner sets for the
finitely many affine maps that arise; such sets exist because the affine group is solvable, and the second paper
constructs them directly.

\begin{lemma}[Finite affine labels]\label{lem:mot-finite-labels}
\cited{the finite step of the proof of~\cite[Theorem~6.2]{SecondPaper2026}, which applies it to open pieces.
Lean-verified there for measurable pieces as
\leanname{TransposeFollowup.AffineMotion.exists_finite_affine_branch_small_cut}; the countable step of that proof,
\leanname{TransposeFollowup.AffineMotion.exists_countable_affine_path_cut_bound}, uses it. Not reproved here.}
Let $(P_t)_{t<k}$ be arbitrary functions on $[0,c_0]$, with $\Lambda_t(p)$, consecutive pairs and $\kappa_t^{R}$ defined
as above (no continuity is needed), and for each $t$ let $J_{t,1},\dots,J_{t,n}$ be Borel subsets of
$[0,c_0]$ on each of which $P_t$ agrees with an affine function; put $R_t=J_{t,1}\cup\dots\cup J_{t,n}$. For all
$\theta>0$ and $\varepsilon>0$ there is a labelling $\beta$ whose classes have measure at most $\theta$ and
$\sum_{t<k}\int\kappa_t^{R_t}\,\mathrm dp<\varepsilon$.
\end{lemma}

In the Lean statement the pieces need only be measurable, they may overlap and may carry constant functions, and the
$P_t$ need not be continuous. At each level it counts the crossing pairs with both points retained, and this count is at
least $\kappa_t^{R_t}(p)$.

The second form of the labels concerns first returns. For $u\in[0,c_0]$ let
$\psi(u)=\inf\{v\in[0,c_0]:v>u,\ P(v)=P(u)\}$, the \emph{first return} of the path to the level of $u$, defined on the
set $\operatorname{dom}\psi$ of those $u$ for which this set is nonempty, and let
$D^+=\{u\in\operatorname{dom}\psi:\psi(u)>u\}$. For a labelling $\beta$, the \emph{crossing set}
$A(\beta)=\{u\in D^+:\beta(\psi(u))\ne\beta(u)\}$ consists of the times whose first return changes the label. For a
family $(P_t)$ we write $\psi_t$, $D^+_t$ and $A_t(\beta)$.

\begin{lemma}[Returns]\label{lem:mot-returns}
\notformalized
\begin{enumerate}[label=(\alph*)]
\item $\operatorname{dom}\psi$ and $D^+$ are Borel sets, and $\psi\colon\operatorname{dom}\psi\to[0,c_0]$ is a Borel
map.
\item $P(\psi(u))=P(u)$ for every $u\in\operatorname{dom}\psi$, and $\psi$ is injective on $D^+$.
\item $\nu(\psi^{-1}(W)\cap D^+)\le\nu(W)$ for every Borel set $W\subseteq[0,c_0]$.
\item If $\Lambda(p)$ is finite, its consecutive pairs are exactly the pairs $(s,\psi(s))$ with $s\in\Lambda(p)\cap D^+$.
\item For every labelling $\beta$ the crossing set $A(\beta)$ is Borel, and $\int\kappa\,\mathrm dp\le\nu(A(\beta))$.
\end{enumerate}
\end{lemma}

\begin{informalproof}
(a) For $a\in\R$, the set $\{u\in\operatorname{dom}\psi:\psi(u)<a\}$ is the projection to the first coordinate of
$K_a=\{(u,v)\in[0,c_0]^2:u<v<a,\ P(u)=P(v)\}$: if $\psi(u)<a$, some $v$ in the set whose infimum is $\psi(u)$ is below
$a$, and conversely. $K_a$ is the intersection of a closed set with an open subset of the compact square, hence
$\sigma$-compact, so its projection is $\sigma$-compact and therefore Borel. Hence
$\operatorname{dom}\psi=\{\psi<c_0+1\}$ is Borel and $\psi$ is Borel on it, and so is
$D^+=\{u\in\operatorname{dom}\psi:\psi(u)>u\}$.
(b) $\psi(u)$ is the infimum of a nonempty subset of the closed set $\Lambda(P(u))\cap[u,c_0]$, so it lies in that set.
Let $u<u'$ be points of $D^+$ with $\psi(u)=\psi(u')=v$. Then $P(u)=P(v)=P(u')$, so $u'$ belongs to the set whose
infimum is $\psi(u)$, and $\psi(u)\le u'<\psi(u')=v$, a contradiction.
(c) The set $\psi^{-1}(W)\cap D^+$ is Borel by (a), and by (b) $\psi$ maps $\psi^{-1}(W)\cap D^+\cap\Lambda(p)$
injectively into $W\cap\Lambda(p)$ for every $p$. Integrate the inequality
$\#(\psi^{-1}(W)\cap D^+\cap\Lambda(p))\le\#(W\cap\Lambda(p))$ over $p$ and apply \cref{lem:mot-levels}(a) to both sides.
(d) If $\Lambda(p)$ is finite and $s\in\Lambda(p)$ is not its largest point, then $\psi(s)$ is the next point of
$\Lambda(p)$, so $\psi(s)>s$; the largest point is not in $\operatorname{dom}\psi$.
(e) With $Z$ as in the definition of a labelling, $\beta$ is constant on each of the finitely many components of
$[0,c_0]\setminus Z$, so it is Borel on $[0,c_0]$, and $A(\beta)$ is the union over $\ell$ of the sets
$\{u\in D^+:\beta(u)=\ell,\ \beta(\psi(u))\ne\ell\}$, which are Borel by (a). By (d),
$\kappa(p)=\#(A(\beta)\cap\Lambda(p))$ whenever $\Lambda(p)$ is finite, and $\kappa(p)=0$ otherwise. So
$\int\kappa\,\mathrm dp\le\nu(A(\beta))$, by monotonicity of the lower integral and \cref{lem:mot-levels}(a).
\end{informalproof}

\begin{lemma}[Finite affine labels, measure form]\label{lem:mot-finite-measure}
\cited{the finite step of the proof of~\cite[Theorem~6.2]{SecondPaper2026}, which applies it to push-forwards of
height measure; Lean-verified there, for every measure $\mu$ as in the statement, as
\leanname{TransposeFollowup.AffineMotion.exists_finite_affine_small_cut}. Not reproved here.}
Let $\Gamma_0$ be a finite set of affine bijections of $\R$, let $K\ge0$, and let $\mu$ be a finite Borel measure on
$[0,c_0]^2$ that is concentrated on the union of the graphs $\{(x,g(x))\}$ of the maps $g\in\Gamma_0$ and satisfies
$\mu(E\times[0,c_0])\le K\lambda(E)$ for every Borel set $E$. For all $\theta>0$ and $\varepsilon>0$ there is a
labelling $\beta$ whose classes have measure at most $\theta$ and $\mu(\{(x,y):\beta(x)\ne\beta(y)\})<\varepsilon$.
\end{lemma}

In the second paper's Lean library, \cref{lem:mot-finite-labels} is derived from this lemma, by pushing the measure of
levels forward to pairs of times in two affine pieces. The two lemmas are the only places below where amenability
enters.

\subsection{Theorems C, C\texorpdfstring{$'$}{'} and Cq}\label{sec:mot-statements}

We first define the class of Theorem~C. Its members are home-based in the sense of \cref{sec:prelim-walks}, up to a
tolerance, and every excursion from the home is made of boundedly many runs at constant rate.

\begin{definition}[Home-based heads with affine runs]\label{def:mot-hbar}
Let $K\ge1$ be an integer and $\rho\ge0$. A function $x\colon\{0,\dots,T\}\to\R$, for instance a walk of length $T+1$, is
\emph{HB-AR$(K,\rho)$} if there are a point $h$, its \emph{home}, and times $0=a_0<a_1<\dots<a_m=T$ such that every
\emph{segment} $[a_{l-1},a_l]$ satisfies:
\begin{enumerate}[label=(\roman*)]
\item (home-based) $\abs{x(a_{l-1})-h}\le\rho$ or $\abs{x(a_l)-h}\le\rho$;
\item (affine runs) there are times $a_{l-1}=s_0\le s_1\le\dots\le s_K=a_l$ and rates $q_1,\dots,q_K\in[-1,1]$ such
that $\abs{x(s)-x(s_{j-1})-q_j(s-s_{j-1})}\le\rho$ whenever $s_{j-1}\le s\le s_j$.
\end{enumerate}
The intervals $[s_{j-1},s_j]$ are the \emph{runs} of the segment. A machine \emph{has HB-AR$(K,\rho)$ work heads at
size $d$} if on every matrix input of size $d$ the trajectory of every work head, up to the halting time, is
HB-AR$(K,\rho)$.
\end{definition}

\begin{remark}[What the class contains]\label{rem:mot-hbar}
A head that stays within distance $\rho/2$ of its home during a time interval contributes a segment with a single run
of rate~$0$, so stays at home cost nothing. A home-based walk in the sense of \cref{sec:prelim-walks}, with home $h$,
whose legs consist of at most $K$ runs at constant rate up to deviation $\rho$ is HB-AR$(K,\rho+2)$: cut the time axis
at the first step of every leg; each resulting segment is contained in the corresponding leg, or exceeds it by one step
at its end. If $x$ is HB-AR$(K,\rho)$ and $\abs{x-x'}\le\rho'$, then $x'$ is HB-AR$(K,\rho+2\rho')$. With $\rho(d)=o(L(d))$ the
class therefore contains the usual home-based schedules: sweeps out and back from a home at the end of the used tape,
with any number of passes; copying excursions with small zigzags; excursions of all sizes in any order; any coupling
between the heads; excursions on both sides of the home; and homes that depend on the size and on the input. Nothing is
assumed about the input head or the output. The class excludes curved macroscopic motion: a macroscopic pause inside an
excursion costs a run, and an excursion whose speed varies along its length, for instance unit moves alternating with
pauses at a slowly varying duty cycle, needs unboundedly many runs.
\end{remark}

For a machine $M$ let $T_M(d)$ be the largest halting time of $M$ on matrix inputs of size $d$ (infinite if some input
does not halt).

\begin{theorem}[Theorem C: home-based heads with affine runs]\label{thm:C}
\notformalized
Let $K\ge1$ and $\rho(d)=o(L(d))$, and let $M$ be a machine that transposes. Suppose that for all large $d$ the machine
$M$ has HB-AR$(K,\rho(d))$ work heads at size $d$. Then $T_M(d)/L(d)\to\infty$. In particular $M$ does not witness~(T).
\end{theorem}

The hypothesis concerns every work head of $M$; for an oblivious machine it is a condition on its single head
schedule. For an adaptive machine it must hold on every input, with $K$ and $\rho$ uniform; the proof uses it only on
the representatives of the popular trace classes.

\paragraph{Why it matters.} Theorem~C is an unconditional impossibility result for a natural class: it needs neither
depth reduction nor the crux. The class allows any number of passes, excursions at all scales in any interleaving, any
coupling of the heads and arbitrary motion of the input head, and it includes adaptive machines. For oblivious machines whose work heads are all
home-based, \cref{cor:B} reaches the same conclusion for every motion, but only if the crux (\cref{def:crux}) holds;
Theorem~C needs no coding statement, for the subclass with affine runs. For a construction it
says that the work heads cannot all be home-based with legs made of boundedly many constant-rate runs. By Theorem~Cq,
in every limit, motion that is not countably-piecewise affine must carry at least half a block unit of travel in
total; for a home-based head such motion needs legs with unboundedly many runs whose speed varies along the leg, as
produced by counters or rulers.

\begin{theorem}[Theorem C$'$: open affine pieces]\label{thm:mot-Cprime}
\notformalized
Let $M$ be a machine that transposes. Suppose that for every $c>0$, every limit $(P_t)_{t<k}$ of $M$ at time constant
$c$ admits relatively open sets $U_t\subseteq[0,c_0]$ such that $P_t$ is affine on each connected component of $U_t$ and
$\nu_t([0,c_0]\setminus U_t)=0$. Then $T_M(d)/L(d)\to\infty$.
\end{theorem}

Theorem~C$'$ allows the non-affine set $[0,c_0]\setminus U_t$ to have positive Lebesgue measure, provided the head is
stationary on almost all of it. Theorem~C follows from the quantitative statement below, and so does C$'$.

\paragraph{Why it matters.} Theorem~C$'$ is the form of Theorem~Cq that is easiest to check on a concrete schedule: a
construction whose limit motion is affine on open intervals wherever the heads move fails, even when the set on which
the heads rest is a fat Cantor set (\cref{rem:mot-fatcantor}).

\begin{theorem}[Theorem Cq: curved motion]\label{thm:Cq}
\notformalized
Let $(P_t)_{t<k}$ be $1$-Lipschitz functions on $[0,c_0]$ that satisfy $(\star)_\emptyset$; by \cref{thm:Mstar} this
holds for every limit of every machine. For each $t$ let $E_{t,0},E_{t,1},\dots$ be Borel subsets of $[0,c_0]$ on each
of which $P_t$ agrees with an affine function, and let $N_t=[0,c_0]\setminus\bigcup_aE_{t,a}$. Then
\[
  \sum_{t<k}\nu_t(N_t)\;\ge\;\tfrac12 .
\]
\end{theorem}

In words: in every limit of a linear-time transposer, the limit paths of the work heads travel a total distance of at
least $\frac12$, in block units, while moving in a way that is not countably-piecewise affine. Since distance $1$
corresponds to about $L/G$ cells, this is about $L/(2G)$ cells of curved travel.

\paragraph{Why it matters.} Theorem~Cq is the quantitative statement behind Theorems~C and~C$'$, and it contains the
Lean-verified affine-time obstruction of the second paper (\cref{cor:mot-lambda}). It tells a construction how much
curved motion it needs: at least half a block unit of travel, whatever else the machine does. It does not exclude
curved motion (\cref{rem:mot-scope}).

\subsection{Proofs of Theorems Cq, C\texorpdfstring{$'$}{'} and C}\label{sec:mot-proofs}

We prove Theorem~Cq (\cref{thm:Cq}) first. The proof is the second paper's truncation argument with the speed measure in
place of Lebesgue measure: finitely many pieces are handled by \cref{lem:mot-finite-labels}, and the returns that involve
the other times are charged by their speed.

\begin{informalproof}
Fix $\theta\in(0,\frac12)$, $\varepsilon>0$ and $n\ge1$, and let $R_{t,n}=E_{t,0}\cup\dots\cup E_{t,n-1}$.
\Cref{lem:mot-finite-labels}, applied to these pieces, gives a labelling $\beta$ whose classes have measure at most
$\theta$ and $\sum_t\int\kappa_t^{R_{t,n}}\,\mathrm dp<\varepsilon$. By \cref{lem:mot-delete},
$\int\kappa_t\,\mathrm dp\le\int\kappa_t^{R_{t,n}}\,\mathrm dp+2\nu_t([0,c_0]\setminus R_{t,n})$ for every $t$, so
$(\star)_\emptyset$, applied to $\beta$, gives
\[
  1\;\le\;\varepsilon+2\theta+2\sum_{t<k}\nu_t\bigl([0,c_0]\setminus R_{t,n}\bigr).
\]
For each $t$ the sets $[0,c_0]\setminus R_{t,n}$ decrease to $N_t$ as $n\to\infty$, and $\nu_t$ is a finite measure, so
$\nu_t([0,c_0]\setminus R_{t,n})\to\nu_t(N_t)$. Letting $n\to\infty$ and then $\theta,\varepsilon\to0$ gives
$\sum_t\nu_t(N_t)\ge\frac12$.
\end{informalproof}

Theorem~C$'$ (\cref{thm:mot-Cprime}) follows from Theorem~Cq and the existence of limits.

\begin{informalproof}
Suppose not. Then there are $c>0$ and sizes $d_i\to\infty$ with $T_M(d_i)\le cL(d_i)$, so $M$ transposes and halts
within $cL(d_i)$ steps on every matrix input of size $d_i$, and by \cref{thm:Mstar}(b) it has a limit $(P_t)_{t<k}$ at
time constant $c$. Let $U_t$ be as in the hypothesis. The components of $U_t$ are countably many relatively open
intervals, on each of which $P_t$ is affine. By \cref{thm:Cq}, applied to these intervals, $\sum_t\nu_t([0,c_0]\setminus
U_t)\ge\frac12$, which contradicts $\nu_t([0,c_0]\setminus U_t)=0$.
\end{informalproof}

The proof of Theorem~C (\cref{thm:C}) rests on the structure of the limits of HB-AR heads: they sit at a home level,
apart from countably many \emph{teeth}, each piecewise affine with boundedly many pieces. The teeth may be dense in time.

\begin{lemma}[Limits of home-based heads with affine runs]\label{lem:mot-teeth}
\notformalized
Let $d_i\to\infty$, $c>0$, $K\ge1$ and $\rho_i=o(L(d_i))$, and let $x_i\colon\{0,\dots,T_i\}\to\Z$ be walks with
$x_i(0)=0$ and $1\le T_i\le cL(d_i)$ that are HB-AR$(K,\rho_i)$. Extend $x_i$ by $x_i(s)=x_i(T_i)$ for $s>T_i$, and suppose
that $P_i(u)=d_i^{-1}\lfloor x_i(\lfloor ud_i\rfloor r_{d_i})/r_{d_i}\rfloor$ converges uniformly on $[0,c_0]$ to a
function $P$. Then there is a level $p\in\R$ with the following properties.
\begin{enumerate}[label=(\alph*)]
\item The set $\{u\in[0,c_0]:P(u)\ne p\}$ is the union of countably many disjoint relatively open intervals of
$[0,c_0]$, the \emph{teeth}.
\item For each tooth $I$ there is a set $\Phi_I\subseteq I$ of at most $2K$ points such that $P$ is affine on each
component of $I\setminus\Phi_I$.
\item For every $y>0$, at most $c_0/(2y)+2$ teeth contain a point $u$ with $\abs{P(u)-p}>y$.
\item $\nu(\{u\in[0,c_0]:P(u)=p\})=0$, where $\nu=\abs{P'}\,\mathrm du$.
\end{enumerate}
\end{lemma}

\begin{informalproof}
Write $r_i=r_{d_i}$ and $R_i=R_{d_i}$. Let $X_i\colon[0,\infty)\to\R$ interpolate $x_i$ linearly between integers, and
let $Q_i(u)=X_i(uR_i)/R_i$; it is $1$-Lipschitz. If $s=\lfloor ud_i\rfloor r_i$, then $0\le uR_i-s<r_i$, so
$\abs{X_i(uR_i)-x_i(s)}\le r_i$, and $\abs{P_i(u)-x_i(s)/R_i}\le1/d_i$. Hence $\abs{P_i-Q_i}\le2/d_i$, so $Q_i\to P$
uniformly on $[0,c_0]$, and $P$ is $1$-Lipschitz (as in~\cite[Lemma~7.1(a)]{FirstPaper2026}).

\emph{Rescaled segments.} Let $h_i$ and $0=a_0<\dots<a_m=T_i$ be a home and segments for $x_i$ as in
\cref{def:mot-hbar}, and put $p_i=h_i/R_i$ and $\varepsilon_i=\rho_i/R_i$; then $\varepsilon_i\le G\rho_i/L(d_i)\to0$.
Divided by $R_i$, the segments become closed intervals that partition $[0,T_i/R_i]$, and we add the interval
$[T_i/R_i,c_0]$, on which $Q_i$ is constant, to the last segment as one more run, of rate~$0$; recall
$T_i/R_i\le cG<c_0$. On a run $[s_{j-1},s_j]$ the function $X_i(s)-X_i(s_{j-1})-q_j(s-s_{j-1})$ is linear between
consecutive integers, so its absolute value is at most $\rho_i$ at all real $s$ of the run. We obtain a partition of
$[0,c_0]$ into closed intervals, the segments, such that each segment has an endpoint at which
$\abs{Q_i-p_i}\le\varepsilon_i$ (for the extended last segment, recall that $Q_i$ is constant on the added interval),
and each segment is the union of at most $K$ consecutive closed intervals, its runs, on each of which
$\abs{Q_i(u)-Q_i(u')-q(u-u')}\le2\varepsilon_i$ for some $q\in[-1,1]$; the last segment has at most $K+1$ runs. The
first segment has an endpoint $s$ with $\abs{x_i(s)-h_i}\le\rho_i$, and $\abs{x_i(s)}\le s\le T_i$, so
$\abs{p_i}\le cG+\varepsilon_i$. Passing to a subsequence, which does not change $P$, we may assume that $p_i\to p$.

(a) The set $\{P\ne p\}$ is relatively open in $[0,c_0]$; its components are disjoint relatively open intervals, and
there are countably many of them.

(b) Let $I$ be a tooth and $[\alpha,\beta]\subseteq I$ a compact interval with $\alpha<\beta$. Then
$\delta=\min_{[\alpha,\beta]}\abs{P-p}>0$, and for large $i$ we have $\sup_{[0,c_0]}\abs{Q_i-P}<\delta/4$,
$\abs{p_i-p}<\delta/4$ and $\varepsilon_i<\delta/2$, so that $\abs{Q_i-p_i}>\delta/2>\varepsilon_i$ on $[\alpha,\beta]$.
Let $J_i$ be the component of $\{u\in[0,c_0]:\abs{Q_i(u)-p_i}>\varepsilon_i\}$ that contains $[\alpha,\beta]$. At most
two segments meet $J_i$: if a segment lies strictly between two segments that meet $J_i$, then both of its endpoints lie
between points of $J_i$, hence in the interval $J_i$, against the home-based condition. So $[\alpha,\beta]$ is covered
by the at most $2K+1$ runs of these two segments, and there are points
$\alpha=\tau_{i,0}\le\tau_{i,1}\le\dots\le\tau_{i,2K+1}=\beta$ and slopes $q_{i,j}\in[-1,1]$ such that
$\abs{Q_i(u)-Q_i(u')-q_{i,j}(u-u')}\le2\varepsilon_i$ for $u,u'\in[\tau_{i,j-1},\tau_{i,j}]$. Along a further
subsequence $\tau_{i,j}\to\tau_j$ and $q_{i,j}\to q_j$. If $\tau_{j-1}<u<u'<\tau_j$, then $u,u'\in[\tau_{i,j-1},\tau_{i,j}]$
for large $i$, and letting $i\to\infty$ gives $P(u')-P(u)=q_j(u'-u)$. So $P$ is affine near every point of
$(\alpha,\beta)\setminus\{\tau_1,\dots,\tau_{2K}\}$. Let $\Phi_I$ be the set of points $u$ in the interior of $I$ such
that $P$ is affine on no neighbourhood of $u$. Every compact interval in $I$ contains at most $2K$ points of $\Phi_I$ in
its interior, so $\Phi_I$ has at most $2K$ points. On the interior of each component of $I\setminus\Phi_I$ the function
$P$ is locally affine, hence affine, and by continuity it is affine on the whole component.

(c) A tooth that does not contain $0$ or $c_0$ has its endpoints at level $p$. If it contains a point $u$ with
$\abs{P(u)-p}>y$, then, $P$ being $1$-Lipschitz, $u$ is at distance more than $y$ from both endpoints, so the tooth has
length more than $2y$. The teeth are disjoint, and at most two of them contain $0$ or $c_0$.

(d) This is \cref{lem:mot-levels}(b).
\end{informalproof}

Theorem~C now follows from the lemma and Theorem~Cq.

\begin{informalproof}
Suppose not. Then there are $c>0$ and sizes $d_i\to\infty$ with $T_M(d_i)\le cL(d_i)$, and by \cref{thm:Mstar}(b) the
machine $M$ has a limit $(P_t)_{t<k}$ at time constant $c$ along a subsequence, with representatives $w_i$. For large
$i$ and each $t$, the trajectory $x_{t,w_i}$ up to the halting time $T_i\le cL(d_i)$ starts at cell $0$ and is
HB-AR$(K,\rho(d_i))$, and $T_i\ge L(d_i)\ge1$ because $M$ writes $L(d_i)$ output bits, at most one per step. So
\cref{lem:mot-teeth} applies to these trajectories. Hence there is a level $p_t$ such that the components of
the sets $I\setminus\Phi_I$, for the teeth $I$ of $P_t$, form countably many relatively open intervals on each of which
$P_t$ is affine, and they cover $[0,c_0]$ except for the set $\{P_t=p_t\}$ and countably many points. Both exceptional
sets are $\nu_t$-null, by \cref{lem:mot-teeth}(d) and \cref{lem:mot-levels}(c). So \cref{thm:Cq}, applied to these
intervals, gives $0\ge\frac12$, a contradiction.
\end{informalproof}

\subsection{Relation to the Lebesgue form}\label{sec:mot-lambda}

\begin{corollary}[Speed-weighted and Lebesgue forms]\label{cor:mot-lambda}
\notformalized
Let $(P_t)_{t<k}$ be $1$-Lipschitz functions on $[0,c_0]$ that satisfy $(\star)_\emptyset$, and let
$U_t\subseteq[0,c_0]$ be relatively open sets on which $P_t$ is locally affine. Then
$\sum_t\nu_t([0,c_0]\setminus U_t)\ge\frac12$, and hence $\sum_t\lambda([0,c_0]\setminus U_t)\ge\frac12$.
\end{corollary}

\begin{informalproof}
The components of $U_t$ are countably many relatively open intervals. On each of them $P_t$ is locally affine, hence
affine: its slope is locally constant on a connected set, hence constant, and then $P_t$ minus a linear function is
locally constant, hence constant. \Cref{thm:Cq} applied to these intervals gives the first inequality, and
$\nu_t\le\lambda$ gives the second.
\end{informalproof}

The second inequality is the quantitative affine-time obstruction of the second paper~\cite[Theorem~6.2]{SecondPaper2026},
which is Lean-verified there for paths satisfying its restricted star condition~\cite[Definition~6.1]{SecondPaper2026}.
That condition includes $(\star)_\emptyset$ as its case $S=\emptyset$ and is formally stronger; the corollary needs only
$(\star)_\emptyset$ and allows relatively open sets, so it slightly generalizes that theorem. Theorem~Cq carries the same
argument, F\o lner labels for finitely many affine pieces followed by truncation, from open sets and Lebesgue measure to
Borel pieces and the speed measure; the area formula charges the omitted returns by speed. The motion exclusions of~\cite[Sections~6 and~7]{SecondPaper2026}, for bounded words in explicit
operation palettes and for contracting recursions with regular stages, rest on that inequality.

The speed-weighted form is stronger exactly where the heads are slow on the non-affine set. For a limit of HB-AR heads,
\cref{lem:mot-teeth} shows that the largest relatively open set on which $P_t$ is locally affine misses only the
boundary of the home-level set $\{P_t=p_t\}$, a closed nowhere dense set, and countably many points. So the Lebesgue
form already excludes every machine with HB-AR work heads for which, in every limit, these boundaries have total
Lebesgue measure less than $\frac12$. Theorem~C is new for limits in which these nowhere dense sets have total measure at least $\frac12$, as
in the next example; they carry no speed, so for Theorem~Cq they cost nothing.

\begin{remark}[A fat Cantor schedule]\label{rem:mot-fatcantor}
Let $c$ satisfy $cG>1$, and let $C\subseteq[0,1]$ be a perfect nowhere dense set with $\lambda(C)=\frac12$ that
contains $0$ and $1$, for instance a Smith--Volterra--Cantor set. Let one work head do the following at size $d$: it stays at its home cell $0$
except during the time intervals $(aR_d,bR_d)$ for the gaps $(a,b)$ of $C$ of length at least $\delta_d$, where
$\delta_d\to0$ and $\delta_dR_d\to\infty$; during such an interval it moves away from home at unit speed until the
midpoint and back at unit speed. Every gap gives two segments of one run each, and the stays at home give segments with
one run of rate~$0$, so the trajectory is HB-AR$(1,1)$. The omitted gaps change the path by at most $\delta_d/2$, so the
rescaled trajectories converge uniformly to $P(u)=\operatorname{dist}(u,C)$ on $[0,1]$, constant afterwards. The
home-level set of $P$ within $[0,1]$ is $C$. No point of $C$ has a neighbourhood on which $P$ is affine: such an
affine function would vanish at that point and, $C$ being perfect, at points of $C$ arbitrarily close to it, hence
vanish identically, and $C$ would contain an interval. So every relatively open set on which $P$ is locally affine misses $C$, and the Lebesgue form gives nothing
because $\lambda(C)=\frac12$. But $\nu(C)=0$, so Theorem~C and Theorem~Cq exclude every machine all of whose work heads move
like this. Alternatively, $P$ is constant on the Borel set $C$, so $C$ can be taken as one of the pieces in Theorem~Cq;
the uncovered set is then countable, and even the Lebesgue-measure bound $\sum_t\lambda(N_t)\ge\frac12$ with Borel pieces,
which follows from Theorem~Cq because $\nu\le\lambda$, excludes this schedule.
\end{remark}

\begin{remark}[What escapes]\label{rem:mot-scope}
Theorem~Cq demands a minimum amount of speed-weighted motion that is not countably-piecewise affine; it says nothing
against curved motion. Theorem~Cr (\cref{thm:Cr}) excludes curved motion only when, in some clock, countably many
affine maps give the first returns at all times except a set carrying less than one block unit of travel.
A machine with one head whose home drifts, or whose motion is not home-based at all, falls outside Theorem~C even if
all its other heads are HB-AR; it is excluded only through Theorems~Cq and~Cr, that is, only if its limits are
countably-piecewise affine on enough of their speed, or have affine returns in a common clock on enough of it. Motion
that satisfies $(\star)$ exists, and by Theorems~Cq and~Cr it is necessarily curved, and in every clock its first
returns are not countably-piecewise affine on a set carrying at least one block unit of travel. The second paper constructs, for every machine, a fixed wrapper
that adds linear time and whose complete runs satisfy, whatever computation follows, the cut inequality
of~\cite[Definition~7.4]{FirstPaper2026} with zero slack, and also the conclusions of Theorem~M* when that computation
runs in linear time~\cite[Theorem~8.1]{SecondPaper2026}; this result is Lean-verified there. The heads of the wrapper
move along rational near-identity M\"obius maps (see~\cite[Section~13]{FirstPaper2026}). So Theorem~Cq is a constraint
on how a linear-time transposer moves its heads, not a lower bound by itself.
\end{remark}

\begin{remark}[How a machine can move along curves]\label{rem:mot-curved}
Proof sketch only: we indicate why curved macroscopic motion can be produced at amortised constant cost per step; this
is a design sketch, not a result. A bounded number of constant-rate runs per leg gives piecewise affine limits, so a
curved macroscopic leg needs unboundedly many runs whose speed, or duty cycle, varies along the leg. One way to produce
this is to hold two numbers $a,b$ of order $L^{\delta}$, for some
$0<\delta<1$, in unary on counter tapes and to alternate $b$ moves of one head with $a$ moves of another, each timed by
a sweep of a counter head; $a$ and $b$ change by single-cell edits scheduled by a slower accumulator. The two heads then
move at average speeds $b/(a+b)$ and $a/(a+b)$, with position error $o(L)$. For a M\"obius profile $f$ the function $1/f'$ is a quadratic
polynomial, so a two-level hierarchy of unary accumulators, a discrete second difference, suffices. This realizes
curved macroscopic motion only; the cell-exact pairings that data movement needs are a separate matter.
\end{remark}

\subsection{Consequences}\label{sec:mot-consequences}

The consequences concern the number of tapes that must interact, affine returns and changes of clock, and hierarchical
schedules; the first two, with Theorem~Cq, give the design constraint of \cref{cor:mot-design}.

\begin{corollary}[One work tape never expands]\label{cor:mot-one-tape}
\notformalized
\begin{enumerate}[label=(\alph*)]
\item Let $P$ be $1$-Lipschitz on $[0,c_0]$. For all $\theta>0$ and $\varepsilon>0$ there is a labelling $\beta$ whose
classes have measure at most $\theta$ and $\int\kappa\,\mathrm dp\le\varepsilon$.
\item Let $(P_t)_{t<k}$ be $1$-Lipschitz functions on $[0,c_0]$ that satisfy $(\star)_\emptyset$. Then at least two of
the $P_t$ are not monotone, and the non-monotone ones do not all have the same level sets: they are not all of the form
$h_t\circ P$ for one of them, $P$, and injective maps $h_t$.
\item Consequently, if $M$ has one work tape and transposes, then $T_M(d)/L(d)\to\infty$. In every limit of a machine
that transposes within $cL(d)$ steps on all inputs, at least two work-head limits are non-monotone, and they do not all
have the same level sets.
\end{enumerate}
\end{corollary}

\begin{informalproof}
(a) Let $\delta>0$ be a band width, to be chosen at the end. By \cref{lem:mot-levels}(b) the level set $\Lambda(p)$ is
finite for almost every $p$, so for almost every $\eta\in[0,1)$ all the levels $p_j=(j+\eta)\delta$, $j\in\Z$, have finite level sets; fix
such an $\eta$. The bands $T_j=\{u\in[0,c_0]:p_j\le P(u)<p_{j+1}\}$ partition $[0,c_0]$, only finitely many of them are
nonempty, and each is a finite union of intervals, because $P$ is continuous and takes each value $p_j$ only finitely
often. In a band with $\lambda(T_j)\le\theta$ make no cut, and put $m_j=0$. In each other band choose
$m_j=\lceil\lambda(T_j)/\theta\rceil-1\le\lambda(T_j)/\theta$ cut times, the first times at which the continuous function
$x\mapsto\lambda(T_j\cap[0,x])$ takes the values $\theta,2\theta,\dots,m_j\theta$. Let $\beta$ label $u\in T_j$ by a
natural number that encodes, injectively, the pair formed by $j$ and the number of cuts of $T_j$ that are at most $u$,
and label every $u\notin[0,c_0]$ by $0$. Then $\beta$ is constant on each of finitely many intervals that partition
$[0,c_0]$, so it is piecewise constant off $\emptyset$ (take $Z$ to be their endpoints). Its classes have measure at most
$\theta$: within $T_j$ they are the parts of $T_j$ before the first cut, between consecutive cuts, and after the last
cut, of measures $\theta,\dots,\theta$ and $\lambda(T_j)-m_j\theta\le\theta$.
Let $p\in[p_j,p_{j+1})$ have a finite level set, and let $s<s'$ be consecutive points of $\Lambda(p)$. Both lie in $T_j$,
so $\beta(s)\ne\beta(s')$ only if a cut of $T_j$ lies in $(s,s']$. These intervals are disjoint for different
consecutive pairs of $\Lambda(p)$, so $\kappa(p)\le m_j$. Hence $\int\kappa\,\mathrm dp\le\delta\sum_jm_j\le\delta
c_0/\theta$, which is $\varepsilon$ for $\delta=\varepsilon\theta/c_0$.

(b) If $P_t$ is monotone, each of its level sets is empty, a point or an interval, so $\kappa_t(p)=0$ for every labelling
and every $p$. Now suppose that the non-monotone $P_t$, if there are any, are all of the form $h_t\circ P$ for one of
them, $P$, and injective maps $h_t$. The interval $I=P([0,c_0])$ is nondegenerate, and $h_t$ is continuous on it: for
closed $F\subseteq\R$ the set $\{q\in I:h_t(q)\in F\}=P(P_t^{-1}(F))$ is compact. Being injective, $h_t$ is strictly
monotone on $I$. So $\Lambda_t(h_t(q))=\Lambda(q)$ for $q\in I$, the other level sets of $P_t$ are empty, and
$\kappa_t(h_t(q))=\kappa(q)$ for every labelling. Take $\theta=\frac14$. By uniform continuity there is $\delta>0$ such
that each of these $h_t$ maps every subinterval of $I$ of length at most $\delta$ onto an interval of length at most
$\theta/(4kc_0)$. Let $\beta$ be the labelling of (a) for $P$ with band width $\delta$, or any labelling with classes of
measure at most $\frac14$ if all $P_t$ are monotone. For each band, $\kappa_t\le m_j$ on the interval
$h_t([p_j,p_{j+1})\cap I)$, and $\kappa_t=0$ off the union of these intervals, so
$\int\kappa_t\,\mathrm dp\le\sum_jm_j\,\theta/(4kc_0)\le1/(4k)$. Hence $\sum_t\int\kappa_t\,\mathrm dp\le\frac14$, and
$(\star)_\emptyset$ gives $1\le\frac14+\frac12$, a contradiction.

(c) If $M$ has one work tape and transposes, and $T_M(d_i)\le cL(d_i)$ along some $d_i\to\infty$, then by
\cref{thm:Mstar}(b) $M$ has a limit, which satisfies $(\star)_\emptyset$ and consists of a single path, against (b). The
second statement is (b) for limits, which satisfy $(\star)_\emptyset$ by \cref{thm:Mstar}(a).
\end{informalproof}

\paragraph{Why it matters.} For one work tape the statement is much weaker than the classical superlinear lower bounds
in neighbouring models: $\Omega(L^{3/2}/(\log L)^{1/2})$ steps with one work tape and no output tape~\cite{DMS1991}, and
$\Theta(L^{5/4})$ for Boolean matrices with one work tape and a write-only output tape~\cite{DietzfelbingerMaass1993}.
Its interest is the reason: label time by thin height bands of the head's own position, and cut each band into pieces of
small measure. Consecutive visits of a level lie in one band and get different labels only across a cut, and each cut
separates at most one consecutive pair at each level; so crossings are rare, however curved the motion. The expansion that $(\star)$ demands must come from the interplay of
the returns of two or more heads moving along non-monotone paths with different level sets. Stationary and monotone heads
contribute nothing, and heads that all move in lockstep, even after a fixed change of scale on each tape, cannot supply
the expansion by themselves.

The next theorem replaces the affine pieces of Theorem~Cq by affine returns, seen through a common change of clock.

\begin{theorem}[Theorem Cr: affine returns]\label{thm:Cr}
\notformalized
Let $(P_t)_{t<k}$ be $1$-Lipschitz functions on $[0,c_0]$ that satisfy $(\star)_\emptyset$. Let $\varphi$ be a
homeomorphism of $[0,c_0]$ onto an interval $J$ such that $\varphi^{-1}$ maps $\lambda$-null sets to $\lambda$-null sets,
for instance the identity; we call such a $\varphi$ a \emph{clock}. Let $\Gamma$ be a countable set of affine bijections
of $\R$. For each $t$ let
\[
  X_t=\{u\in D^+_t:\varphi(\psi_t(u))\ne g(\varphi(u))\text{ for every }g\in\Gamma\}
\]
be the set of times whose first return is not given, in the clock $\varphi$, by a map in $\Gamma$; it is Borel by
\cref{lem:mot-returns}(a). Then
\[
  \sum_{t<k}\nu_t(X_t)\;\ge\;1 .
\]
\end{theorem}

In words: in every limit of a linear-time transposer, in every clock and for every countable set of affine maps, the
times at which a head's next return to its current level is not given by one of these maps carry at least one block
unit of travel in total.

\begin{informalproof}
Translating $J$, and conjugating the maps in $\Gamma$ by the same translation, we may assume that $J=[0,c_1]$. The
finite measure $A\mapsto\lambda(\varphi^{-1}(A))$ on $J$ vanishes on $\lambda$-null sets, so by the Radon--Nikodym
theorem it has a density $\rho$ with respect to $\lambda$, which we take Borel and finite everywhere, and for every
$\theta>0$ there is $\delta>0$ with $\lambda(\varphi^{-1}(A))\le\theta$ whenever $\lambda(A)\le\delta$. Let
$\Gamma_1\subseteq\Gamma_2\subseteq\cdots$ be finite sets with union $\Gamma$. For $n\ge1$ let $B_{t,n}$ be the set of
$u\in D^+_t$ with $\rho(\varphi(u))\le n$ and $\varphi(\psi_t(u))=g(\varphi(u))$ for some $g\in\Gamma_n$; it is Borel by
\cref{lem:mot-returns}(a).

Fix $\theta\in(0,\frac12)$, $\varepsilon>0$ and $n\ge1$, and choose $\delta$ for $\theta$. Push the restriction of
$\nu_t$ to $B_{t,n}$ forward along $u\mapsto(\varphi(u),\varphi(\psi_t(u)))$ and sum over $t$. This gives a finite Borel
measure $\mu$ on $J^2$, concentrated on the graphs of the maps in $\Gamma_n$, and by \cref{lem:mot-levels}(c), for
every Borel set $E\subseteq J$,
\[
  \mu(E\times J)=\sum_t\nu_t\bigl(B_{t,n}\cap\varphi^{-1}(E)\bigr)
  \le k\,\lambda\bigl(\varphi^{-1}(E\cap\{\rho\le n\})\bigr)=k\int_{E\cap\{\rho\le n\}}\rho\,\mathrm d\lambda
  \le kn\,\lambda(E).
\]
\Cref{lem:mot-finite-measure}, on $J$ and with $\delta$ in place of $\theta$, gives a labelling $\beta'$ whose classes
have measure at most $\delta$ and $\mu(\{(x,y):\beta'(x)\ne\beta'(y)\})<\varepsilon$. Let $\beta=\beta'\circ\varphi$ on
$[0,c_0]$ and $\beta=0$ elsewhere. Since $\varphi$ is a homeomorphism, $\beta$ is piecewise constant off $\emptyset$,
with the finite set $\varphi^{-1}(Z')$, where $Z'$ is that of $\beta'$. Its classes are the preimages under $\varphi$ of
those of $\beta'$, so by the choice of $\delta$ they have measure at most $\theta$. A point $u\in B_{t,n}$ lies in $A_t(\beta)$ exactly when
$\beta'(\varphi(u))\ne\beta'(\varphi(\psi_t(u)))$, so
$\sum_t\nu_t(A_t(\beta)\cap B_{t,n})=\mu(\{(x,y):\beta'(x)\ne\beta'(y)\})<\varepsilon$. By $(\star)_\emptyset$ and
\cref{lem:mot-returns}(e),
\[
  1\;\le\;\sum_t\int\kappa_t\,\mathrm dp+2\theta\;\le\;\sum_t\nu_t(A_t(\beta))+2\theta
  \;<\;\sum_t\nu_t(D^+_t\setminus B_{t,n})+\varepsilon+2\theta .
\]
As $n\to\infty$, the sets $B_{t,n}$ increase to $D^+_t\setminus X_t$, so $\nu_t(D^+_t\setminus B_{t,n})\to\nu_t(X_t)$,
and $1\le\varepsilon+2\theta+\sum_t\nu_t(X_t)$. Letting $\theta,\varepsilon\to0$ proves the theorem.
\end{informalproof}

For $\Gamma=\emptyset$ the theorem says $\sum_t\nu_t(D^+_t)\ge1$. By \cref{lem:mot-levels}(a) and~(b) and
\cref{lem:mot-returns}(d) the left side is the integrated level excess $\sum_t\int(\#\Lambda_t(p)-1)^+\,\mathrm dp$. The
second paper also bounds it below by $1$, for limits of linear-time runs that satisfy the first paper's cut inequality
with zero slack~\cite[Section~8]{SecondPaper2026}, and the bound follows directly from $(\star)_\emptyset$, since
$\kappa_t\le(\#\Lambda_t-1)^+$. The constant is $1$, not the $\frac12$ of Theorem~Cq, because each crossing is charged
once, to the earlier of its two visits, whereas deleting a visit can break two consecutive pairs.

A countable union of countable sets is countable, so the sets $X_t$ have a smallest version. If countable sets
$\Gamma_j$ make $\sum_t\nu_t(X_t)$ tend to its infimum over all countable $\Gamma$, their union $\Gamma^*$ attains it, and
then, for every countable $\Gamma$, $\nu_t$-almost every time in the set $X_t$ for $\Gamma^*$ also lies in the set $X_t$
for $\Gamma$. So in every clock the sets $X_t$ for $\Gamma^*$ carry at least one block unit of travel in total, and on no
part of them that carries travel are the first returns, seen in that clock, countably-piecewise affine.

The hypotheses on the clock matter. A Cantor function of a null Cantor set is continuous and nondecreasing but constant
on each gap, and almost every time lies in a gap; by \cref{lem:mot-returns}(c), $\nu_t$-almost no first return lands in
the Cantor set. With it in place of $\varphi$, countably many translations give the first return of $\nu_t$-almost every
time, so $\sum_t\nu_t(X_t)=0$ for every family of paths, and the conclusion fails.

Affine pieces, possibly after a change of scale on the tape, produce affine returns.

\begin{lemma}[Returns between affine pieces]\label{lem:mot-transitions}
\notformalized
Let $P$ be $1$-Lipschitz on $[0,c_0]$, let $\varphi\colon[0,c_0]\to\R$ be a Borel map, and let $h\colon\R\to\R$ be
injective. Let $E_a$ ($a\in\N$) be Borel subsets of $[0,c_0]$ and $f_a$ affine maps with $P=h\circ f_a\circ\varphi$ on
$E_a$, and let $N=[0,c_0]\setminus\bigcup_aE_a$. Let $\Gamma$ be the countable set of the maps $f_b^{-1}\circ f_a$ for
which neither $f_a$ nor $f_b$ is constant. Then the set of $u\in D^+$ with $\varphi(\psi(u))\ne g(\varphi(u))$ for every
$g\in\Gamma$, which is Borel by \cref{lem:mot-returns}(a), has $\nu$-measure at most $2\nu(N)$.
\end{lemma}

\begin{informalproof}
If $f_a$ is constant, then $P$ is constant on $E_a$, so $E_a$ lies in a level set and $\nu(E_a)=0$ by
\cref{lem:mot-levels}(b). Let $E^*$ be the union of the $E_a$ with $f_a$ not constant. The Borel set
$N^*=[0,c_0]\setminus E^*$ contains $N$ and is contained in $N$ together with the $\nu$-null sets $E_a$ with $f_a$
constant, so $\nu(N^*)=\nu(N)$. Let $u\in D^+$ with $u\in E_a$ and $\psi(u)\in E_b$, where neither $f_a$ nor $f_b$ is
constant. By \cref{lem:mot-returns}(b), $h(f_b(\varphi(\psi(u))))=P(\psi(u))=P(u)=h(f_a(\varphi(u)))$, and since $h$ is
injective, $\varphi(\psi(u))=f_b^{-1}(f_a(\varphi(u)))$. So the set in question is contained in
$N^*\cup(\psi^{-1}(N^*)\cap D^+)$, which has $\nu$-measure at most $2\nu(N)$ by \cref{lem:mot-returns}(c).
\end{informalproof}

\begin{corollary}[A common clock change and changes of scale do not help]\label{cor:mot-clock}
\notformalized
Let $(P_t)_{t<k}$ be $1$-Lipschitz functions on $[0,c_0]$ that satisfy $(\star)_\emptyset$. Let $\varphi$ be a
homeomorphism of $[0,c_0]$ onto an interval $J$ such that $\varphi^{-1}$ maps $\lambda$-null sets to $\lambda$-null sets,
for instance an increasing homeomorphism with absolutely continuous inverse, and for each $t$ let $h_t$ be an injective
map of $\R$ into itself, for instance a homeomorphism of $\R$. For each $t$ let $E_{t,a}$ ($a\in\N$) be Borel subsets
of $J$ on each of which $P_t\circ\varphi^{-1}=h_t\circ f$ for an affine function $f$, and let
$N_t=J\setminus\bigcup_aE_{t,a}$. Then $\sum_t\nu_t(\varphi^{-1}(N_t))\ge\frac12$.
\end{corollary}

\begin{informalproof}
For each $t$ apply \cref{lem:mot-transitions} to $P_t$, $\varphi$, $h_t$ and the Borel pieces $\varphi^{-1}(E_{t,a})$,
whose union misses exactly $\varphi^{-1}(N_t)$, and let $\Gamma$ be the union over $t$ of the countable sets of affine
maps that it produces. The sets $X_t$ of \cref{thm:Cr} for $\varphi$ and $\Gamma$ lie in the sets that the lemma bounds
for head $t$, because $\Gamma$ contains the maps produced for $t$. So $\nu_t(X_t)\le2\nu_t(\varphi^{-1}(N_t))$, and
\cref{thm:Cr} gives $1\le2\sum_t\nu_t(\varphi^{-1}(N_t))$.
\end{informalproof}

With $\varphi$ and all $h_t$ the identity this is Theorem~Cq, which therefore also follows from Theorem~Cr. The corollary
covers, for example, a global time profile that slows all heads down and speeds them up together, a shared duty cycle,
and heads whose speed depends in a fixed way on their position on the tape.

\paragraph{Why it matters.} Curvature that is common to all heads, or that comes from a fixed change of scale
on each tape, does not create expansion: if one change of clock and such changes of scale make every head
countably-piecewise affine, up to sets carrying little speed, then $(\star)$ fails. Non-affine motion is therefore
necessary but not sufficient. What is required, by Theorem~Cr, is that in no common clock are the first returns of the
heads given by countably many affine maps, except at times that carry less than one block unit of travel. Behind this
is amenability: the affine maps of $\R$ form a solvable group, and \cref{lem:mot-finite-measure} builds its labels from
F\o lner sets of finitely many of them. Head-dependent clocks, $P_t=f_t\circ\varphi_t$ with countably-piecewise affine
$f_t$ and different $\varphi_t$, are not covered; a M\"obius clock on one head and affine motion on another produce M\"obius transitions,
which is the regime of the monotone expanders of Bourgain and Yehudayoff~\cite{BourgainYehudayoff2013}.

Next we consider hierarchical schedules in which every move is a sweep at constant speed and every excursion returns
exactly to where it started. They may be arbitrarily deep and interleaved.

\begin{definition}[Stroke trees]\label{def:mot-stroke}
Let $K,K'\ge1$ be integers, $\gamma\in(0,1)$, $\eta_0>0$, $\rho\ge0$ and $R>0$, and let $x\colon\{0,\dots,T\}\to\Z$ be a
walk of length $T+1$. For integers $0\le a\le b\le T$ the \emph{time interval} $[a,b)$ is $\{a,\dots,b-1\}$; a finite family of time
intervals \emph{tiles} $[a,b)$ if its members are nonempty and pairwise disjoint and their union is $[a,b)$. A time
interval $[a,b)$ is a \emph{pause} if $x$ is constant on $\{a,\dots,b\}$. A \emph{stroke tree} for $x$ with parameters
$(K,K',\gamma,\eta_0)$, tolerance $\rho$ and scale $R$ is a finite rooted forest whose nodes are \emph{strokes} and
\emph{V's}, each with a time interval, such that:
\begin{enumerate}[label=(\roman*)]
\item the intervals of the roots tile $[0,T)$, and every root is a stroke;
\item every stroke $S$ has a direction $\sigma_S\in\{-1,1\}$, and its interval is tiled by at most $K$ \emph{runs}, any
number of pauses, and the intervals of its children, which are V's; a run is a time interval $[a,b)$ such that
$\sigma_Sx$ is nondecreasing on $\{a,\dots,b\}$ and, for some $q\in[0,1]$, $\abs{x(s)-x(a)-\sigma_Sq(s-a)}\le\rho$ for
$a\le s\le b$;
\item every V, with interval $[a,b)$, satisfies $x(a)=x(b)$, and its interval is tiled by any number of pauses and the
intervals of its children, which are at most $K'$ strokes;
\item every root either has extent at least $\eta_0R$ or has no runs, in which case it is a \emph{home stroke}; and every
stroke $C$ that is a child of a V of a stroke $S$ other than a home stroke satisfies
$\operatorname{ext}(C)\le\gamma\operatorname{ext}(S)$.
\end{enumerate}
Here the \emph{extent} $\operatorname{ext}(S)$ of a stroke $S$ is the sum of $\abs{x(b)-x(a)}$ over its runs $[a,b)$.
\end{definition}

Every stroke moves monotonically through at most $K$ runs, with any number of excursions (V's) inserted between runs and
at its ends; every excursion returns exactly to its starting cell and consists of at most $K'$ smaller strokes. A home
stroke models a fixed home from which excursions of any size start.

\begin{proposition}[Stroke trees are affine or stationary]\label{prop:mot-strokes}
\notformalized
Let $K,K'\ge1$, $\gamma\in(0,1)$, $\eta_0>0$ and $\rho(d)=o(L(d))$. Let $M$ be a machine and $t$ one of its work heads,
and suppose that for all large $d$, on every matrix input of size $d$, the trajectory of head $t$ up to the halting time
has a stroke tree with parameters $(K,K',\gamma,\eta_0)$, tolerance $\rho(d)$ and scale $R_d$. Let $(P_s)_{s<k}$ be a
limit of $M$ at some time constant, and let $U$ be the union of the open intervals in $(0,c_0)$ on which $P_t$ is affine.
Then $P_t$ is affine on each component of $U$, and $\nu_t([0,c_0]\setminus U)=0$. Consequently, if every work head of
$M$ satisfies the hypothesis and $M$ transposes, then $T_M(d)/L(d)\to\infty$.
\end{proposition}

\begin{informalproof}
Write $P=P_t$, let $d_i$ and $w_i$ be as in \cref{def:mot-limit}, let $x_i=x_{t,w_i}$, frozen after its halting time
$T_i$, let $R_i=R_{d_i}$, and let $Q_i(u)=X_i(uR_i)/R_i$, where $X_i$ interpolates $x_i$ linearly between integers. As
in the proof of \cref{lem:mot-teeth}, $Q_i\to P$ uniformly on $[0,c_0]$. The set $U$ is open, $P$ is affine on each of
its components because it is locally affine there, and $Z=[0,c_0]\setminus U$ is closed.

\emph{Claim: $\abs{P(v)-P(w)}\le\lambda([w,v]\cap U)$ for all $0\le w<v\le c_0$.} Fix $w<v$ and
$\eta\in(0,\min(\eta_0,c_0))$, and fix a large $i$; we drop the index $i$ where this causes no ambiguity. Call a stroke
\emph{big} if its extent is at least $\eta R$ and \emph{small} otherwise; home strokes have no runs and play no role. Let
$w''=\min(\lceil wR\rceil,T)$ and $v''=\min(\lfloor vR\rfloor,T)$. Since $X$ is $1$-Lipschitz and constant after $T$,
$\abs{Q(v)-Q(w)}\le(\abs{x(v'')-x(w'')}+2)/R$. For large $i$ either $w''<v''$, or $w''=v''=T$ and $x(v'')=x(w'')$; so
assume $w''<v''$, and let $I=[w'',v'')$.

Every time $s\in I$ has a \emph{chain}: the root containing $s$; then the run, pause or V of that root containing $s$;
then, if this is a V, the pause or stroke of that V containing $s$; and so on, down to a run or a pause. Call a V
\emph{inside} if its interval is contained in $I$. The \emph{tile} of $s$ is the interval of the first inside V on its
chain, or, if there is none, the run or pause at the end of its chain. Two tiles are disjoint or equal: nested intervals
of the forest belong to a node and a descendant, and a tile that is an inside V contains the tiles of none of its
descendants, since their chains pass through it. So the sets $H\cap I$, for the tiles $H$, tile $I$ into intervals, and
$x(v'')-x(w'')$ is the sum of the displacements of $x$ over them. Inside V's and pauses have displacement $0$. A tile
that is a run belongs to a stroke $S$ none of whose ancestor V's is inside; call such strokes \emph{exposed}. Since $x$ is
monotone with unit steps on a run, its displacement over a subinterval $[a',b')$ of a run $[a,b)$ is at most
$\min(\abs{x(b)-x(a)},\,b'-a')$. Hence
\[
  \abs{x(v'')-x(w'')}\le\sum_{S\text{ exposed}}c(S),\qquad
  c(S)\le\min\bigl(\operatorname{ext}(S),\ \#(I\cap\mathrm{runs}(S))\bigr),
\]
where $\mathrm{runs}(S)$ is the union of the runs of $S$.

\emph{Small exposed strokes.} An exposed stroke that is not a root has a parent V that meets $I$ but is not inside;
this V contains $w''-1$ and $w''$, or it contains $v''-1$ and $v''$. The V's of the first kind contain $w''$, so they lie
on the chain of $w''$, and they are its first $J$ V's for some $J$, say $V_1\supset V_2\supset\dots\supset V_J$, with
$V_j$ a child of the stroke $S_{j-1}$ of that chain and $S_j$ a child of $V_j$. For $j\ge2$ the stroke $S_{j-1}$ is not a
root, hence not a home stroke, and repeated use of (iv) gives $\operatorname{ext}(S_{j-1})\le\gamma^{j-2}
\operatorname{ext}(S_1)\le\gamma^{j-2}T$. So every child of $V_j$ has extent at most $\gamma^{j-1}T\le\gamma^{j-1}c_0R$,
because $T\le cL\le c_0R$. Each $V_j$ has at most $K'$ children, so the small exposed children of $V_1,\dots,V_J$
contribute at most
\[
  K'\eta R+K'R\sum_{j\ge1}\min(\eta,\gamma^jc_0)\;\le\;K'\eta R\Bigl(1+\log_{1/\gamma}\frac{c_0}{\eta}+\frac1{1-\gamma}\Bigr):
\]
at most $\log_{1/\gamma}(c_0/\eta)$ terms of the sum equal $\eta$, and the others form a geometric series with first
term below $\eta$. The V's of the second kind lie on the chain of $v''-1$, and their small children contribute at most
the same amount. Every small exposed stroke is a child of one of these V's, because the roots are big or home strokes.
So the small exposed strokes contribute at most $\delta(\eta)R$, where
$\delta(\eta)=2K'\eta(1+\log_{1/\gamma}(c_0/\eta)+1/(1-\gamma))$ tends to $0$ with $\eta$.

\emph{Big strokes.} Runs of different strokes are disjoint, and a big stroke spends at least $\eta R$ steps in its runs.
Since $T\le c_0R$, there are at most $c_0/\eta$ big strokes and at most $m=\lfloor Kc_0/\eta\rfloor$ runs of big
strokes. List these runs as $[a_{i,l},b_{i,l})$, $l\le m$, padding with empty intervals, and let
$J_{i,l}=[a_{i,l}/R_i,\,b_{i,l}/R_i]$. Since $[w''/R,v''/R]\subseteq[w,v]$, the big exposed strokes contribute at most
$\sum_l\#(I\cap[a_{i,l},b_{i,l}))\le R_i\sum_l\lambda([w,v]\cap J_{i,l})$. Altogether
\[
  \abs{Q_i(v)-Q_i(w)}\;\le\;\delta(\eta)+\sum_{l\le m}\lambda\bigl([w,v]\cap J_{i,l}\bigr)+\frac2{R_i}.
\]

\emph{Limit.} By (ii), as in the proof of \cref{lem:mot-teeth}, $\abs{Q_i(u)-Q_i(u')-q_{i,l}(u-u')}\le2\rho(d_i)/R_i$
for $u,u'\in J_{i,l}$, with slopes $q_{i,l}\in[-1,1]$. Along a subsequence the endpoints of the $J_{i,l}$ and the
slopes $q_{i,l}$ converge, to intervals $J_l\subseteq[0,c_0]$ and slopes $q_l$. The $J_{i,l}$ have pairwise disjoint
interiors, hence so do the $J_l$. On the interior of each $J_l$ the function $P$ is affine with slope $q_l$, as in the
proof of \cref{lem:mot-teeth}, so this interior lies in $U$. Moreover $\lambda([w,v]\cap J_{i,l})\to\lambda([w,v]\cap
J_l)$, and $\sum_l\lambda([w,v]\cap J_l)\le\lambda([w,v]\cap U)$ because the interiors of the $J_l$ are disjoint and
lie in $U$. Letting $i\to\infty$ gives $\abs{P(v)-P(w)}\le\delta(\eta)+\lambda([w,v]\cap U)$, and letting $\eta\to0$
proves the claim.

\emph{Conclusion.} At almost every point $u$ of $Z$ the function $P$ is differentiable and $u$ is a Lebesgue density point
of $Z$. For such $u$ and small $\tau>0$ the claim gives
$2\tau\abs{P'(u)}+o(\tau)=\abs{P(u+\tau)-P(u-\tau)}\le\lambda([u-\tau,u+\tau]\setminus Z)=o(\tau)$, so $P'(u)=0$. Hence $\nu_t(Z)=0$. The
last sentence of the proposition follows from \cref{thm:mot-Cprime}, with $U_s$ the set $U$ of head $s$.
\end{informalproof}

\paragraph{Why it matters.} Stroke trees model hierarchical copy-and-return schedules of unbounded depth, with roots of
macroscopic size or a fixed home. The proposition says that, however deep and however interleaved, such motion is
affine or stationary in the limit, and hence excluded when all work heads move this way. The non-affine sets
$[0,c_0]\setminus U$ are nowhere dense, since $P_t'=0$ almost everywhere on them. The Lebesgue form already excludes stroke trees whose non-affine sets have total
Lebesgue measure less than $\frac12$; the proposition is needed when that measure is at least $\frac12$, as for
Cantor-type arrangements of V's. A construction must break one of
the hypotheses, for instance with strokes made of unboundedly many runs of varying density, which are curved; V's that
drift; V's with unboundedly many strokes, or nested strokes whose extents do not contract geometrically; or roots that
are not macroscopic, concatenated with varying density.

\begin{corollary}[The motion a linear-time transposer needs]\label{cor:mot-design}
\notformalized
Let $M$ transpose within $cL(d)$ steps on every matrix input of every size $d$, for instance a machine that
witnesses~(T) with constant $c/2$, and let $(P_t)_{t<k}$ be a limit of $M$ at time constant $c$; such limits exist along
a subsequence of every sequence of sizes. Then:
\begin{enumerate}[label=(\alph*)]
\item (curved) for all Borel sets $E_{t,a}$ on each of which $P_t$ is affine, the uncovered sets $N_t$ satisfy
$\sum_t\nu_t(N_t)\ge\frac12$;
\item (head-dependent) this survives every common change of clock and every fixed change of scale on each tape: for
every homeomorphism $\varphi$ of $[0,c_0]$ onto an interval $J$ such that $\varphi^{-1}$ maps null sets to null sets and
every countable set $\Gamma$ of affine bijections of $\R$, the times $u\in D^+_t$ with
$\varphi(\psi_t(u))\notin\{g(\varphi(u)):g\in\Gamma\}$ carry travel at least $1$ in total, that is, the sum over $t$
of their $\nu_t$-measures is at least $1$; in particular, for all
injective maps $h_t$ of $\R$ and all Borel sets $E_{t,a}\subseteq J$ on each of which $P_t\circ\varphi^{-1}=h_t\circ f$
for an affine $f$, the uncovered sets $N_t\subseteq J$ satisfy $\sum_t\nu_t(\varphi^{-1}(N_t))\ge\frac12$;
\item (interacting tapes) at least two of the $P_t$ are not monotone, and these do not all have the same level sets;
in particular
$k\ge2$.
\end{enumerate}
\end{corollary}

\begin{informalproof}
Limits exist by \cref{thm:Mstar}(b), and they satisfy $(\star)_\emptyset$ by \cref{thm:Mstar}(a). Then (a) is
\cref{thm:Cq}, (b) is \cref{thm:Cr,cor:mot-clock} and (c) is \cref{cor:mot-one-tape}(b).
\end{informalproof}

\paragraph{Why it matters.} This is the motion a construction must use: curved, with curvature that differs between
heads, on at least two work tapes whose returns interact. The monotone expanders of Bourgain and
Yehudayoff~\cite{BourgainYehudayoff2013}, driven along near-identity M\"obius maps as in the architecture
of~\cite[Section~13]{FirstPaper2026}, are consistent with all three conditions. The corollary constrains only the
macroscopic limit of the motion. Structure at mesoscopic scales, such as nested excursions much shorter than $L$, is
invisible in the limit, and a curved macroscopic host can carry any mesoscopic structure; so the motion constraints say
nothing about depth reduction (\cref{sec:depth}), which concerns all scales at once. In this sense the results of this
section are complementary to the decomposition of \cref{sec:decomposition}: they exclude classes of macroscopic motion
without any coding hypothesis, while the decomposition applies to every motion of an oblivious machine but needs depth
reduction and the crux.

\section{The decomposition: depth reduction plus an index-coding crux}\label{sec:decomposition}

This section reduces the oblivious case of~(T) to two statements of different kinds. For an oblivious machine and a
fixed input length, the run is a circuit: the value computed at each step is a fixed function of the input bit under
the input head, of the previous state, and of the symbols under the work heads, and each of these symbols was written
at the previous visit of its head to its cell. Revealing the values at a set $\Gamma$ of steps cuts this circuit. If
every path of the circuit that avoids $\Gamma$ is short, each output bit depends only on the revealed values and on
few input bits, so the revealed values form an index code for the transpose demand with small side information.

\Cref{sec:dec-dag} sets up the circuit and proves the cone lemma. \Cref{sec:dec-D} states the two hypotheses, depth
reduction at linear cost (depth-TCP) and a lower bound for index codes on the resulting side information (the crux),
and shows that together they exclude oblivious linear-time transposers (Theorem~D). \Cref{sec:dec-rates} records what
the crux must deliver. \Cref{sec:dec-B} proves depth reduction for walks of linear leg area, in particular for
home-based walks, for every coupling of the heads (Theorem~B), and \cref{sec:dec-corB} turns its cuts into genies,
which reduces home-based oblivious machines to the crux (Corollary~B). None of the results of this section is formalized; each comes with a complete proof.

\subsection{The dependency DAG of an oblivious run}\label{sec:dec-dag}

Let $M$ be a machine as in \cref{sec:prelim-machines}, with state set $\{0,\dots,q\}$ and $k\ge1$ work tapes over the
alphabet $\{0,\dots,g\}$, and let $L\ge1$ be an input length at which $M$ is oblivious, with halting time~$T$. On
every input $w\in\{0,1\}^L$ the machine executes steps $0,\dots,T-1$. Write $x_0(t)$ for the position of the input
head and $x_j(t)$ for the position of work head~$j$ ($1\le j\le k$) at time~$t$; by obliviousness these positions,
and the set of steps that append an output bit, do not depend on~$w$. Step~$t$ reads the symbols under the heads at
time~$t$, writes the cells $x_1(t),\dots,x_k(t)$, moves the heads, and possibly appends an output bit. As in
\cref{sec:prelim-walks}, $\prev_j(t)=\max\{s<t:x_j(s)=x_j(t)\}$ when this set is nonempty.

The \emph{value} $v_t(w)$ of step~$t$ on input~$w$ consists of the state after step~$t$, the $k$ symbols written at
step~$t$, and the output bit appended at step~$t$ (or a symbol for ``none''). It takes at most $2^{\beta(M)}$ values,
where
\[
  \beta(M)=\clog{(q+1)}+k\,\clog{(g+1)}+2
\]
is the \emph{value width} of~$M$; the last two bits encode the output (none, $0$ or $1$).

\begin{definition}[Dependency DAG]\label{def:depgraph}
The \emph{dependency DAG} $H_M$ of $M$ at length~$L$ has vertex set $[0,T)=\{0,\dots,T-1\}$ and the following edges,
all pointing backwards in time: a \emph{state edge} $t\to t-1$ for every $1\le t<T$, and a \emph{cell edge}
$t\to\prev_j(t)$, labelled~$j$, for every work tape~$j$ and every $t$ for which $\prev_j(t)$ is defined. The
\emph{length} of an edge $t\to s$ is $t-s$.
\end{definition}

Every vertex has out-degree at most $k+1$. For each~$j$ the map $t\mapsto\prev_j(t)$ is injective: if
$\prev_j(t)=\prev_j(t')=s$ with $t<t'$, then $x_j(t')=x_j(s)=x_j(t)$ with $s<t<t'$, so $\prev_j(t')\ge t>s$. Hence
the symbol written on tape~$j$ at step~$s$ is read along at most one edge, the state after step~$s$ only along
$s+1\to s$, and the output bit along none. The DAG depends only on the head schedule of~$M$ at length~$L$, not on the
input or on the transition table.

\begin{lemma}[Local recursion]\label{lem:dec-local}
\notformalized
For every $t<T$ there is a function $\Delta_t$ such that for every input $w\in\{0,1\}^L$,
\[
  v_t(w)=\Delta_t\bigl(v_{t-1}(w),\,w_{x_0(t)},\,v_{\prev_1(t)}(w),\dots,v_{\prev_k(t)}(w)\bigr),
\]
where an undefined argument is replaced by a constant: by the initial state if $t=0$, by the end marker if
$x_0(t)\notin[0,L)$, and by the blank symbol if $\prev_j(t)$ is undefined.
\end{lemma}

\begin{informalproof}
Step~$t$ applies the transition function to the state at time~$t$, the input symbol under the input head and the
$k$ symbols under the work heads, and $v_t(w)$ is read off from its result. The state at time~$t$ is the initial
state if $t=0$ and the state component of $v_{t-1}(w)$ otherwise. The input symbol is $w_{x_0(t)}$, or the end marker
if $x_0(t)\notin[0,L)$. Every step writes the cell under each work head, so at time~$t$ the cell under work head~$j$
holds the symbol written on tape~$j$ at step $\prev_j(t)$, a component of $v_{\prev_j(t)}(w)$, and it is blank if
the head has not visited it before time~$t$; this is the cell-contents lemma of the first paper, Lean-verified
there~\cite[Lemma~2.3]{FirstPaper2026}. The positions $x_0(t)$ and the steps $\prev_j(t)$ do not depend on~$w$, so neither
does~$\Delta_t$.
\end{informalproof}

\begin{definition}[Genies, cones and depth]\label{def:genie}
A \emph{genie} is a set $\Gamma\subseteq[0,T)$ of revealed steps. For an integer $\beta\ge\beta(M)$, its
\emph{revealed word} $G_\Gamma(w)=(v_s(w))_{s\in\Gamma}$ is written with $\beta$ bits per value; this defines a map
$G_\Gamma\colon\{0,1\}^L\to\{0,1\}^{\beta\card\Gamma}$, and $\beta\card\Gamma$ is the \emph{cost} of~$\Gamma$. A path
$t_0\to t_1\to\dots\to t_m$ in $H_M$ has \emph{length}~$m$, and it is \emph{$\Gamma$-free} if
$t_1,\dots,t_m\notin\Gamma$ (its first vertex may lie in~$\Gamma$). The \emph{cone} $C_\Gamma(t)$ is the set of
endpoints of $\Gamma$-free paths starting at~$t$, so $t\in C_\Gamma(t)$. Its \emph{boundary} $\partial_\Gamma(t)$ is
the set of $s\in\Gamma$ such that some edge of $H_M$ leads from a vertex of $C_\Gamma(t)$ to~$s$, and
$J_\Gamma(t)=x_0(C_\Gamma(t))\cap[0,L)$ is the set of input positions read in the cone. The genie has \emph{depth at
most~$D$} if every $\Gamma$-free path has length at most~$D$.
\end{definition}

So revealing the value $v_s$ blocks every edge into~$s$.

\begin{lemma}[Cone lemma]\label{lem:cone}
\notformalized
Let $\Gamma$ be a genie and $t<T$.
\begin{enumerate}[label=(\alph*)]
\item If two inputs $w,w'\in\{0,1\}^L$ agree on $J_\Gamma(t)$ and satisfy $v_s(w)=v_s(w')$ for every
$s\in\partial_\Gamma(t)$, then $v_t(w)=v_t(w')$. Thus $v_t$ is a function of the values at the boundary and of the
input bits read in the cone.
\item If $\Gamma$ has depth at most~$D$, then
$\card{J_\Gamma(t)}\le\card{C_\Gamma(t)}\le\sum_{i=0}^{D}(k+1)^i\le(k+1)^{D+1}$.
\end{enumerate}
\end{lemma}

\begin{informalproof}
(a) By strong induction on~$t$. Let $w,w'$ be as stated; we show that the arguments of $\Delta_t$ in
\cref{lem:dec-local} agree on $w$ and~$w'$. Constant arguments agree trivially. If $x_0(t)\in[0,L)$, then
$x_0(t)\in J_\Gamma(t)$ because $t\in C_\Gamma(t)$, so $w_{x_0(t)}=w'_{x_0(t)}$. Now let $t\to t'$ be an edge, with
$t'=t-1$ or $t'=\prev_j(t)$. If $t'\in\Gamma$, then $t'\in\partial_\Gamma(t)$, since the edge starts at
$t\in C_\Gamma(t)$, and $v_{t'}(w)=v_{t'}(w')$ by hypothesis. If $t'\notin\Gamma$, then every $\Gamma$-free path
from~$t'$, preceded by the edge $t\to t'$, is a $\Gamma$-free path from~$t$. Hence $C_\Gamma(t')\subseteq C_\Gamma(t)$,
so $J_\Gamma(t')\subseteq J_\Gamma(t)$ and $\partial_\Gamma(t')\subseteq\partial_\Gamma(t)$, and the induction
hypothesis for $t'<t$ gives $v_{t'}(w)=v_{t'}(w')$. So all arguments agree, and $v_t(w)=v_t(w')$.

(b) Every vertex of $C_\Gamma(t)$ is the endpoint of a $\Gamma$-free path from~$t$ of length at most~$D$, and since
out-degrees are at most $k+1$, at most $(k+1)^i$ paths of length~$i$ start at~$t$. The last inequality holds because
$k\ge1$.
\end{informalproof}

\begin{corollary}[Output bits]\label{cor:dec-output}
\notformalized
Suppose that $M$ appends exactly $L$ output bits at length~$L$, the $\ell$-th one ($0\le\ell<L$) at step~$t_\ell$,
and let $\Gamma$ be a genie. Put $I_\Gamma(\ell)=J_\Gamma(t_\ell)$. If $G_\Gamma(w)=G_\Gamma(w')$ and $w,w'$ agree on
$I_\Gamma(\ell)$, then output bit~$\ell$ is the same on $w$ and on~$w'$. If $\Gamma$ has depth at most~$D$, then
$\card{I_\Gamma(\ell)}\le(k+1)^{D+1}$.
\end{corollary}

\begin{informalproof}
Output bit~$\ell$ is a component of $v_{t_\ell}$, and $\partial_\Gamma(t_\ell)\subseteq\Gamma$, so
$G_\Gamma(w)=G_\Gamma(w')$ gives $v_s(w)=v_s(w')$ for every $s\in\partial_\Gamma(t_\ell)$. Apply \cref{lem:cone}.
\end{informalproof}

\begin{remark}[Field-level genies]\label{rem:dec-field}
A genie may also reveal single \emph{fields}: the state after a step, or the symbol written on one tape. Each field
is read along at most one edge, so revealing it blocks exactly that edge, and the proof of \cref{lem:cone} applies
verbatim to paths that avoid the blocked edges. We do not need this refinement. Below, an edge $t\to s$ is blocked by
putting~$s$ into~$\Gamma$, at a cost of $\beta$ bits, which also bounds the width of a single field.
\end{remark}

\subsection{Two hypotheses and Theorem~D}\label{sec:dec-D}

The first hypothesis, depth-TCP (TCP stands for \emph{trajectory cut problem}), says that a genie of linear cost
reduces the depth of every dependency DAG of linear size to a constant.

\begin{definition}[depth-TCP]\label{def:depth-tcp}
Let $k\ge1$, $D\ge0$ and $\beta\ge2$ be integers and $C,\eps>0$ reals. The statement
$\mathrm{depth}\text{-}\TCP(k,C,\beta,\eps,D)$ says: for every machine~$M$ with $k$ work tapes and value width
$\beta(M)\le\beta$ there is~$d_0$ such that, for every $d\ge d_0$ at which $M$ is oblivious at length $L=d^2b(d)$
with halting time $T\le C(L+1)$, the dependency DAG $H_M$ at length~$L$ has a genie $\Gamma$ of cost
$\beta\card\Gamma\le\eps L$ and depth at most~$D$.
\end{definition}

This is a statement about head schedules: whether a genie of a given size and depth exists depends only on the DAG.
By the trade-off of \cref{prop:tradeoff} below, applied to a machine that copies its input to its output, depth-TCP
with $C\ge1$ can hold only if $\eps\ge\beta/(D+1)$.

For the second hypothesis recall the transpose map. Bit~$s$ of entry $(i,j)$ of a $d\times d$ matrix with $b$-bit
entries sits at input position $(di+j)b+s$, and in the encoding of the transpose at position
$\operatorname{tpos}((di+j)b+s)=(dj+i)b+s$~\cite[Section~2.2 and Lemma~2.5]{FirstPaper2026}. The map $\operatorname{tpos}$ is an involution of
$[0,L)$, so on input~$w$ the output bit of a transposer at position~$\ell$ is $w_{\operatorname{tpos}(\ell)}$.

\begin{definition}[Index codes and the crux]\label{def:crux}
Let $d\ge1$ and $L=d^2b(d)$. A \emph{side-information system} at size~$d$ is a family $I=(I(\ell))_{\ell\in[0,L)}$
of subsets of $[0,L)$: receiver~$\ell$ knows the input bits at the positions in $I(\ell)$ and demands
$w_{\operatorname{tpos}(\ell)}$. A map $G\colon\{0,1\}^L\to\{0,1\}^m$ is an \emph{index code} for~$I$ with $m$ bits
if for every~$\ell$ and all $w,w'\in\{0,1\}^L$ with $G(w)=G(w')$ that agree on $I(\ell)$ we have
$w_{\operatorname{tpos}(\ell)}=w'_{\operatorname{tpos}(\ell)}$.

For a class $\mathcal S$ of side-information systems and a real $\rho\ge0$, the statement $\Crux(\mathcal S,\rho)$
says: for every $\eta>0$ there is~$d_0$ such that every index code for every system of $\mathcal S$ at a size
$d\ge d_0$ has at least $(\rho-\eta)L$ bits. In other words, index codes for systems of~$\mathcal S$ have at least
$\rho L-o(L)$ bits, with the $o(L)$ uniform over~$\mathcal S$.

The class $\mathcal S_{k,C,\beta,D}$ consists of the systems $I_\Gamma=(I_\Gamma(\ell))_{\ell\in[0,L)}$ of
\cref{cor:dec-output}, where $d\ge1$; $M$ is a machine with $k$ work tapes and $\beta(M)\le\beta$ that is oblivious
at length $L=d^2b(d)$ with halting time $T\le C(L+1)$ and appends exactly $L$ output bits; and $\Gamma$ is a genie
of depth at most~$D$ in~$H_M$.
\end{definition}

The class $\mathcal S_{k,C,\beta,D}$ depends on each machine only through its head schedule and its output times; no
transition table enters. Its sets have at most $(k+1)^{D+1}$ elements, and it grows with~$D$. The crux is a statement
about \emph{all} index codes, linear or not, and it must be: the revealed word of a genie is an arbitrary function of
the input, determined by the transition table. A lower bound for linear index codes alone would cover only machines
whose revealed words are linear functions of the input, such as machines with an $\F_2$-linear data path.

\begin{remark}[Two readings of the crux]\label{rem:dec-crux}
(a) \emph{The crux restricts head schedules.} Call receiver~$\ell$ of a system~$I$ \emph{trivial} if
$\operatorname{tpos}(\ell)\in I(\ell)$, and let $\mathrm{triv}(I)$ be the number of trivial receivers. Broadcasting
the demanded bits of the non-trivial receivers is an index code with $L-\mathrm{triv}(I)$ bits, so
$\Crux(\mathcal S,\rho)$ implies $\mathrm{triv}(I)\le(1-\rho)L+o(L)$ for the systems $I\in\mathcal S$. For
$\mathcal S_{k,C,\beta,D}$ this says that at most $(1-\rho)L+o(L)$ output positions~$\ell$ have the property that the input
head reads position $\operatorname{tpos}(\ell)$ at the end of a $\Gamma$-free path that starts at the step $t_\ell$ appending
output bit~$\ell$.

(b) \emph{Storage codes.} For a system~$I$ let $N_I(p)=I(\operatorname{tpos}(p))$, the side information
of the receiver that demands $w_p$. If $G$ is an index code for~$I$ with $m$ bits, then every fibre $G^{-1}(z)$
is a binary storage code for $N_I$ (\cref{def:storage-code}): two words of the fibre that agree on $N_I(p)$
agree at~$p$. Some fibre has at least $2^{L-m}$ words. Hence, if every binary storage code for every $N_I$,
$I\in\mathcal S$, has at most $2^{(1-\rho)L+o(L)}$ words, uniformly, then $\Crux(\mathcal S,\rho)$ holds; in
particular, if every $N_I$ satisfies the rate bound~$c$ (\cref{def:polymatroid}), then $\Crux(\mathcal S,1-c)$ holds
by \cref{thm:entropic-rate}. Conversely, a linear storage code $U\subseteq\F_2^L$ for $N_I$ of dimension $\lambda L$
gives the linear index code $w\mapsto w+U\in\F_2^L/U$ with $(1-\lambda)L$ bits: if $w-w'\in U$ vanishes on
$N_I(p)$, it agrees there with $0\in U$, so it vanishes at~$p$. Restricted to linear codes, the crux at
rate~$\rho$ is therefore equivalent to the statement that binary linear storage codes for the $N_I$, $I\in\mathcal S$,
have dimension at most $(1-\rho)L+o(L)$, uniformly (the kernel of a linear index code is a linear storage code). For
general codes the two notions are related by the duality recalled in \cref{sec:prelim-codes}.
\end{remark}

\begin{theorem}[Theorem D: decomposition]\label{thm:D}
\notformalized
Let $k\ge1$, $D\ge0$ and $\beta\ge2$ be integers, $C>0$, and $0<\eps<\rho$. Suppose that the statements
$\mathrm{depth}\text{-}\TCP(k,C,\beta,\eps,D)$ and $\Crux(\mathcal S_{k,C,\beta,D},\rho)$ hold. Then for every
machine~$M$ with $k$ work tapes and value width at most~$\beta$ there is~$d_0$ such that for no $d\ge d_0$ is $M$
oblivious at length $L=d^2b(d)$ with halting time at most $C(L+1)$ and at the same time a transposer of all
$d\times d$ matrices with $b(d)$-bit entries. In particular, no oblivious machine with $k$ work tapes and value width
at most~$\beta$ transposes within $C(L+1)$ steps for all~$d$, or even for infinitely many~$d$.
\end{theorem}

\begin{informalproof}
Fix~$M$. Let $d_1$ be given by $\mathrm{depth}\text{-}\TCP(k,C,\beta,\eps,D)$ for~$M$, let $d_2$ be given by
$\Crux(\mathcal S_{k,C,\beta,D},\rho)$ for $\eta=(\rho-\eps)/2$, and put $d_0=\max(d_1,d_2)$. Suppose that for some
$d\ge d_0$ the machine is oblivious at length $L=d^2b(d)$ with halting time $T\le C(L+1)$ and transposes all
$d\times d$ matrices. Every word of length~$L$ encodes such a matrix (\cite[Lemma~2.5]{FirstPaper2026}, Lean-verified there), so on every input~$w$
the output has length~$L$ and its bit~$\ell$ is $w_{\operatorname{tpos}(\ell)}$; in particular, $M$ appends exactly
$L$ output bits. By depth-TCP there is a genie~$\Gamma$ with $\beta\card\Gamma\le\eps L$ and depth at most~$D$, so
$I_\Gamma\in\mathcal S_{k,C,\beta,D}$. By \cref{cor:dec-output}, if $G_\Gamma(w)=G_\Gamma(w')$ and $w,w'$ agree on
$I_\Gamma(\ell)$, then $w_{\operatorname{tpos}(\ell)}=w'_{\operatorname{tpos}(\ell)}$: the revealed word $G_\Gamma$ is
an index code for $I_\Gamma$ with $\beta\card\Gamma$ bits. The crux gives
$\beta\card\Gamma\ge(\rho-\eta)L=\tfrac12(\rho+\eps)L>\eps L$, a contradiction.
\end{informalproof}

\begin{corollary}[The two hypotheses refute oblivious (T)]\label{cor:dec-obl}
\notformalized
Suppose that for all integers $k\ge1$ and $\beta\ge2$ and every real $C>0$ there are reals $0<\eps<\rho$ and an
integer $D\ge0$ such that $\mathrm{depth}\text{-}\TCP(k,C,\beta,\eps,D)$ and $\Crux(\mathcal S_{k,C,\beta,D},\rho)$
hold. Then for every oblivious machine~$M$ and every constant~$c$ there is~$d_0$ such that for no $d\ge d_0$ does $M$
transpose all $d\times d$ matrices within $c(L+1)$ steps. In particular, no oblivious machine witnesses~(T).
\end{corollary}

\begin{informalproof}
If $c=0$ there is nothing to prove, since a machine that halts at time~$0$ writes no output. A machine without work
tapes first receives one idle work tape, which keeps it oblivious and changes neither its running time nor its output.
So let $M$ have $k\ge1$ work tapes. It is oblivious at every input length, and if it transposes all $d\times d$ matrices
within $c(L+1)$ steps, its halting time at length $L=d^2b(d)$ is at most $c(L+1)$. Apply \cref{thm:D} with $C=c$ and
$\beta=\beta(M)\ge2$.
\end{informalproof}

\paragraph{Why it matters.} Theorem~D isolates exactly two statements that together refute the oblivious case
of~(T): a structural statement about tape dependency DAGs, which mentions neither the data nor the transition table,
and an information-theoretic statement about index codes, which sees a machine only through its head schedule and
its output times. An impossibility proof for oblivious machines can therefore be split into a depth-reduction part
and a coding part; for
instance, a crux at rate~$\tfrac12$ for every depth~$D$ would combine with depth reduction at a single
$\eps<\tfrac12$. Theorem~B and Corollary~B below settle the depth-reduction part for home-based heads; in general
both parts are open (\cref{sec:depth,sec:codes}). On the construction side, Theorem~D says what an oblivious
transposer~$M$ with $k$ work tapes and running time $C(L+1)$ must do; let $\beta=\beta(M)$. At every size at which it
transposes and for all $\eps>0$ and $D\ge0$, either $H_M$ has no genie of cost at most $\eps L$ and depth at
most~$D$, or the system $I_\Gamma\in\mathcal S_{k,C,\beta,D}$ of such a genie has an index code with at most $\eps L$
bits, namely~$G_\Gamma$; by \cref{rem:dec-crux}(b), some fibre of $G_\Gamma$ is a storage code of rate at least
$1-\eps$ for $N_{I_\Gamma}$.

Theorem~D is an elementary composition of the cone lemma with definitions tailored to it, the multitape analogue of
the route from small-depth computation to index coding by depth reduction recalled in the introduction. Its content is
the formulation of the two hypotheses and the bookkeeping of genie cost against depth.

\begin{remark}[Relation to the first paper]\label{rem:dec-first}
The first paper proves, Lean-verified, that the plain form of its Tape Coding Conjecture, which forbids causal codes
for complete exchange on the block networks of tape diagrams of linear length, at the profile capacities, refutes the
oblivious case of~(T)~\cite[Theorem~9.9]{FirstPaper2026}. Theorem~D is not a refinement of that reduction, and we know no
implication, in either direction, between the plain Tape Coding Conjecture and the conjunction of depth-TCP and the crux:
depth-TCP concerns the step-level dependency DAG and says nothing about codes on block networks, and the crux
concerns index codes for cone side information rather than complete exchange.
\end{remark}

\subsection{What the crux must deliver}\label{sec:dec-rates}

Theorem~D needs the crux at a rate~$\rho$ above the relative cost~$\eps$ of the genie. The next proposition bounds that
cost from below in terms of the depth.

\begin{proposition}[Forced trade-off]\label{prop:tradeoff}
\notformalized
Let $M$ be oblivious at length~$L$ with halting time~$T$, and let $\Gamma$ be a genie of depth at most~$D$. Then
$\card\Gamma\ge\lfloor(T-1)/(D+1)\rfloor$. If moreover $T\ge L$, as for every machine that appends $L$ output bits
and in particular for every transposer, and $\beta\card\Gamma\le\eps L$, then
\[
  \eps\;\ge\;\frac{\beta}{D+1}-\frac{\beta}{L},\qquad\text{equivalently}\qquad D+1\;\ge\;\frac{\beta L}{\eps L+\beta}.
\]
\end{proposition}

\begin{informalproof}
Let $n=\lfloor(T-1)/(D+1)\rfloor$ and $a_i=i(D+1)$ for $0\le i<n$. Then $a_i+D+1\le n(D+1)\le T-1$, and the path
$a_i+D+1\to a_i+D\to\dots\to a_i$ of state edges has length $D+1$, so it is not $\Gamma$-free: one of
$a_i,\dots,a_i+D$ lies in~$\Gamma$. These $n$ windows are disjoint, so $\card\Gamma\ge n\ge(T-1-D)/(D+1)$. With
$T\ge L$ and $\beta\card\Gamma\le\eps L$ this gives $\eps L\ge\beta(L-D-1)/(D+1)=\beta L/(D+1)-\beta$, which is
either form of the claim.
\end{informalproof}

\paragraph{Why it matters.} In Theorem~D the crux is used at a rate $\rho>\eps$, and for a transposer
\cref{prop:tradeoff} gives $\eps\ge\beta/(D+1)-\beta/L$. So at depth~$D$ the crux must deliver a rate of at least
about $\beta/(D+1)$, and since $\beta\ge2$, genies of cost $\eps L$ have depth $D\ge2/\eps-1-o(1)$: the cheaper the
genie, the deeper the side information that the crux has to handle. The next proposition shows that the size of the
cones alone never suffices.

\begin{proposition}[The counting bound cannot replace the crux]\label{prop:dec-counting}
\notformalized
Let $I$ be a side-information system at size~$d$ with $\card{I(\ell)}\le B$ for all~$\ell$. Call a set~$W$ of
non-trivial receivers (\cref{rem:dec-crux}(a)) \emph{acyclic} if there are no distinct $\ell_1,\dots,\ell_r\in W$, $r\ge2$, with
$\operatorname{tpos}(\ell_{i+1})\in I(\ell_i)$ for all~$i$ (indices modulo~$r$), and let $\MAIS(I)$ be the largest
size of an acyclic set.
\begin{enumerate}[label=(\alph*)]
\item Every index code for~$I$ has at least $\MAIS(I)$ bits.
\item $\MAIS(I)\ge(L-\mathrm{triv}(I))/(B+1)$.
\item Suppose that a machine~$M$ with $k$ work tapes, oblivious at length $L$ with halting time $T\ge L$, has a genie
$\Gamma$ of depth at most~$D$ and cost $\beta\card\Gamma\le\eps L$, where $\beta\ge2$, $0<\eps\le1$ and
$L\ge\beta/\eps$. Then $(k+1)^{D+1}>1/\eps$. Consequently, for every system~$I$ at this size whose sets have at most
$B=(k+1)^{D+1}$ elements, in particular for every system of $\mathcal S_{k,C,\beta,D}$ (\cref{cor:dec-output}), the
lower bound $(L-\mathrm{triv}(I))/(B+1)$ given by (a) and~(b) is smaller than $\eps L$; it never contradicts the
existence of an index code with $\eps L$ bits, such as the revealed word $G_\Gamma$ when $M$ is a transposer.
\end{enumerate}
\end{proposition}

\begin{informalproof}
(a) This is the acyclic-subgraph bound of Bar-Yossef et al.~\cite[Theorem~3]{BBJK2011}; we include the short proof in
our notation. Let $W$ be acyclic, and order it as
$\ell_1,\dots,\ell_r$ so that $\operatorname{tpos}(\ell_j)\in I(\ell_i)$ with $j\ne i$ implies $j<i$; this is
possible because the digraph on~$W$ with an arc from $\ell$ to $\ell'$ when $\operatorname{tpos}(\ell')\in I(\ell)$
has no cycle. Let $Z$ be the set of inputs that vanish outside $\operatorname{tpos}(W)$; it has $2^{\card W}$
elements because $\operatorname{tpos}$ is injective. We show that an index code~$G$ is injective on~$Z$. Let
$w,w'\in Z$ with $G(w)=G(w')$, and suppose $w$ and $w'$ agree at $\operatorname{tpos}(\ell_j)$ for all $j<i$. A
position in $I(\ell_i)$ lies outside $\operatorname{tpos}(W)$, where both inputs vanish, or equals
$\operatorname{tpos}(\ell_j)$ for some $j<i$, because $\operatorname{tpos}(\ell_i)\notin I(\ell_i)$ (receiver
$\ell_i$ is non-trivial). So $w$ and $w'$ agree on $I(\ell_i)$, hence at $\operatorname{tpos}(\ell_i)$. By induction
$w=w'$, so $2^m\ge2^{\card W}$.

(b) Order the non-trivial receivers uniformly at random, and let $W$ consist of those receivers~$\ell$ that precede
every non-trivial $\ell'\ne\ell$ with $\operatorname{tpos}(\ell')\in I(\ell)$. Then $W$ is acyclic, since along a
cycle in~$W$ each receiver would precede the next. As $\operatorname{tpos}$ is injective, a receiver~$\ell$ has at
most $\card{I(\ell)}\le B$ such receivers~$\ell'$, so it lies in~$W$ with probability at least $1/(B+1)$, and
$\mathbb E\card W\ge(L-\mathrm{triv}(I))/(B+1)$.

(c) By \cref{prop:tradeoff}, $D+1\ge\beta L/(\eps L+\beta)\ge\beta L/(2\eps L)=\beta/(2\eps)\ge1/\eps$, using
$\beta\le\eps L$ and $\beta\ge2$. Hence $(k+1)^{D+1}\ge2^{D+1}\ge2^{1/\eps}>1/\eps$, and
$(L-\mathrm{triv}(I))/(B+1)\le L/(B+1)<\eps L$.
\end{informalproof}

\paragraph{Why it matters.} The combined bound $(L-\mathrm{triv}(I))/(B+1)$ of (a) and~(b) uses only the sizes of the
sets of the side information and the number of trivial receivers. Part~(c) shows that it is always below $\eps L$: the
depth that a genie of cost $\eps L$ must reach is so large that the size bound $(k+1)^{D+1}$ on the cones is far bigger
than $1/\eps$. No argument that uses only the sizes can do better, because splitting the receivers into groups of
$B+1$, each receiver knowing the bits demanded by the others in its group, gives a system with sets of size~$B$ and an
index code of $\lceil L/(B+1)\rceil$ bits, one XOR per group. A proof of the crux must use the structure of the side
information produced by tape schedules.

\begin{remark}[The rate needed]\label{rem:dec-rates}
For each depth~$D$, let $\rho^*(D)$ be the supremum of all rates~$\rho$ such that the statement
$\Crux(\mathcal S_{k,C,\beta,D},\rho)$ holds. It is non-increasing in~$D$, because the classes grow with~$D$. Suppose that depth-TCP holds with parameters
$(k,C,\beta,\eps,D(\eps))$ for all small $\eps>0$. By Theorem~D, oblivious machines with these parameters are then
excluded as soon as $\rho^*(D(\eps))>\eps$ for a single~$\eps$. If $D(\eps)\le c\,\eps^{-a}$ for constants $c$
and~$a$ (by \cref{prop:tradeoff}, necessarily $a\ge1$ if depth-TCP is to apply to transposers), a crux with
$\rho^*(D)\ge D^{-1/a+\delta}$ for some $0<\delta<1/a$ suffices, because then
\[
  \rho^*(D(\eps))\;\ge\;\bigl(c\,\eps^{-a}\bigr)^{-1/a+\delta}\;=\;c^{-1/a+\delta}\,\eps^{1-a\delta}\;>\;\eps
\]
for all small~$\eps$. For heads of leg area at most $AT$ (home-based heads have $A=1$), Corollary~B below gives depth
$D(\eps)=2^{O(k^3(1+A)\beta C/\eps)}$; there a crux at a single depth~$D$ with a rate exceeding about
$456k^3(1+A)\beta C/\log_2D$ suffices. Both requirements are far
above the rate, at most $(k+1)^{-D-1}$, that the counting bound of \cref{prop:dec-counting} gives.
\end{remark}

\subsection{Depth reduction for home-based walks: Theorem~B}\label{sec:dec-B}

Theorem~B is a statement about walks; machines reappear in \cref{sec:dec-corB}. Throughout this subsection
$x_1,\dots,x_k$ are walks of length~$T$ (\cref{sec:prelim-walks}), indexed by \emph{heads} $h\in\{1,\dots,k\}$. They
may be coupled arbitrarily: no relation between them is assumed.

\paragraph{The walk DAG and cuts.} The \emph{walk DAG} of $x_1,\dots,x_k$ has vertex set $[0,T)$ and two kinds of
edges: \emph{state edges} $t\to t-1$ for $1\le t<T$, and \emph{long hops} $u\to\prev_h(u)$, one for every head~$h$
and every time~$u$ with $u-\prev_h(u)\ge2$ (the hop $(\prev_h(u),u)$ of \cref{sec:prelim-walks}, read backwards). A
pause $u\to u-1$ of a head joins the same two vertices as the state edge $u\to u-1$ and is identified with it; in
general an edge is determined by its two endpoints, so a pair that is a long hop of several heads is a single edge. A
\emph{cut} is a set~$F$ of edges, and it has \emph{depth at most~$D$} if every path that uses no edge of~$F$ contains
at most~$D$ long hops. For the work-head walks of an oblivious machine~$M$ at length~$L$, the walk DAG has the
vertices of $H_M$ (\cref{def:depgraph}), and every edge of $H_M$ joins the endpoints of an edge of the walk DAG: a
cell edge of length~$1$ is a pause, and a cell edge of length at least~$2$ is a long hop.

\paragraph{Bands, legs and leg area.} For $z\in\Z$ let $\lambda(z)=\lfloor\log_2\abs z\rfloor$ if $\abs z\ge2$ and
$\lambda(z)=0$ otherwise. The \emph{band} of head~$h$ at time~$t$ is $\lambda_h(t)=\lambda(x_h(t))$, a function of
the cell. Bands, scales, floors and the exponents~$\theta$ of radii $2^\theta$ are integers throughout. For integers
$\theta\ge1$, $\lambda(z)\ge\theta$ if and only if $\abs z\ge2^\theta$; and since $\abs{x_h}$ changes by at most~$1$ per
step, so does~$\lambda_h$. The \emph{scale} of a leg (\cref{sec:prelim-walks}) is the largest band over it. For
$\theta\ge0$ let $N_h(\theta,t)$ be the number of legs of head~$h$ of scale at least~$\theta$ that begin at or before
time~$t$, and let $N_h(\theta)=N_h(\theta,T-1)$. The \emph{leg area} of head~$h$ is
\[
  S_h=\sum_{\text{legs $\ell$ of $h$}}2^{\mathrm{scale}(\ell)},
\]
and the \emph{area hypothesis} with constant $A>0$ is
\[
  \text{(A)}\qquad S_h\le A\,T\quad\text{for every head }h.
\]
Every leg contributes at least~$1$, so head~$h$ has at most $S_h$ legs. The statements below use only that the legs
of each walk partition $[0,T)$ into intervals on which the walk is weakly monotone.

\begin{lemma}[Leg area of home-based walks]\label{lem:dec-area}
\notformalized
If every leg of a walk~$x$ of length~$T$ begins or ends at a cell~$z$ with $\abs z\le1$, then every leg~$\ell$
satisfies $2^{\mathrm{scale}(\ell)}\le\card\ell$, and the leg area of~$x$ is at most~$T$. In particular, home-based
walks satisfy~(A) with $A=1$.
\end{lemma}

\begin{informalproof}
If the scale~$s$ of $\ell$ is~$0$, then $2^s=1\le\card\ell$. Otherwise $\ell$ contains a time~$t$ with
$\abs{x(t)}\ge2^s$ and has an endpoint~$e$ with $\abs{x(e)}\le1$. Since $x$ moves by at most~$1$ per step,
$\abs{t-e}\ge2^s-1$, and the interval~$\ell$ contains the $\abs{t-e}+1\ge2^s$ times between $e$ and~$t$. Summing
over the legs, which partition $[0,T)$, bounds the leg area by~$T$.
\end{informalproof}

The walk DAG does not change when a single walk is translated, so everything below applies to heads with individual
fixed home cells, the bands of each head being measured from its home.

\begin{lemma}[Crossings and spanning hops]\label{lem:dec-facts}
\notformalized
Let $x$ be a walk of length~$T$ with leg area~$S$, let $\theta\ge1$, and let $N(\theta)$ be the number of legs of~$x$
of scale at least~$\theta$. A \emph{crossing} of radius~$2^\theta$ at time~$t$ ($1\le t<T$) is \emph{outward} if
$\abs{x(t-1)}<2^\theta\le\abs{x(t)}$ and \emph{inward} if $\abs{x(t-1)}\ge2^\theta>\abs{x(t)}$.
\begin{enumerate}[label=\textup{(F\arabic*)}]
\item The walk has at most $2N(\theta)$ outward crossings, and at most $5N(\theta)$ crossings in total, of
radius~$2^\theta$. Moreover $\sum_{\theta\ge1}2^\theta N(\theta)\le2S$.
\item For every time~$\tau$ and every integer $R\ge1$, at most $2R-1$ hops $u\to\prev(u)$ of~$x$ \emph{span}~$\tau$,
that is, satisfy $\prev(u)<\tau\le u$, and have $\abs{x(u)}<R$.
\end{enumerate}
\end{lemma}

\begin{informalproof}
(F1) Assign each outward crossing~$t$ to the leg~$\ell$ that contains~$t$; the scale of~$\ell$ is at least~$\theta$
because $\abs{x(t)}\ge2^\theta$. Either $t$ is the first time of~$\ell$, which happens for at most one crossing per
leg, or $t-1\in\ell$ as well. In the second case suppose that $x$ is non-decreasing on~$\ell$; the non-increasing case
follows by the symmetry $z\mapsto-z$. If $x(t)\le0$, then $x(t-1)\le x(t)\le0$ and $\abs{x(t-1)}\ge\abs{x(t)}$,
which is impossible; so $x(t)\ge2^\theta$, and $x\ge2^\theta$ at all later times of~$\ell$, so no later outward
crossing $t'$ with $t'-1\in\ell$ is assigned to~$\ell$. Thus each leg of scale at least~$\theta$ receives at most two
outward crossings and every other leg none, which gives the first bound. Crossings of radius~$2^\theta$ alternate
between outward and inward, as each changes whether $\abs x\ge2^\theta$. So the inward crossings number at most the
outward ones plus one, and the extra one occurs only if $\abs{x(0)}\ge2^\theta$, in which case the first leg has
scale at least~$\theta$ and $N(\theta)\ge1$. Altogether there are at most $2N(\theta)+3N(\theta)$ crossings. Finally,
$\sum_{\theta\ge1}2^\theta N(\theta)=\sum_\ell\sum_{1\le\theta\le\mathrm{scale}(\ell)}2^\theta
\le\sum_\ell2^{\mathrm{scale}(\ell)+1}=2S$.

(F2) For such a hop, $x$ visits the cell $x(u)$ at time $\prev(u)<\tau$ and at no time strictly between $\prev(u)$
and~$u$, so $u$ is the first visit to $x(u)$ at or after~$\tau$. Hence each of the $2R-1$ cells~$z$ with $\abs z<R$
is the cell of at most one such hop.
\end{informalproof}

\begin{lemma}[Legs and long hops]\label{lem:legs}
\notformalized
Let $u\to p$ be a long hop of head~$h$, so $p=\prev_h(u)\le u-2$. The leg of~$h$ that contains~$u$ begins after~$p$
and has scale at least $\lambda_h(u)$. Hence $N_h(\theta,u)\ge N_h(\theta,p)+1$ for every $\theta\le\lambda_h(u)$.
\end{lemma}

\begin{informalproof}
Pick~$t$ with $p<t<u$. By the definition of $\prev_h$, head~$h$ is not at the cell $x_h(u)=x_h(p)$ at time~$t$. If
$p$ and~$u$ lay in one leg, $x_h$ would be weakly monotone on an interval containing $p<t<u$, so $x_h(t)$ would lie
between the equal values $x_h(p)$ and $x_h(u)$, a contradiction. So the leg containing~$u$ is a later interval than
the leg containing~$p$, and it begins after~$p$; its scale is at least $\lambda_h(u)$ because it contains~$u$. For
$\theta\le\lambda_h(u)$ this leg is counted in $N_h(\theta,u)$ but not in $N_h(\theta,p)$, and every leg counted in
$N_h(\theta,p)$ is counted in $N_h(\theta,u)$.
\end{informalproof}

\Cref{lem:legs} fails for pauses, whose two endpoints may lie in one leg; this is why pauses are treated as state edges.

\begin{definition}[The band label and its cut]\label{def:dec-label}
Fix integers $a\ge1$ (the \emph{group width}), $r\in\{0,\dots,a-1\}$ (the \emph{offset}), $K\ge2$ (the
\emph{modulus}) and $\omega_1,\dots,\omega_k\in\{0,\dots,K-1\}$.
\begin{itemize}[leftmargin=2em]
\item The \emph{group} of head~$h$ at time~$t$ is $g_h(t)=\lfloor(\lambda_h(t)+r)/a\rfloor$, and the \emph{floor} of
a group $G\ge0$ is $\theta(G)=\max(aG-r,0)$, the smallest band in group~$G$. The floors of the groups $G\ge1$ are the
integers $\theta\ge1$ with $\theta\equiv-r\pmod a$.
\item The \emph{occupied groups} at time~$t$ are the distinct values among $g_1(t),\dots,g_k(t)$, listed from the top
as $G_1(t)>G_2(t)>\dots>G_{m(t)}(t)$, with floors $\theta_i(t)=\theta(G_i(t))$; here $m(t)\le k$.
\item The \emph{shared counts} are $n_i(t)=\sum_{h=1}^kN_h(\theta_i(t),t)$ for $i\le m(t)$: the number of
legs, of all heads, of scale at least $\theta_i(t)$ begun by time~$t$.
\item The \emph{label} is
\[
  \Psi(t)=\sum_{i=1}^{m(t)}\bigl((n_i(t)+\omega_i)\bmod K\bigr)\,K^{k-i}\;\in[0,K^k).
\]
\item The \emph{cut} $F=F(a,r,K,\omega)$ consists of the long hops $u\to p$ with $\Psi(p)\ge\Psi(u)$ and the state
edges $t\to t-1$ with $\Psi(t-1)>\Psi(t)$.
\end{itemize}
\end{definition}

So $\Psi(t)$ is the integer whose base-$K$ digits, most significant first, are $\psi_i(t)=(n_i(t)+\omega_i)\bmod K$
for $i\le m(t)$ and $0$ for $m(t)<i\le k$, and labels compare as these digit strings compare lexicographically.
Whether an edge is cut depends only on its endpoints. For a fixed~$\theta$, the shared count $\sum_hN_h(\theta,t)$ is
non-decreasing in~$t$; this monotonicity and \cref{lem:legs} are what the label uses.

\begin{theorem}[Theorem B: depth reduction for home-based walks]\label{thm:B}
\notformalized
Let $x_1,\dots,x_k$ be walks of length~$T$, coupled arbitrarily, that satisfy the area hypothesis~(A) with
constant~$A$. Let $a\ge1$ and $K\ge2$ be integers, and let $r\in\{0,\dots,a-1\}$ and $\omega\in\{0,\dots,K-1\}^k$ be
uniformly random and independent. Then the cut $F=F(a,r,K,\omega)$ has depth at most $K^k-1$, and
\begin{equation}\label{eq:dec-B}
  \mathbb E_{r,\omega}\card F\;\le\;\frac{(20k^2A+k)\,T}{a}+\Bigl(9+\frac{2^{a+2}}a\Bigr)\frac{k^2AT}{K}
  \;\le\;20k^2(1+A)\,T\Bigl(\frac1a+\frac{2^a}{aK}\Bigr)+\frac{20k^2AT}{K}.
\end{equation}
Consequently:
\begin{enumerate}[label=(\alph*)]
\item for every $D\ge4^k$, the choice $K=\lfloor D^{1/k}\rfloor$ and $a=\lfloor\log_2D/(2k)\rfloor$, with suitable
offsets, gives a cut of depth at most $D-1$ with $\card F\le114k^3(1+A)T/\log_2D$;
\item for every $0<\eps\le1$, some cut of size at most $\eps T$ has depth less than
$2^{\lceil114k^3(1+A)/\eps\rceil}$;
\item in particular, for every coupling of $k$ home-based walks, $\eps T$ edges suffice to reach depth less than
$2^{\lceil228k^3/\eps\rceil}=2^{O(k^3/\eps)}$.
\end{enumerate}
\end{theorem}

\begin{informalproof}
\emph{Depth.} Along a path that avoids~$F$, $\Psi$ does not increase along state edges and strictly decreases along
long hops. Since $\Psi$ takes values in $[0,K^k)$, the path contains at most $K^k-1$ long hops.

\emph{Notation.} Fix the walks, $a$ and~$K$; probabilities and expectations are over $r$ and~$\omega$. We bound the
number of cut long hops counted with their head, as pairs $(h,u)$, which is at least the number of cut long-hop
edges, and the number of cut state edges. Consider a long hop $u\to p$ of head~$h$, with cell $z=x_h(u)=x_h(p)$.
Since bands are functions of the cell, $h$ has the same band $\lambda=\lambda(z)$ and the same group
$G=g_h(u)=g_h(p)$ at $p$ and at~$u$. Let $\theta^+=\theta(G+1)=a(G+1)-r\ge1$ be the floor of the next group: a head
is in a group above~$G$ at time~$t$ exactly when its cell has absolute value at least $2^{\theta^+}$, and
$\abs z<2^{\theta^+}$. Let $U(t)$ be the set of occupied groups above~$G$ at time~$t$. Whether $U(p)=U(u)$ holds,
and all counts $n_i$, depend on~$r$ only.

\emph{Step 1: $U(p)=U(u)$.} Then $G$ has the same position $j=\card{U(p)}+1$ in the lists at $p$ and at~$u$, and
$G_i(p)=G_i(u)$, with a common floor~$\theta_i$, for every $i\le j$. Let $\Delta_i=n_i(u)-n_i(p)$ for $i\le j$. For
$i<j$, $\Delta_i\ge0$ by monotonicity at a fixed floor. For $i=j$ we have $\theta_j=\theta(G)\le\lambda=\lambda_h(u)$,
so \cref{lem:legs} gives $N_h(\theta_j,u)\ge N_h(\theta_j,p)+1$, and with monotonicity for the other heads, $\Delta_j\ge1$.
Say that digit~$i$ \emph{wraps} if $\lfloor(n_i(p)+\omega_i)/K\rfloor\ne\lfloor(n_i(u)+\omega_i)/K\rfloor$. If no
digit $i\le j$ wraps, then $\psi_i(p)=\psi_i(u)-\Delta_i$ for every $i\le j$, so the first position at which the
digit strings of $p$ and~$u$ differ is at most~$j$ and carries the smaller digit at~$p$. Thus $\Psi(p)<\Psi(u)$, and
the hop is not cut; the digits below position~$j$ play no role. Given~$r$, since $\omega_i$ is uniform on
$\{0,\dots,K-1\}$, digit~$i$ wraps with probability at most $\Delta_i/K$. So, given~$r$, a hop of Step~1 is cut with
probability at most $\sum_{i\le j}\Delta_i/K$.

\emph{Step 2: $U(p)\ne U(u)$.} We count every such hop as cut and charge it to a crossing. Let $G'>G$ be a group that
lies in exactly one of $U(p)$ and $U(u)$. Since $h$ is in group~$G$ at both times, at one of the two times $G'$ is
occupied by a head $h'\ne h$ that is not in~$G'$ at the other time. Groups change by at most one per step, and a
head moves from a group $G''$ to $G''+1$ exactly when its band reaches $\theta(G''+1)$, that is, at an outward
crossing of radius $2^{\theta(G''+1)}$; it moves back at an inward crossing of the same radius. If $h'$ is in~$G'$
at time~$p$, consider its first move out of~$G'$ after~$p$; if $h'$ is in~$G'$ at time~$u$, consider its last move
into~$G'$ at or before~$u$. In either case the move happens at a step $\tau-1\to\tau$ with $p<\tau\le u$, and it is a
crossing of radius $2^{\theta(G'+1)}$ or $2^{\theta(G')}$. So $h'$ has a crossing of radius $2^{\theta^*}$ at a time
$\tau\in(p,u]$, where $\theta^*$ is a floor with $\theta^*\ge\theta(G')\ge\theta^+$.
Charge the hop to one such crossing $(h',\theta^*,\tau)$, chosen by a fixed rule. Every hop $(h,u)$ charged to
$(h',\theta^*,\tau)$ satisfies $\prev_h(u)<\tau\le u$ and $\abs{x_h(u)}=\abs z<2^{\theta^+}\le2^{\theta^*}$, so by
(F2) with $R=2^{\theta^*}$, applied to each of the $k$ heads, at most $k\,2^{\theta^*+1}$ hops are charged to each
crossing. By (F1), head~$h'$ has at most $5N_{h'}(\theta)$ crossings of radius~$2^\theta$. Hence, for fixed~$r$,
\[
  \#\{\text{hops of Step 2}\}\;\le\;\sum_{h'=1}^k\ \sum_{\theta\ge1\text{ a floor}}5N_{h'}(\theta)\cdot k\,2^{\theta+1}.
\]
Each $\theta\ge1$ is a floor for exactly one value of~$r$, hence with probability $1/a$, and by (F1) and~(A),
\[
  \mathbb E\,\#\{\text{hops of Step 2}\}\;\le\;\frac{10k}a\sum_{h'}\sum_{\theta\ge1}2^\theta N_{h'}(\theta)
  \;\le\;\frac{10k}a\sum_{h'}2S_{h'}\;\le\;\frac{20k^2AT}a.
\]

\emph{Step 3: wraps.} By Step~1, the expected number of cut hops of Step~1 is at most
$\frac1K\mathbb E\sum\sum_{i\le j}\Delta_i$, where the outer sum runs over the hops of Step~1. For such a hop and
$i\le j$, $\Delta_i$ is the number of legs, of all heads, that have scale at least~$\theta_i$ and begin in $(p,u]$.
We exchange the order of summation: for each leg~$\ell$, of any head, beginning at time~$\tau_\ell$ with
scale~$s_\ell$, we count the hops that span~$\tau_\ell$ in the sense of (F2), that is, satisfy $p<\tau_\ell\le u$. By the argument
of (F2), for each head and each cell at most one hop of that head spans~$\tau_\ell$ and ends at that cell.

\emph{Own digit ($i=j$).} A hop spanning~$\tau_\ell$ contributes~$\ell$ to its $\Delta_j$ only if $\theta(G)\le
s_\ell$, where $G$ is its group. Since $\theta(G)$ is the largest floor at most~$\lambda$ (or~$0$ if there is none),
this holds if $\lambda\le s_\ell$, and for $\lambda>s_\ell$ it holds exactly when no floor lies in $(s_\ell,\lambda]$.
The floors at least~$1$ form an arithmetic progression of difference~$a$ with a uniformly random offset, so the latter
has probability $\max(0,1-(\lambda-s_\ell)/a)$. There are $2^{s+2}-1$ cells of band at most~$s$, and $2^{\lambda+1}$
cells of band $\lambda\ge1$. So the expected number of hops counted for~$\ell$ is at most
\[
  k\Bigl(2^{s_\ell+2}+\sum_{i=1}^{a-1}2^{s_\ell+i+1}\Bigl(1-\frac ia\Bigr)\Bigr)
  \;\le\;k\Bigl(2^{s_\ell+2}+\frac{2^{s_\ell+a+2}}a\Bigr),
\]
because $\sum_{i=1}^{a-1}2^i(a-i)=2^a\sum_{j=1}^{a-1}j2^{-j}<2^{a+1}$. Summing over the legs of all heads, with
$\sum_\ell2^{s_\ell}=\sum_{h'}S_{h'}\le kAT$, gives $\mathbb E\sum\Delta_j\le4k^2AT(1+2^a/a)$.

\emph{Higher digits ($i<j$).} Here $G_i>G$, so $\theta_i\ge\theta^+$ is a floor at least~$1$ that exceeds the band
of the hop, and distinct~$i$ give distinct floors. So, for each leg~$\ell$, the pairs consisting of a hop
spanning~$\tau_\ell$ and an index $i<j$ with $\theta_i\le s_\ell$ number at most the sum, over the floors
$1\le\theta\le s_\ell$, of the number of hops spanning~$\tau_\ell$ whose cell has absolute value less than~$2^\theta$.
By (F2) this number is at most $k(2^{\theta+1}-1)$, and since consecutive floors differ by $a\ge1$, the sum is at most
$k\sum_{\theta=1}^{s_\ell}2^{\theta+1}<k\,2^{s_\ell+2}$. Summing over the legs gives
$\sum\sum_{i<j}\Delta_i\le4k^2AT$ for every~$r$.

Together, the expected number of cut hops of Step~1 is at most
$(4k^2AT(1+2^a/a)+4k^2AT)/K=(8+2^{a+2}/a)\,k^2AT/K$.

\emph{Step 4: state edges.} Consider the state edge $t\to t-1$. If no head changes group at step $t-1\to t$, the
occupied groups and their floors are the same at $t-1$ and~$t$, and $n_i(t-1)\le n_i(t)$ for every~$i$; as in
Step~1, the edge is cut only if some digit wraps, which, given~$r$, has probability at most
$\sum_i(n_i(t)-n_i(t-1))/K$. Here $n_i(t)-n_i(t-1)$ is the number of heads that begin, at time~$t$, a leg of scale at
least~$\theta_i$. A head begins at most one leg at a time and there are at most $k$ digits, so the sum over $i$
and~$t$ is at most $k$ times the number of legs, which is at most $\sum_{h'}S_{h'}\le kAT$; so these cuts number at
most $k^2AT/K$ in expectation. Otherwise some head changes group at step $t-1\to t$, and we count the edge as cut. A
head changes group at a step only if its band moves between two values $\lambda$ and $\lambda+1$ and $\lambda+1$ is a
floor, which has probability $1/a$; so the expected number of such steps is at most $kT/a$.

\emph{Total.} Adding Steps 2, 3 and~4,
\[
  \mathbb E\card F\;\le\;\frac{20k^2AT}a+\Bigl(8+\frac{2^{a+2}}a\Bigr)\frac{k^2AT}K+\frac{k^2AT}K+\frac{kT}a,
\]
which is the first bound of \eqref{eq:dec-B}; the second follows from $k\le20k^2$, $9\le20$ and
$2^{a+2}A\le20(1+A)2^a$. Some offsets $r,\omega$ achieve at most the expectation.

\emph{Parameters.} For~(a) put $y=\log_2D/(2k)\ge1$. Then $a=\lfloor y\rfloor\ge y/2$, and $K=\lfloor2^{2y}\rfloor
\ge2^{2y-1}$ because $2^{2y}\ge4$; moreover $K^k\le D$. Using $2^{-y}\le1/y$ for $y>0$,
\[
  \frac1a\le\frac2y=\frac{4k}{\log_2D},\qquad\frac{2^a}K\le2^{1-y}\le\frac2y=\frac{4k}{\log_2D},\qquad
  \frac1K\le2^{1-2y}\le2^{-y}\le\frac1y=\frac{2k}{\log_2D}.
\]
Inserting these into the first bound of \eqref{eq:dec-B}, with $2^{a+2}/(aK)\le4\cdot2^a/K$, gives
\[
  \mathbb E\card F\;\le\;\bigl(80k^3A+4k^2+18k^3A+16k^3A\bigr)\frac T{\log_2D}\;\le\;\frac{114k^3(1+A)\,T}{\log_2D},
\]
and the depth is at most $K^k-1\le D-1$. For~(b) take $\log_2D=\lceil114k^3(1+A)/\eps\rceil$, which is at least~$2k$.
Part~(c) is~(b) with $A=1$, by \cref{lem:dec-area}.
\end{informalproof}

\paragraph{Why it matters.} Theorem~B proves the walk form of the first hypothesis of Theorem~D for all heads of
linear leg area, in particular for home-based heads, and for every coupling: it makes no assumption on hop lengths, on
independence, or on how the heads move relative to each other. So, for an oblivious machine running in time at most
$C(L+1)$, a genie of cost $\eps L$ and depth at most $D(\eps)$, depending on $k$, $A$, $\beta$ and~$C$, exists for every
$\eps\in(0,1]$ (see the paragraph after the proof of \cref{cor:B}) unless the leg area of some work head is not $O(T)$
for any choice of home cell. That requires many legs whose duration is much shorter than
$2^{\mathrm{scale}}$; by \cref{lem:dec-area} such legs neither begin nor end near the home cell, so the head turns
around far from every fixed home. \Cref{sec:depth} studies walks of this kind.

\subsection{Home-based oblivious machines reduce to the crux}\label{sec:dec-corB}

Theorem~B cuts edges, counts only long hops, and pays one unit per cut edge. Two conversions turn its cut into a genie
in the sense of \cref{def:genie}: an edge $t\to s$ is blocked by revealing~$s$, at a cost of $\beta$ bits, and the
runs of short edges are cut by revealing periodic steps.

\begin{lemma}[Long-edge conversion]\label{lem:dec-longedge}
\notformalized
Let $M$ be oblivious at length~$L$ with halting time~$T$, let $\mu\ge1$ and $P>\mu$ be integers, and let $\Gamma_0$
be a genie such that every $\Gamma_0$-free path has at most $D'$ edges of length greater than~$\mu$. Then the genie
\[
  \Gamma=\Gamma_0\cup\{t\in[0,T):t\bmod P\ge P-\mu\}
\]
has depth at most $D'+(D'+1)(P-\mu)\le(D'+1)P-1$, and $\card\Gamma\le\card{\Gamma_0}+\mu\lceil T/P\rceil$; for
$\mu=1$, $\card\Gamma\le\card{\Gamma_0}+\lfloor T/P\rfloor$.
\end{lemma}

\begin{informalproof}
Let $t_0\to t_1\to\dots\to t_m$ be $\Gamma$-free. It is $\Gamma_0$-free, so at most $D'$ of its edges have length
greater than~$\mu$, and removing them leaves at most $D'+1$ maximal runs of consecutive edges of length at most~$\mu$.
Let $t_a\to\dots\to t_b$ be such a run. Its vertices $t_{a+1}>\dots>t_b$ are not in~$\Gamma$, so each lies in an
interval $[nP,nP+P-\mu)$ with $n\ge0$. If $t_i$ lies in such an interval and $t_{i+1}<nP$ for consecutive
vertices of the run, then $n\ge1$ and $t_{i+1}\ge t_i-\mu\ge nP-\mu$, so $t_{i+1}\bmod P\ge P-\mu$ and
$t_{i+1}\in\Gamma$, which is impossible. Hence $t_{a+1},\dots,t_b$ are distinct elements of one interval of $P-\mu$
consecutive integers, and the run has $b-a\le P-\mu$ edges. So $m\le D'+(D'+1)(P-\mu)$, and this is at most
$(D'+1)P-1$ as $\mu\ge1$. Each of the $\lceil T/P\rceil$ periods $[nP,(n+1)P)$ that meet $[0,T)$ contains $\mu$
revealed residues; for $\mu=1$ the revealed steps are $P-1,2P-1,\dots$, of which $\lfloor T/P\rfloor$ are less
than~$T$.
\end{informalproof}

\begin{remark}[Why blocks of $\mu$ steps]\label{rem:dec-jitter}
For $\mu=1$ one revealed step per period suffices, because a run of edges of length~$1$ visits consecutive times.
Single revealed steps do not cut runs of edges of length at most~$\mu$ when $\mu\ge2$: a work head that oscillates
between two adjacent cells has $\prev_j(t)=t-2$ for all $t\ge2$, and if $P$ is even, the path $t\to t-2\to t-4\to
\cdots$ through even times avoids every step $t\equiv P-1\pmod P$. A depth-reduction statement that treats such short
revisits (jitter, or revisits within a window of $\mu$ steps) as free must therefore be combined with
\cref{lem:dec-longedge} for that~$\mu$, which reveals blocks of $\mu$ consecutive steps at $\mu$ times the cost.
Theorem~B needs only $\mu=1$, because it counts every hop of length at least~$2$ as long.
\end{remark}

\begin{corollary}[Corollary B: home-based oblivious machines reduce to the crux]\label{cor:B}
\notformalized
Let $k\ge1$ and $\beta\ge2$ be integers, $C\ge1$, $A>0$ and $0<\eps\le1$, and put
\[
  D=D(\eps)=2^{\lceil456k^3(1+A)\beta C/\eps\rceil}\cdot\lceil4\beta C/\eps\rceil=2^{O(k^3(1+A)\beta C/\eps)}.
\]
Suppose that $\Crux(\mathcal S_{k,C,\beta,D},\rho)$ holds for some $\rho>\eps$. Then for every machine~$M$ with $k$
work tapes and value width at most~$\beta$ there is~$d_0$ such that for no $d\ge d_0$ do the following three
conditions hold together at $L=d^2b(d)$: $M$ is oblivious at length~$L$ with halting time $T\le C(L+1)$; each
work-head walk~$x_j$ of~$M$ on $[0,T)$ has leg area at most $AT$ measured from some cell~$c_j$ (that is, the walk
$x_j-c_j$ satisfies~(A)); and $M$ transposes all $d\times d$ matrices with $b(d)$-bit entries. In particular, if an
oblivious machine with $k$ work tapes and value width at most~$\beta$, whose work heads are home-based at every input
length, transposes within $C(L+1)$ steps for infinitely many~$d$, then for every $\eps\in(0,1]$ the crux
$\Crux(\mathcal S_{k,C,\beta,D(\eps)},\rho)$, with $A=1$ in $D(\eps)$, fails for every rate $\rho>\eps$.
\end{corollary}

\begin{informalproof}
Put $D_1=2^{\lceil456k^3(1+A)\beta C/\eps\rceil}$ and $P=\lceil4\beta C/\eps\rceil$, so that $D=D_1P$; note that
$D_1\ge4^k$, because $456k^3(1+A)\beta C/\eps\ge2k$. Fix~$M$, let $d_0$ be given by
$\Crux(\mathcal S_{k,C,\beta,D},\rho)$ for $\eta=(\rho-\eps)/2$, and suppose that the three conditions hold at some
$d\ge d_0$. As in the proof of \cref{thm:D}, $M$ appends exactly $L$ output bits, and on every input~$w$ its output
bit~$\ell$ is $w_{\operatorname{tpos}(\ell)}$.

Apply \cref{thm:B}(a) with depth parameter~$D_1$ to the translated work-head walks $x_1-c_1,\dots,x_k-c_k$ on
$[0,T)$, whose walk DAG is that of $x_1,\dots,x_k$: some cut~$F$ of depth at most $D_1-1$ has
$\card F\le114k^3(1+A)T/\log_2D_1\le\eps T/(4\beta C)$. Let $\Gamma_0$ be the set of
steps~$s$ such that some edge of~$F$ ends at~$s$; so $\card{\Gamma_0}\le\card F$. Every edge of a $\Gamma_0$-free path
in $H_M$ ends outside~$\Gamma_0$, so its endpoints are not those of an edge of~$F$. Read in the walk DAG, the path
therefore avoids~$F$, and its edges of length at least~$2$ are long hops; so it has at most $D_1-1$ edges of length
greater than~$1$. By \cref{lem:dec-longedge} with $\mu=1$, the genie
$\Gamma=\Gamma_0\cup\{t<T:t\equiv P-1\pmod P\}$ has depth at most $D_1P-1<D$, and
\[
  \beta\card\Gamma\;\le\;\beta\card F+\beta\lfloor T/P\rfloor\;\le\;\frac{\eps T}{4C}+\frac{\eps T}{4C}
  \;\le\;\frac{\eps(L+1)}2\;\le\;\eps L,
\]
using $T\le C(L+1)$ and $L\ge1$. So $I_\Gamma\in\mathcal S_{k,C,\beta,D}$, and by \cref{cor:dec-output} the revealed
word $G_\Gamma$ is an index code for $I_\Gamma$ with at most $\eps L$ bits, while the crux requires at least
$(\rho-\eta)L>\eps L$ bits. This contradiction proves the first statement. For the second, a work head that is
home-based with home cell~$c_j$ satisfies~(A) with $A=1$ when measured from~$c_j$, by \cref{lem:dec-area} applied to
$x_j-c_j$; a machine that transposes within $C(L+1)$ steps at length~$L$ has halting time at most $C(L+1)$ there; so a
machine as described would violate the first statement at infinitely many~$d$ if the crux held at some rate
$\rho>\eps$.
\end{informalproof}

The genie~$\Gamma$ built in this proof does not use that $M$ transposes. It shows that the conclusion of
$\mathrm{depth}\text{-}\TCP(k,C,\beta,\eps,D(\eps))$ holds for every machine~$M$ with $k$ work tapes and value width at
most~$\beta$, at every length $L=d^2b(d)$ at which $M$ is oblivious with halting time at most $C(L+1)$ and its
work-head walks satisfy~(A), measured from some home cells.

\paragraph{Why it matters.} Theorem~D needs two hypotheses. For home-based heads Theorem~B supplies the first, so
Corollary~B reduces home-based oblivious machines to the crux alone, whatever the coupling of their heads. Read as a
requirement on the crux, it suffices to prove $\Crux(\mathcal S_{k,C,\beta,D},\rho)$ at a single depth $D=D(\eps)$
with a rate $\rho>\eps$, that is, a rate exceeding about $456k^3(1+A)\beta C/\log_2D$ at depth~$D$ (\cref{rem:dec-rates});
this is far beyond the rate, at most $(k+1)^{-D-1}$, that the counting bound of \cref{prop:dec-counting} gives, but, for large~$D$, far below
rate~$\tfrac12$. Read as a design constraint, an oblivious linear-time transposer must either have a work head whose
leg area, measured from any home cell, is not $O(T)$, or produce, for every $\eps\in(0,1]$ and at every size at which it transposes, an index
code with at most $\eps L$ bits for its own cone side information at depth $D(\eps)$, and hence a storage code of
rate at least $1-\eps$ (\cref{rem:dec-crux}(b)).

\begin{remark}[Adaptive machines]\label{rem:dec-adaptive}
The argument uses obliviousness twice: the dependency DAG and the times of the output bits must be the same for all
inputs. We do not treat adaptive machines here. We expect that an extension of Corollary~B to adaptive machines whose heads are
home-based on every input would need additional hypotheses on head motion far from home and a stronger form of the
crux, for index codes that need to be correct only on a large set of inputs. No result of this paper depends on such an extension.
\end{remark}

\section{Depth reduction for interacting walks}\label{sec:depth}

\Cref{thm:D} reduces the oblivious case of~(T) to two hypotheses. This section is about the structural one: deleting a
small fraction of the edges of the walk DAG of a head schedule so that every remaining path contains few long hops.
\Cref{thm:B} does this for home-based walks, for every coupling. Here we state the general problem at a fixed
fraction~$\eps$ (\cref{prob:tcp}), rule out packing certificates for lower bounds (\cref{prop:packing}), give
evidence against per-head band bookkeeping (\cref{rem:dep-N,rem:dep-Nprime}, proof sketches only), and study what we
regard as the sharpest open instance: two deeply nested walks, shifted against each other in time (\cref{def:dnp}). One
such walk alone is easy (\cref{lem:dep-visit}), and the pair with $m$ levels of nesting is depth-reducible: clock labels cut
it at a depth exponential in~$m$ (\cref{prop:dep-clock}). On a cyclic cover of the pair the cost of every label obeys exact
identities (Lemma~Zq, \cref{lem:chain,lem:transfer}); with them we prove, for the pair with one level of nesting, that
the labels that count exactly along one head pay a price per hop on the cover that does not depend on the depth
(\cref{thm:E,thm:E1}). Whether the pairs are depth-robust as $m$ grows remains open (\cref{conj:Aprime}).

\subsection{The problem}\label{sec:dep-problem}

We use the walk DAG of \cref{sec:dec-B}: for walks $x_1,\dots,x_k$ of length~$T$, its vertices are the times
$[0,T)$, its edges are the state edges $t\to t-1$ and the long hops $u\to\prev_h(u)$, a pause is identified with the
state edge joining the same two times, and an edge is determined by its endpoints. A \emph{cut} is a set of edges; it
has \emph{depth at most~$D$} if every path that uses no edge of the cut contains at most $D$ long hops.

\begin{problem}[$\TCP(k)$ at fixed $\eps$]\label{prob:tcp}
Let $k\ge1$. Is there, for every $\eps>0$, an integer $D=D(k,\eps)$ such that for every~$T$ and all walks
$x_1,\dots,x_k$ of length~$T$, coupled arbitrarily, the walk DAG has a cut with at most $\eps T$ edges and depth at
most~$D$?
\end{problem}

A family of $k$-tuples of walks is \emph{depth-robust} if for some $\eps>0$ no single~$D$ serves all its members.
Cuts are conveniently described by labels. For an integer $D\ge0$ and a map $\Psi\colon[0,T)\to\{0,\dots,D\}$, call a
long hop $u\to p$ \emph{violated} if $\Psi(u)\le\Psi(p)$, call a time $t\ge1$ a \emph{descent} if
$\Psi(t)<\Psi(t-1)$, and let $\cost(\Psi)$ be the number of violated long hops plus the number of descents.

\begin{lemma}[Cuts are labels]\label{lem:dep-label}
\notformalized
For all walks $x_1,\dots,x_k$ of length~$T$ and every $D\ge0$, the smallest size of a cut of depth at most~$D$ equals
the smallest value of $\cost(\Psi)$ over all maps $\Psi\colon[0,T)\to\{0,\dots,D\}$.
\end{lemma}

\begin{informalproof}
Given $\Psi$, let $F$ consist of the violated long hops and of the state edges $t\to t-1$ at descents~$t$; so
$\card F=\cost(\Psi)$. Along every edge $a\to b$ outside~$F$ we have $\Psi(b)\le\Psi(a)$, with strict inequality for
long hops. Along a path avoiding~$F$, $\Psi$ therefore never increases and drops at every long hop; as it takes values
in $\{0,\dots,D\}$, the path contains at most $D$ long hops. Conversely, let $F$ be a cut of depth at most~$D$, and
let $\Psi(t)$ be the largest number of long hops on a path that starts at~$t$ and avoids~$F$; so
$0\le\Psi(t)\le D$. For an edge $a\to b$ not in~$F$, prepending it to a path from~$b$ gives
$\Psi(a)\ge\Psi(b)+1$ if the edge is a long hop and $\Psi(a)\ge\Psi(b)$ otherwise. Hence every violated long hop and
every state edge at a descent of~$\Psi$ belongs to~$F$, and $\cost(\Psi)\le\card F$.
\end{informalproof}

\paragraph{Why the problem matters.} \Cref{prob:tcp} is the walk-level form of the structural hypothesis of \cref{thm:D}. If it
has a positive answer for~$k$, the proof of \cref{cor:B}, with the cut provided by \cref{prob:tcp} in place of the one
provided by \cref{thm:B}, shows that $\mathrm{depth}\text{-}\TCP(k,C,\beta,\eps,D)$ holds for all $C\ge1$,
$\beta\ge2$ and $0<\eps\le1$, with $D=(D(k,\eps/(4\beta C))+1)\lceil4\beta C/\eps\rceil$; \cref{thm:D} then reduces
the oblivious case
of~(T) to the crux alone, as \cref{cor:B} does for home-based heads. The problem is posed for all walks, and not only
for the head schedules of transposers, because a machine can run any schedule that its finite control generates on
extra work tapes without changing its output; such dummy traffic enters the dependency DAG. A family that
oblivious machines can generate in linear time, and that is depth-robust at every $\eps$ below the crux rates in
question, would therefore make the structural hypothesis fail for some machines, and the route of \cref{thm:D} would
need a version of it that sees only the part of the DAG able to carry the transposition. On the
construction side, \cref{thm:D} says that an oblivious linear-time transposer either beats the crux or has a head
schedule without genies of small cost and bounded depth; up to constants this means that its walk DAG has no cut of
size $\eps T$ and bounded depth, since a genie~$\Gamma$ yields a cut of at most $(k+1)\card\Gamma$ edges and the
conversions of \cref{sec:dec-corB} go the other way. The walks of this section show what the second alternative could
look like.

For $k=1$ the problem is easy. Call a long hop $(p,u)$ of a single walk~$x$ an upper hop if $x>x(p)$ on $(p,u)$ and a
lower hop if $x<x(p)$ there. Two upper hops never cross as intervals of time: if $p<p'<u<u'$, then $x(p')>x(p)$
because $p'\in(p,u)$, while $u\in(p',u')$ and $x(u)=x(p)<x(p')$, contradicting that $x>x(p')$ on $(p',u')$; the same
holds for lower hops. Every long hop is an upper or a lower hop, since $x\ne x(p)$ on $(p,u)$ and $x$ moves by at most~$1$ per
step. Drawing the upper hops on one side of the time axis and the lower hops on the other shows that
the walk DAG of one walk is planar; its maximum degree is at most~$4$. By the planar separator theorem of Lipton and
Tarjan~\cite{LiptonTarjan1979}, applied recursively as in~\cite{LiptonTarjan1980}, some set of $O(T/\sqrt M)$ vertices
leaves components of at most $M$ vertices. Cutting the at most four edges at each of these vertices cuts
$O(T/\sqrt M)$ edges, and a path that avoids the cut stays inside one component, so it has at most $M-1$ long hops;
this gives $D(1,\eps)=O(\eps^{-2})$. For $k\ge2$ the union of the walks' hop families is in general not planar, and only special
classes are solved, chief among them the home-based walks of \cref{thm:B}.

\paragraph{Relation to segregators.} Paul, Pippenger, Szemer\'edi and Trotter~\cite{PPST1983} showed that the
computation graphs of multitape machines have small segregators: deleting $O(n/\log^*n)$ vertices leaves every vertex
with at most $O(n/\log^*n)$ predecessors. This is the combinatorial core of their separation of deterministic from
nondeterministic linear time. \Cref{prob:tcp} is a depth version of the same problem with constant parameters.
Through the conversions of \cref{sec:dec-corB}, a positive answer gives genies of cost $\eps L$ and bounded
depth~$D'$, and by the cone lemma (\cref{lem:cone}) every step then has at most $(k+1)^{D'+1}$ predecessors reached
without passing through the genie: the deleted set is a fixed fraction of the steps, far more than $n/\log^*n$, but
the remaining cones have bounded size instead of size $n/\log^*n$. Santhanam stated that if, for every~$r$, the
multi-pushdown graphs $H_r(n)$, a class containing the computation graphs of multitape machines, had segregators
deleting $o(n/\log n)$ vertices and leaving every vertex at most $o(n)$ predecessors, then deterministic and
nondeterministic time $n\log n$ would differ~\cite[Thm.~2.7]{Santhanam2001}, and indicated that its proof is analogous
to that of his Theorem~2.5. Since the deleted
part in \cref{prob:tcp} is linear, a positive answer would not by itself give a new separation of complexity classes.

\subsection{Uniform recursive walks}\label{sec:dep-walks}

The instances of this section are built from one family of walks in which every level of a recursive nesting
receives the same total time. Since \cref{prob:tcp} concerns all schedules, these walks are relevant even though the
results of \cref{sec:motion} constrain the macroscopic motion of transposers.

\begin{definition}[Uniform recursive walks]\label{def:dep-walk}
Let $g\ge1$ and $m\ge0$ be integers, and let $s_0$ be an integer with $s_i:=s_0-ig\ge1$ for $0\le i\le m$. A
\emph{level-$m$ stroke} of direction $\sigma\in\{-1,1\}$ is the sequence of $2^{s_m}$ steps equal to~$\sigma$. For
$i<m$, a \emph{level-$i$ stroke} of direction~$\sigma$ is the concatenation, for $j=0,\dots,2^{g-1}-1$, of a
\emph{segment} of $2^{s_{i+1}+1}$ steps equal to~$\sigma$ followed by a \emph{V}: a level-$(i+1)$ stroke of
direction~$-\sigma$ (the \emph{back stroke}) followed by a level-$(i+1)$ stroke of direction~$\sigma$ (the
\emph{forward stroke}). The walk $U(s_0,m,g)$ is the walk $x\colon\Z\to\Z$ with $x(0)=0$, periodic with period
$P:=4(m+1)2^{s_0}$, whose steps $x(t+1)-x(t)$, $t\in[0,P)$, are those of four level-$0$ strokes of directions
$1,-1,-1,1$, in this order.
\end{definition}

By induction on $m-i$, a level-$i$ stroke has $(m-i+1)2^{s_i}$ steps and moves the head by $\sigma2^{s_i}$: its
segments contribute $2^{g-1}\cdot2^{s_{i+1}+1}=2^{s_i}$, and each V returns the head to where it started. So every
period has $4\cdot2^{s_0}=P/(m+1)$ steps belonging to segments of level-$i$ strokes for each $i<m$, and as many
belonging to level-$m$ strokes, and the walk never pauses. Each step belongs to exactly one segment or level-$m$
stroke; we say that the stroke emits it, and a time~$t$ belongs to the stroke that emits the step from $t$ to $t+1$.
The \emph{span} of a stroke or of a V is the set of times belonging to it, and its \emph{end time} is the first time
after the span. The terms \emph{during}, \emph{first time} and \emph{last time} refer to spans.

If a level-$i$ stroke of direction~$\sigma$, $i<m$, starts at cell~$a$, its segment~$j$ ends at the \emph{grid point}
$b_j=a+\sigma(j+1)2^{s_{i+1}+1}$, and the V that follows is \emph{attached} at~$b_j$. By
\cref{lem:dep-structure}(a) below, this V stays in its \emph{range}, the interval of cells between its \emph{base}
$b_j$ and its \emph{valley cell} $b_j-\sigma2^{s_{i+1}}$; the \emph{level} of the V is the level of its strokes, its
\emph{apex time} is the first time of its back stroke and its \emph{valley time} the first time of its forward
stroke. The \emph{open half-segments} of the stroke are the open intervals of cells between consecutive points of
$a+\sigma2^{s_{i+1}}\Z$ that lie between $a$ and $a+\sigma2^{s_i}$; the \emph{open range} of a level-$i$ stroke is the set of
cells strictly between $a$ and $a+\sigma2^{s_i}$, and for a level-$m$ stroke we use it in place of the open
half-segments; the open range of a V is its range without the base and the valley cell. A \emph{turn} is a time~$t$ with
$x(t+1)-x(t)\ne x(t)-x(t-1)$, and its \emph{level} is the level of the stroke to which $t$ belongs. Intervals of
cells are written without regard to orientation, and ``beyond~$b$ in direction~$\sigma$'' means $\sigma(c-b)>0$.

\begin{lemma}[Structure of $U$]\label{lem:dep-structure}
\notformalized
Let $x=U(s_0,m,g)$.
\begin{enumerate}[label=(\alph*)]
\item During a level-$i$ stroke~$S$ of direction~$\sigma$ starting at cell~$a$, including its end time, the head
stays between $a$ and $a+\sigma2^{s_i}$, and it is at $a+\sigma2^{s_i}$ at the end time. The first and the last step
of~$S$ equal~$\sigma$. If $i<m$, the part of~$S$ up to the end time of the V at~$b_j$ stays between $a$ and~$b_j$, and
the part after it stays between $b_j$ and $a+\sigma2^{s_i}$. During a V, including its end time, the head stays in
the range of the V.
\item The turns are exactly the apex and valley times of all V's, each of the level of its V, and, in every period,
the first times of the second and of the fourth level-$0$ stroke, which have level~$0$; at these two times the head
is at $2^{s_0}$ and at $-2^{s_0}$.
\item During a V of level~$i$, every cell strictly inside its range is visited exactly $2(m-i+1)$ times and each of
its two endpoints exactly $m-i+1$ times. In every period, every cell~$c$ with $\abs c<2^{s_0}$ is visited exactly
$2(m+1)$ times and each of $\pm2^{s_0}$ exactly $m+1$ times.
\item For every stroke~$S$ whose first time~$t_S$ is nonnegative, the number of times $t\in[0,t_S)$ with $x(t)=c$
is the same for all cells~$c$ of any one open half-segment of~$S$ (of the open range of~$S$ if $S$ has level~$m$).
\item Let $1\le j\le m$. If no turn of level less than~$j$ lies strictly inside a hop $(p,u)$, then
$\abs{x(t)-x(p)}\le2^{s_j}$ for all $t\in[p,u]$.
\end{enumerate}
\end{lemma}

\begin{informalproof}
(a) Induction on $m-i$. A level-$m$ stroke is monotone. For $i<m$, the segments move the head monotonically from~$a$
to $a+\sigma2^{s_i}$, reaching $b_j$ at the end of segment~$j$. By the induction hypothesis, the back stroke of the V
at~$b_j$ stays between $b_j$ and $b_j-\sigma2^{s_{i+1}}$ and ends at the valley cell, and the forward stroke stays in
the same range and ends at~$b_j$. This range lies between $b_j-\sigma2^{s_{i+1}+1}$, the start of segment~$j$, and
$b_j$, which gives all the claims about~$S$. The first step of~$S$ is a segment step, and its last step is the last
step of its last forward stroke, which equals~$\sigma$ by induction.

(b) Inside a segment and inside a level-$m$ stroke all steps are equal. Let $t$ be a time at which the steps from
$t-1$ to~$t$ and from $t$ to $t+1$ are emitted by different strokes. Let $Y$ be the largest segment or stroke whose
span ends at $t-1$, and $Z$ the largest segment or stroke whose span begins at~$t$. They are consecutive parts of one
stroke or of the sequence of level-$0$ strokes, so the pair $(Y,Z)$ is of one of four kinds, and by~(a) the last step
of~$Y$ and the first step of~$Z$ are their directions. (1)~$Y$ is a segment of a stroke of direction~$\sigma$ and $Z$
the back stroke of the V attached at its end: the steps are $\sigma$ and~$-\sigma$, and $t$ is the apex time, a turn
of the level of the back stroke. (2)~$Y$ is a back stroke and $Z$ the forward stroke of the same V: the steps are
$-\sigma$ and~$\sigma$, and $t$ is the valley time, a turn of the level of the forward stroke. (3)~$Y$ is a forward
stroke and $Z$ the next segment of the stroke to which its V is attached: both steps are~$\sigma$, and $t$ is not a
turn. (4)~$Y$ and~$Z$ are consecutive level-$0$ strokes: their directions $1,-1,-1,1$, followed by~$1$ in the next
period, give a turn exactly when $Z$ is the second or the fourth stroke of a period, and then the head is at
$2^{s_0}$ or $-2^{s_0}$ by~(a).

(c) By induction on $m-i$ we prove the claim about V's. A V of level~$m$ consists of two monotone runs of $2^{s_m}$
steps in opposite directions: the back stroke visits every cell of the range except the valley cell, the forward
stroke every cell except the base. Let $i<m$, and suppose by symmetry that the back stroke goes down, so that the
range is $[v,b]$ with $v=b-2^{s_i}$. Put $s=s_{i+1}$ and $K=2^{g-1}$, and let $v=\gamma_0<\gamma_1<\dots<\gamma_K=b$
be the points of $v+2^{s+1}\Z$ in $[v,b]$; since $2^{s_i}=K2^{s+1}$, these are also the points of $b-2^{s+1}\Z$ in
$[v,b]$. The segments of the back stroke visit $b,b-1,\dots,v+1$ once each, and its V's have level $i+1$, are attached
at $\gamma_{K-1},\dots,\gamma_0$ and have ranges $[\gamma_k,\gamma_k+2^s]$. The segments of the forward stroke visit
$v,v+1,\dots,b-1$ once each, and its V's are attached at $\gamma_1,\dots,\gamma_K$ and have ranges
$[\gamma_k-2^s,\gamma_k]$. By the induction hypothesis each of these V's visits the cells strictly inside its range
$2(m-i)$ times and its endpoints $m-i$ times, and the spans of the segments and of these V's partition the span of the
V. A cell strictly between $\gamma_k$ and $\gamma_k+2^s$, or strictly between $\gamma_k+2^s$ and $\gamma_{k+1}$, lies
strictly inside exactly one of the attached V's and is visited $1+1+2(m-i)$ times. The cell $\gamma_k+2^s$
($0\le k<K$) is an endpoint of two attached V's, the back stroke's at $\gamma_k$ and the forward stroke's at
$\gamma_{k+1}$, and so is every $\gamma_k$ with $0<k<K$; these cells are visited $1+1+2(m-i)$ times. The base~$b$ is
visited once by the back stroke's segments and $m-i$ times by the forward stroke's V at $\gamma_K=b$, and the valley
cell~$v$ once by the forward stroke's segments and $m-i$ times by the back stroke's V at $\gamma_0=v$. This proves
the claim for V's. In a period, the first two level-$0$ strokes have, over $[0,2^{s_0}]$, the same segments and
attached V's as the forward and back strokes of a V over this range (in the other order), so the count for V's applies
verbatim and shows that they visit every $c$ with $0<c<2^{s_0}$ exactly $2(m+1)$ times and $0$ and $2^{s_0}$ exactly $m+1$
times; the last two strokes do the same over $[-2^{s_0},0]$.

(d) For a stroke~$S$ with $t_S\ge0$ write $\mathrm{pre}_S(c)$ for the number of times $t\in[0,t_S)$ with $x(t)=c$.
By~(a) and~(c), during a level-$i$ stroke with $i<m$ the number of visits to~$c$ is constant on each open
half-segment: the segments visit each cell between the first cell (included) and the last cell (excluded) once, and
the V at~$b_j$ adds $2(m-i)$ visits on the open half-segment of segment~$j$ next to~$b_j$ and none elsewhere inside
open half-segments. During a level-$m$ stroke every cell of the open range is visited once. We now argue down the
tree of strokes; since time~$0$ starts a level-$0$ stroke, the parent of every stroke~$S$ with $t_S\ge0$ also starts at a
nonnegative time. Let $S$ be a level-$0$ stroke of period~$n$. If it is the first or the third stroke of the period,
$\mathrm{pre}_S(c)=2(m+1)n$ on its open range by~(c). If it is the second (fourth) stroke, $\mathrm{pre}_S$ is the sum
of $\mathrm{pre}_{S'}$ and of the visits during~$S'$, where $S'$ is the first (third) stroke of the period, and $S$
and~$S'$ have the same open half-segments because $2^{s_0}$ is a multiple of $2^{s_1}$ (for $m=0$, the same open range). Next let $S$ be the back
stroke of a V attached at~$b_j$ to a stroke~$S'$. The open range of the V is an open half-segment of~$S'$, and on it
$\mathrm{pre}_S=\mathrm{pre}_{S'}+1$: by~(a) the part of~$S'$ before segment~$j$ stays between the first cell of~$S'$
and the start of segment~$j$, and segment~$j$ visits each cell of the open range once. So $\mathrm{pre}_S$ is constant
on the open range of~$S$, which contains all its open half-segments. If $S$ is the forward stroke of that V and $B$ its
back stroke, $\mathrm{pre}_S$ is the sum of $\mathrm{pre}_B$ and of the visits during~$B$, and $B$ and~$S$ have the
same open half-segments, or the same open range at level~$m$, because their common displacement is a multiple of the
spacing of their half-segments. Hence $\mathrm{pre}_S$ is constant on the open half-segments, or the open range, of every
stroke.

(e) Let $c=x(p)$. As $u-p\ge2$ and the walk moves by $\pm1$, the first time $\tau^*\in(p,u)$ at which $\abs{x(t)-c}$
is largest is a turn, and it suffices to show $\abs{x(\tau^*)-c}\le2^{s_j}$. The level of~$\tau^*$ is at least
$j\ge1$, so by~(b) $\tau^*$ is the apex or valley time of a V of level at least~$j$, and therefore belongs to the span
of a V of level exactly~$j$, say~$\mathcal V$: a V of level greater than~$j$ lies inside a stroke of level at
least~$j$, and such a stroke lies inside a V of level~$j$. Let $\mathcal V$ be attached to the stroke~$S'$, of
level~$j-1$ and direction~$\sigma$, at its base~$b$, so that its range is $[v,b]$ with $v=b-\sigma2^{s_j}$. By~(a),
$x(\tau^*)\in[v,b]$, so it suffices to show $c\in[v,b]$. Suppose not. By~(a) the head is in $[v,b]$ during
$\mathcal V$ and at its end time~$t_E$, so $p$ precedes the apex time~$t_A$ of~$\mathcal V$ and $u$ follows~$t_E$.

\emph{Case 1: $c$ lies beyond~$b$ in direction~$\sigma$.} By~(a) the part of~$S'$ before~$\mathcal V$ stays between
the first cell of~$S'$ and~$b$, so $p$ precedes the first time of~$S'$. If $j\ge2$, $S'$ is a back or forward stroke
of a V of level $j-1$, whose first time is an apex or valley time, a turn of level $j-1$ in $(p,t_A)$. If $j=1$ and
$S'$ is the second or fourth level-$0$ stroke of its period, its first time is a turn of level~$0$ in $(p,t_A)$. If
$S'$ is the first (third) level-$0$ stroke, then $\sigma=1$ and $c>b>0$ ($\sigma=-1$ and $c<b<0$); the two preceding
level-$0$ strokes stay in $[-2^{s_0},0]$ (in $[0,2^{s_0}]$), so $p$ precedes both of them, and the first time of the
later one, the fourth (second) level-$0$ stroke, is a turn of level~$0$ in $(p,t_A)$.

\emph{Case 2: $c$ lies beyond~$v$ in direction~$-\sigma$.} Let $A_{j-1}=S'$, and for $1\le l\le j-1$ let $A_{l-1}$ be
the stroke to which the V containing $A_l$ is attached; $A_l$ has level~$l$. From~$t_E$ on, by~(a), the head stays
beyond~$b$ in direction~$\sigma$, or at~$b$, until the end time of~$A_{j-1}$. If $A_{j-1}$ is a back stroke, its end
time is the valley time of its V, a turn of level $j-1$. If $A_{j-1}$ is a forward stroke, it has the direction
$\sigma$ of $A_{j-2}$, its end time is the end time of a V attached to~$A_{j-2}$ at the last cell of~$A_{j-1}$, and
by~(a) the head stays beyond that cell, hence beyond~$b$, until the end time of~$A_{j-2}$. Repeating this, if some
$A_l$ with $1\le l\le j-1$ is a back stroke and $l^*$ is the largest such~$l$, the head stays beyond~$b$ (or at~$b$)
from~$t_E$ until the end time of~$A_{l^*}$, which is a valley time, a turn of level $l^*\le j-1$. Otherwise all
$A_l$ have direction~$\sigma$, and the head stays beyond~$b$ (or at~$b$) until the end time of the level-$0$
stroke~$A_0$. If $A_0$ is the first (third) stroke of its period, that end time is a turn of level~$0$. If it is the
second (fourth) stroke, the next level-$0$ stroke continues in direction~$\sigma$ and stays beyond~$b$, and its end
time is a turn of level~$0$. In every case the turn follows~$t_E$, and the head is not at~$c$ between $t_E$ and the
turn, so the turn precedes~$u$.

In both cases a turn of level less than~$j$ lies strictly inside $(p,u)$, contrary to the hypothesis.
\end{informalproof}

\subsection{Union labels and the nesting barrier}\label{sec:dep-nesting}

The labels of \cref{thm:B} are lexicographic: they carry counts of events at a few scales and compare them from the
top. A natural way to extend this to arbitrary walks is to let each head decide, at each cell and time, which bands of
scales (groups of consecutive scales; in \cref{sec:dec-B} a band was a single scale) it currently needs, and to carry, for every band needed by some head, a count of the events of that band (a
\emph{union label}). Two requirements make such a label work. A hop must find its own band among the tracked bands at
both of its ends, so that it is decided there; and a change of the bands tracked by one head breaks the comparison
along the hops of the other heads that span it, of which there can be about $2^\theta$ when $\theta$ is the bottom
scale of the changed band, so changes at high bands must be rare. We make the per-head part precise and sketch, without a
full proof, why it cannot be met with bounded load on deeply nested walks. We do not claim that such rules would
suffice for \cref{prob:tcp}; the point is that they appear not to exist.

\begin{definition}[Band rules]\label{def:dep-band}
Let $x$ be a walk of length~$T$ without pauses. For a turn~$\tau$ at which $x$ has a local maximum, a time~$t$ is
\emph{higher} than~$\tau$ if $x(t)>x(\tau)$, or $x(t)=x(\tau)$ and $t<\tau$; the \emph{prominence} $\mathrm{pr}(\tau)$
is the least $d\ge1$ such that on one side of~$\tau$ the walk reaches a time higher than~$\tau$ without going below
$x(\tau)-d$ on the way, and $x(\tau)-\min x$ if no time is higher than~$\tau$. For local minima exchange above and
below, an earlier time at the same height counting as lower. Fix integers $h_0,W\ge1$ and offsets $r\in[0,h_0)$,
$r'\in[0,W)$. A scale $\theta\ge0$ lies in \emph{band} $\lfloor(\lfloor(\theta+r)/h_0\rfloor+r')/W\rfloor$, and
$\theta(b)$ is the least scale of band~$b$. The band of a turn is the band of $\lfloor\log_2\mathrm{pr}(\tau)\rfloor$.
The \emph{owner} of a hop $(p,u)$ is the first time in $(p,u)$ at which $\abs{x(t)-x(p)}$ is largest, a turn, and the
\emph{own band} of the hop is the band of its owner. For a band $b\ge1$, the \emph{$b$-gaps} are the intervals
$[\tau,\tau')$ between consecutive turns of band at least~$b$ (and the interval before the first such turn). A
\emph{band rule} assigns to every cell~$c$ and time~$t$ a set $\mathcal M(c,t)$ of bands containing~$0$. It is
\emph{consistent} if for every $b\ge1$ the membership $b\in\mathcal M(c,t)$ depends only on~$c$ and on the $b$-gap
containing~$t$; its \emph{load} is $\max_t\card{\mathcal M(x(t),t)}$; $\mathrm{Viol}$ is the number of hops whose own
band is missing from $\mathcal M(x(p),p)$ or from $\mathcal M(x(u),u)$; and
\[
  D_{\mathrm{sum}}=\sum_{t}\;\sum_{b\ge1,\ b\in\mathcal M(x(t),t)\,\triangle\,\mathcal M(x(t+1),t+1)}2^{\theta(b)},
\]
while $D_{\max}$ counts, for each~$t$, only the largest band of the symmetric difference.
\end{definition}

A union label for $k$ heads tracks, at time~$t$, the union of the sets $\mathcal M_h(x_h(t),t)$ of the heads; with
load~$P_0$ per head it has at most $kP_0$ digits, so its depth is bounded in terms of $k$, $P_0$ and the moduli. Its
cost is bounded by the heads' violations, by their changes weighted as above (a change at band~$b$ is charged to the
hops of other heads that span it), and by wrap terms. This is the setting of the following two statements.

\begin{remark}[The nesting barrier for consistent band rules]\label{rem:dep-N}
Proof sketch only: we state the result precisely, but its proof exists only at the level of the following sketch;
exact linear programs support it numerically. Let
$g\ge\max(Wh_0,2)$ and $s_m\ge Wh_0+1$, let $x=U(s_0,m,g)$ on a whole number of periods, $T$ steps in all, and put
$M=T/(m+1)$. Then for every choice of the offsets, every consistent band rule of load at most~$P_0$ satisfies
\[
  \mathrm{Viol}+\gamma D_{\mathrm{sum}}\;\ge\;(m+2-P_0)M-2^{1-g}T-2^{1-s_m}M-O(m2^{s_0}),\qquad\gamma=2^{Wh_0-g};
\]
if the rule need not contain band~$0$, replace $m+2-P_0$ by $m+1-P_0$. \emph{Sketch.} Up to boundary effects the
prominences of the turns of~$U$ are $2^{s_0+1}$ and $2^{s_1},\dots,2^{s_m}$, one value per level of nesting (at a
double extremum the later turn carries the prominence of the deeper level), and for every offset their bands
$\beta_0>\beta_1>\dots>\beta_m$ are distinct, with $\theta(\beta_i)\ge s_i-Wh_0+1$. For a level-$i$ stroke~$S$ let
$\mathrm{gap}(S)$ be the $\beta_i$-gap containing its first segment step, and let $y(c,S)\in\{0,1\}$ be the membership
of~$\beta_i$ at $(c,\mathrm{gap}(S))$, well defined by consistency. (1)~For every cell~$c$ that the segments of~$S$
visit inside $\mathrm{gap}(S)$, the first visit of~$c$ in $\mathrm{gap}(S)$ ends a hop of own band~$\beta_i$, which
$y(c,S)=0$ violates; distinct pairs $(c,S)$ give distinct hops, except for a fraction at most $q_02^{-g}$ of each
level, $q_0=1/(2(1-2^{-g}))$, which comes from the last V's of strokes whose gap ends at their own arrival, and except
for the first cell of each back stroke. (2)~Every switch of $y(\cdot,S)$ between adjacent cells is a change of
$\beta_i$ at a segment step inside $\mathrm{gap}(S)$, of weight $2^{\theta(\beta_i)}\ge2^{1-Wh_0}\cdot2^{s_i}$, and
distinct pairs (stroke, switch) give distinct pairs (step, band). (3)~A cell~$c$ of a level-$m$ stroke with ancestors
$S_0\supset S_1\supset\dots\supset S_m$ is visited by the segments of~$S_m$ at a time that lies in $\mathrm{gap}(S_i)$
for every~$i$, except at exceptional levels, where $S_m$ lies in the last V of an ancestor whose gap ends early; so
the load gives $\sum_iy(c,S_i)\le P_0-1+\#\{\text{exceptional levels of }c\}$, and the expected number of exceptional
levels of a uniform such cell is at most $2q_02^{-g}m$. (4)~Averaging~(3) over a uniform level-$m$ stroke and a
uniform cell of it, which gives every block (a segment with its V) of every ancestor the same mass, and comparing
with~(1) and~(2) yields the bound; the exceptional terms add up to at most $2^{1-g}T$ because $g\ge2$.

Per head, a union label would need, heuristically, $\mathrm{Viol}\le\eps(N+T)$ and $D_{\max}\le\eps(N+T)$, where $N\le T$ is the number of
long hops. As $D_{\mathrm{sum}}\le D_{\max}/(1-2^{-Wh_0})$, the bound excludes this for all $P_0$, $W$ and~$h_0$ once
$\eps<0.37$ (take $g=Wh_0+4$, $m=10P_0$ and $s_m$ large), and once $\eps<0.21$ at $g=Wh_0+2$. Numerically, the exact
optimum of $\mathrm{Viol}+\gamma D_{\mathrm{sum}}$ over all consistent rules, computed by linear programming on
seventeen instances with $g\in\{1,2,3,4\}$ and $m\le5$, was $(m+2-P_0)M'-1$, where $M'$ is the mass per level in the
window used; the bound is loose only by its correction terms.
\end{remark}

\begin{remark}[The nesting barrier without consistency]\label{rem:dep-Nprime}
Proof sketch only: under the hypotheses of \cref{rem:dep-N}, let $t\mapsto\mathcal M(t)$ be any choice of sets of bands
containing~$0$, with $\card{\mathcal M(t)}\le P_0$ and no consistency requirement; it may depend on the whole walk, and
it may be randomised or fractional. Let $\mathrm{Viol}_u$ be the number of hops whose own band is missing from
$\mathcal M(u)$ at the later endpoint~$u$, and let $D_{\mathrm{sum}}$ count the changes of non-zero bands with weights
$2^{\theta(b)}$. Then $\mathrm{Viol}_u+\gamma D_{\mathrm{sum}}$ satisfies the bound of \cref{rem:dep-N}. Every band rule,
read along the walk as $\mathcal M(t)=\mathcal M(x(t),t)$, is such a choice, with the same load and changes and with
$\mathrm{Viol}_u\le\mathrm{Viol}$; so relaxing consistency does not help, and the thresholds $0.37$ and $0.21$ remain.
\emph{Sketch.} Steps~(2) and~(3) of the sketch of \cref{rem:dep-N} are the only uses of consistency, and both are replaced
by a trichotomy in time. For a level-$i$ stroke~$S$ and its block~$j$ (segment~$j$ with its V), let $J$ be the
block's time window, from just after the start of segment~$j$ to the start of segment~$j+1$ (the first and last
windows extend to the ends of $\mathrm{gap}(S)$); the windows of one level are disjoint. Exactly one of the following
holds: $\beta_i$ changes at a step inside~$J$, at cost $\gamma2^{\theta(\beta_i)}\ge2^{s_{i+1}+1}$, the number of
cells of segment~$j$; $\beta_i$ is absent throughout~$J$, so all the hops entering the cells of segment~$j$, whose
later endpoints lie in~$J$, are violated; or $\beta_i$ is present throughout~$J$, in particular at the deepest visits
inside the V of block~$j$, where it uses load. The load at the deepest visit of a cell then bounds the number of
levels~$i$ of the third kind by $P_0-1$, up to exceptional levels, and the averaging of step~(4) concludes. A
layer-cake decomposition of $1-\mathcal M$ gives a second proof, as a max-flow certificate. Numerically, when the own
band is required at both endpoints of a hop, the optimum over such choices of sets was lower than over consistent
rules only by a boundary term on the instances computed, and a bound with only boundary losses is conjectured; when it
is required at the later endpoint alone, the optimum was lower by a small multiple of~$M$.
\end{remark}

\paragraph{Why they matter.} \Cref{rem:dep-N,rem:dep-Nprime} indicate, with proof sketches only, that the union-label
route to \cref{prob:tcp} fails in its natural form: per-head band bookkeeping with load independent of the depth of
nesting, own-band decisions and $2^\theta$-weighted changes would fail already on a single walk, offline and for every
offset. The obstruction does not
concern $\TCP(1)$, which holds by planarity; what fails is a per-head format that would have handled all couplings at
once. A repair must decide some hops at a band other than their own, as the labels of \cref{thm:B} do, or weigh
changes differently, which is legitimate only against other heads that are not adversarial, or use a different
mechanism altogether.

\paragraph{One nested walk is nevertheless easy.} A single digit that ignores bands handles all levels of one walk at
once.

\begin{lemma}[One nested walk is easy]\label{lem:dep-visit}
\notformalized
Let $x=U(s_0,m,g)$ on $[0,T)$, where $T$ is a multiple of~$P$, let $q\ge1$, and let $\kappa(t)$ be the number of times
$s\in[0,t)$ with $x(s)=x(t)$. For $\omega\in\{0,\dots,q-1\}$ put $\Psi_\omega(t)=(\kappa(t)+\omega)\bmod q$. Then
\[
  \frac1q\sum_{\omega=0}^{q-1}\cost(\Psi_\omega)\;\le\;\frac{\#\{\text{long hops}\}}q+\frac{6T}{(m+1)2^{s_m}}.
\]
In particular, for every $\eps>0$ the walk DAG has a cut of at most $\eps T$ edges and depth at most
$\lceil2/\eps\rceil-1$ as soon as $(m+1)2^{s_m}\ge12/\eps$.
\end{lemma}

\begin{informalproof}
\emph{Hops.} For a hop $(p,u)$ we have $\kappa(u)=\kappa(p)+1$, so $\Psi_\omega(u)\le\Psi_\omega(p)$ if and only if
$\Psi_\omega(p)=q-1$, which holds for exactly one~$\omega$. Averaged over~$\omega$, the number of violated hops is
$\#\{\text{long hops}\}/q$ (all hops of~$U$ are long).

\emph{Descents.} A descent at~$t$ requires $\kappa(t)\ne\kappa(t-1)$; we count such~$t$. Let the step from $t-1$ to~$t$
be emitted by a segment of a level-$i$ stroke~$S$, $i<m$, running from $a_j$ to the grid point~$b_j$ in
$2^{s_{i+1}+1}$ steps, so that its middle cell is $a_j+\sigma2^{s_{i+1}}$. If neither $x(t-1)$ nor $x(t)$ is one of
$a_j$, $a_j+\sigma2^{s_{i+1}}$ and~$b_j$, both cells lie in one open half-segment of~$S$, both times belong to the
segment, and $\kappa(t-1)=\kappa(t)$: indeed, for a segment time~$s$ at such a cell~$c$, the visits to~$c$ before~$s$
are those before the first time of~$S$, because by \cref{lem:dep-structure}(a) the earlier parts of~$S$ stay between
the first cell of~$S$ and~$a_j$ and the segment reaches~$c$ only at~$s$; and the number of earlier visits is the same
for all cells of the open half-segment by \cref{lem:dep-structure}(d). So at most four steps emitted by the segment
change~$\kappa$: those leaving~$a_j$, entering and leaving the middle cell, and entering~$b_j$. In the same way a
level-$m$ stroke from~$a$ to $a+\sigma2^{s_m}$ emits at most two such steps, those leaving~$a$ and entering
$a+\sigma2^{s_m}$. A period contains $4\cdot2^{gi}$ strokes of level~$i$, so the number of changes of~$\kappa$ per
period is at most
\[
  4\sum_{i<m}4\cdot2^{gi}\cdot2^{g-1}+2\cdot4\cdot2^{gm}\;\le\;\frac{8\cdot2^g}{2^g-1}\,2^{gm}+8\cdot2^{gm}
  \;\le\;24\cdot2^{gm}=\frac{6P}{(m+1)2^{s_m}},
\]
using $2^{gm}=2^{s_0-s_m}$. Each change of~$\kappa$ is at most one descent of~$\Psi_\omega$, for every~$\omega$, which
gives the second term.

For the last claim take $q=\lceil2/\eps\rceil$: the first term is at most $\eps T/2$, the second at most $\eps T/2$,
and some~$\omega$ does at least as well as the average; \cref{lem:dep-label} turns $\Psi_\omega$ into a cut.
\end{informalproof}

\paragraph{Why it matters.} One deeply nested walk costs essentially nothing beyond the floor $1/q$ per hop, at every
depth of nesting, and larger~$m$ even helps. The label is a single anchored digit, since $\kappa$ rises by exactly~$1$
along every hop whatever happens in between, and it is not a band rule, so it escapes \cref{rem:dep-N,rem:dep-Nprime}. The
difficulty of the instance studied next lies entirely in the interaction of two heads.

\subsection{The deep-nesting pair}\label{sec:dep-dnp}

\Cref{thm:B} covers walks of linear leg area, in particular walks whose legs all begin or end near a home. In
$U(s_0,m,g)$ most legs do neither: the nesting happens far from home, at every level. Two copies of~$U$ shifted
against each other in time give the instance that we regard as the sharpest open case of \cref{prob:tcp} for $k=2$.

\begin{definition}[The deep-nesting pair]\label{def:dnp}
Let $s_0$, $m$, $g$ be as in \cref{def:dep-walk} and $P=4(m+1)2^{s_0}$. The \emph{golden shift} $\xi$ is the least odd
integer at least $\lfloor0.381966\,P\rfloor$. The \emph{deep-nesting pair} $\mathrm{DNP}(s_0,m,g)$ consists of the
walks $x_1=U(s_0,m,g)$ and $x_2(t)=x_1(t+\xi)$, $t\in\Z$. For $r\ge1$, its \emph{$r$-fold cover} is the pair on the
cyclic time axis $\Z_T$, $T=rP$. For $t\in\Z_T$ and a head~$h$, $\prev_h(t)$ is the residue modulo~$T$ of the previous
visit, on~$\Z$, of head~$h$ to $x_h(t)$; this defines bijections $\prev_1,\prev_2$ of $\Z_T$, and the \emph{hops} of
head~$h$ on the cover are the pairs $(\prev_h(t),t)$, $t\in\Z_T$, so there are $n_{\mathrm h}=2T$ hops in all. A
\emph{label} on the cover is a map $G\colon\Z_T\to\{0,\dots,q-1\}$ with $q\ge1$. A hop $(p,u)$ is \emph{violated} if
$G(u)\le G(p)$, a time~$t$ is a \emph{descent} if $G(t)<G(t-1)$, and $\cost(G)$ is the number of violated hops of both
heads plus the number $\mathrm{desc}(G)$ of descents. Hops are indexed by head and later endpoint, so a pair that is a
hop of both heads would count twice: $\cost(G)$ counts violations per head, unlike $\cost(\Psi)$ on the walk DAG;
by \cref{rem:dep-linear} no such pair exists when $s_m\ge2$.
\end{definition}

The constant $0.381966$ is $2-\phi$ to six decimals, $\phi$ the golden ratio, and we also call $\mathrm{DNP}(s_0,m,g)$
the \emph{golden pair}. It is meant to keep the two hierarchies
misaligned, but because of the truncation it fails to do so when $15625$ divides $m+1$: then $\xi\equiv1\pmod{2^{s_0-3}}$,
and the hierarchies are aligned at every scale up to $2^{s_0-3}$; for some other~$m$ they nearly align at several
scales (\cref{rem:dep-Em}). Only the arithmetic of \cref{thm:E1}, where $m=1$,
depends on its exact value.

The cover is a model of the walk
DAG of many periods in which the identities below are exact; \cref{rem:dep-linear} transfers the main bound back to
the walk DAG. A label on the cover plays the role of a label of depth $q-1$; by \cref{lem:dep-label}, labels lose
nothing on the walk DAG. A cell~$c$ is \emph{interior} if $\abs c<2^{s_0}$; by \cref{lem:dep-structure}(c), on the
$r$-fold cover each head visits every interior cell $r(2m+2)$ times and each of the two tips $\pm2^{s_0}$ $r(m+1)$
times.

For $t\in\Z_T$, identified with its representative in $[0,T)$, the \emph{visit index} $\kappa_h(t)$ is the number of
times $s\in[0,t)$ with $x_h(s)=x_h(t)$. The \emph{error} of a hop $(p,u)$ is
$e(p,u)=(G(u)-G(p)-1)\bmod q\in\{0,\dots,q-1\}$, and $G$ is \emph{exact} on the hop if the error is~$0$, that is,
$G(u)\equiv G(p)+1\pmod q$. The labels $(\kappa_1+\omega)\bmod q$ of \cref{lem:dep-visit} are exact on every head-$1$
hop at an interior cell when $q$ divides $r(2m+2)$. The \emph{level} of a hop $(p,u)$ of head~$h$ is the least level
of a turn of~$x_h$ strictly inside $(p,u)$ (on the cover, inside any lift of the hop to~$\Z$); every hop contains a
turn, its owner, so levels are defined.

\paragraph{Anchored counts.} The visit index is anchored: it rises by exactly one along every hop of its head. A
count restricted to a band of levels rises by exactly one along the hops of that band and is unchanged along the
others.

\begin{lemma}[Level-band anchored counts]\label{lem:dep-band}
\notformalized
For a head~$h$, a set~$B$ of levels and a time $t\ge0$, let $U_h^B(t)$ be the number of times $s\in[0,t)$ with
$x_h(s)=x_h(t)$ such that the hop of head~$h$ that begins at~$s$ has its level in~$B$. Then for every hop $(p,u)$ of
head~$h$ with $p\ge0$, $U_h^B(u)-U_h^B(p)$ is $1$ if the level of $(p,u)$ lies in~$B$ and $0$ otherwise.
\end{lemma}

\begin{informalproof}
The visits to $x_h(u)$ in $[0,u)$ are those in $[0,p)$ and $p$ itself, because $p$ and~$u$ are consecutive visits.
The visit~$p$ is counted in $U_h^B(u)$ exactly when the hop that begins at~$p$, which is $(p,u)$, has its level in~$B$.
\end{informalproof}

With $B$ all levels, $U_h^B$ is the visit index. Along the \emph{chain} of a cell, the visits of one head (here head~$1$) to it
in time order, a label exact on the head-$1$ hops whose
level lies in~$B$ rises by one modulo~$q$ wherever $U_1^B$ rises, so modulo~$q$ it differs from $U_1^B$ by a function
of the cell and of the number of earlier head-$1$ hops at that cell with level outside~$B$; this is the class used in
\cref{rem:dep-Em}.

\begin{remark}[Numerical: how often band counts change]\label{rem:dep-bandcount}
Numerically, on all instances checked, $U_1^{\{0,\dots,j\}}$ is constant during every V
of level greater than~$j$ at all times at which the head is not at the base of the V; at the base it may change, so the
plain statement ``constant during the V'' is false. The total number of its changes was of the order of
$T/((m+1)2^{s_{j+1}})+T(m-j)/((m+1)2^{s_j})$, so the gain $1/(m+1)$ of \cref{lem:dep-visit} is lost for $j<m$.
Nothing below uses this remark.
\end{remark}

\begin{lemma}[No monotone anchoring]\label{lem:dep-monotone}
\notformalized
Let $x$ be a time shift of $U(s_0,m,g)$, let $1\le i<m$, and let $\mathcal V$ be a V of level~$i$ with valley
cell~$v$ whose forward stroke has direction~$\sigma$. For $1\le\rho<2^{s_i}$ let $p_\rho$ be the last visit to
$v+\sigma\rho$ during the back stroke and $u_\rho$ its first visit during the forward stroke; $(p_\rho,u_\rho)$ is a
hop. Let $K$ be an integer function on the span of~$\mathcal V$ with $K(t)\ge K(t-1)$ whenever $t-1$ and~$t$ belong to
the span, and with $K(u')\ge K(p')+1$ for every hop $(p',u')$ of level $i+1$ with both endpoints in the span. Then
\[
  K(u_\rho)-K(p_\rho)\;\ge\;\rho/2^{s_{i+1}}-3/2\qquad(1\le\rho<2^{s_i});
\]
in particular $K(u_\rho)-K(p_\rho)\ge2^g-2$ for $\rho=2^{s_i}-1$.
\end{lemma}

\begin{informalproof}
After its last visit to $v+\sigma\rho$ the back stroke ends at~$v$, so by continuity it stays on the side of~$v$; the
forward stroke starts at~$v$ and first reaches $v+\sigma\rho$ at~$u_\rho$. So $p_\rho,u_\rho$ are consecutive visits.
Put $s=s_{i+1}$ and let $\gamma_k=v+\sigma k2^{s+1}$. By \cref{lem:dep-structure}(a), the back stroke has V's of
level $i+1$ attached at $\gamma_k$ with ranges $[\gamma_k,\gamma_k+\sigma2^s]$, $0\le k<2^{g-1}$, and the forward
stroke has V's attached at $\gamma_{k+1}$ with ranges $[\gamma_{k+1}-\sigma2^s,\gamma_{k+1}]$. If
$(2k+1)2^s<\rho$, the back stroke's V at~$\gamma_k$ lies strictly on the side of~$v$ of $v+\sigma\rho$; it is
attached after the back stroke has passed $v+\sigma\rho$, and nothing after it in the back stroke returns to
$v+\sigma\rho$, so it lies inside $(p_\rho,u_\rho)$. There are at least $(\rho/2^s-1)/2$ such~$k$. If
$(k+1)2^{s+1}<\rho$, the forward stroke reaches $\gamma_{k+1}$ and completes its V there before it reaches
$v+\sigma\rho$, so this V lies inside $(p_\rho,u_\rho)$ as well; there are at least $\rho/2^{s+1}-1$ such~$k$. In all,
at least $\rho/2^s-3/2$ V's of level $i+1$ lie inside $(p_\rho,u_\rho)$, and their spans are disjoint.

In each of them, with valley cell $v'$ and forward direction $\sigma'$, take the hop between the last visit of the back
stroke and the first visit of the forward stroke to $v'+\sigma'$. As $2^s\ge2$, this cell is not the base, so the hop
lies after the apex time; by \cref{lem:dep-structure}(b) the only turns strictly inside it are the valley time of this
V and turns of V's nested in it, so the hop has level $i+1$. Listing the chosen hops in time order and joining them by
the stretches of time between them, $K$ gains at least~$1$ along each chosen hop and does not decrease along the
stretches. Hence $K(u_\rho)-K(p_\rho)$ is at least the number of chosen hops. For $\rho=2^{s_i}-1$ the bound is
$2^g-2^{-s}-3/2\ge2^g-2$.
\end{informalproof}

\paragraph{Clock labels.} Counting turns instead of visits gives digits that never decrease. One digit per scale
suffices for both heads, because the two walks of a deep-nesting pair have the same scales.

\begin{proposition}[Clock labels]\label{prop:dep-clock}
\notformalized
Let $x_1$, $x_2$ be time shifts of $U(s_0,m,g)$, for instance a deep-nesting pair, on $[0,T)$, where $T$ is a multiple
of~$P$. For $0\le i\le m$ let $N^{(i)}(t)$ be the number of turns of level~$i$ of either head at times in $[0,t]$.
For an
integer $B\ge2$ and $\omega\in\{0,\dots,B-1\}^{m+1}$ let
\[
  \Psi_\omega(t)=\sum_{i=0}^{m}B^{m-i}\bigl((N^{(i)}(t)+\omega_i)\bmod B\bigr)\in\{0,\dots,B^{m+1}-1\}.
\]
Then the average of $\cost(\Psi_\omega)$ over all~$\omega$ is at most $12T/B$. Consequently, for every $\eps>0$ the walk
DAG of the pair has a cut of at most $\eps T$ edges and depth at most $\lceil12/\eps\rceil^{m+1}-1$.
\end{proposition}

\begin{informalproof}
The vector $(N^{(0)}(t),\dots,N^{(m)}(t))$ is coordinatewise non-decreasing in~$t$, and $\Psi_\omega(t)$ is the integer
whose base-$B$ digits, most significant first, are $(N^{(i)}(t)+\omega_i)\bmod B$.

\emph{Hops.} Let $(p,u)$ be a hop of head $h'\in\{1,2\}$ inside $[0,T)$. Its owner is a turn strictly inside $(p,u)$,
so some $N^{(i)}$ increases from $p$ to~$u$; let $j$ be the least such~$i$ and $\Delta=N^{(j)}(u)-N^{(j)}(p)\ge1$. The
digits of index less than~$j$ agree at $p$ and~$u$. If $(N^{(j)}(p)+\omega_j)\bmod B<B-\Delta$, digit~$j$ increases and
$\Psi_\omega(u)>\Psi_\omega(p)$ whatever the later digits are; for uniform $\omega_j$ the residue
$(N^{(j)}(p)+\omega_j)\bmod B$ is uniform, so the hop is violated with probability at most $\Delta/B$. Now $\Delta$ is
the number of turns of level~$j$ of either head at times in $(p,u]$, and no turn of level less than~$j$ of either head
lies in $(p,u]$. Fix a turn~$\tau$ of level~$j$ and consider the hops $(p,u)$ of head~$h'$ with $\tau\in(p,u]$ and
without turns of head~$h'$ of level less than~$j$ strictly inside. By \cref{lem:dep-structure}(e) if $j\ge1$, and
because $\abs{x_{h'}}\le2^{s_0}$ if $j=0$, the cell of every such hop lies within distance $2^{s_j}$ (within
$2^{s_0+1}$ if $j=0$) of $x_{h'}(\tau)$; and at each cell at most one hop of head~$h'$ has $\tau\in(p,u]$, because the
hops of one head at one cell are disjoint half-open intervals of time. So the turn~$\tau$ is counted, over both
heads, at most $2(2^{s_j+1}+1)$ times if $j\ge1$ and at most $2(2^{s_0+1}+1)$ times if $j=0$. By
\cref{lem:dep-structure}(b), each head has, per period, two turns of level~$0$ and $4\cdot2^{gj}$ turns of level~$j$
for $1\le j\le m$ (a period contains $2\cdot2^{gj}$ V's of level~$j$). Summing over the turns of both heads in one
period gives at most
\[
  4\cdot2(2^{s_0+1}+1)+\sum_{j=1}^{m}8\cdot2^{gj}\cdot2(2^{s_j+1}+1)
  =32(m+1)2^{s_0}-16\cdot2^{s_0}+8+16\sum_{j=1}^m2^{gj}\;\le\;8P+8\;\le\;9P,
\]
because $16\sum_{j\le m}2^{gj}\le32\cdot2^{gm}\le16\cdot2^{s_0}$ (as $s_m\ge1$) and $P\ge8$. So the expected number of
violated hops is at most $9T/B$.

\emph{Descents.} If no turn occurs at~$t$, then $\Psi_\omega(t)=\Psi_\omega(t-1)$. Otherwise the coordinates increase
by one for each head that turns at~$t$, so by at most two in all; the first changed digit increases by $1$ or~$2$, and
$\Psi_\omega$ descends only if this digit wraps, which has probability at most (number of turns at~$t$)$/B$. The two
heads have at most
$4+8\sum_{j\le m}2^{gj}\le4+8\cdot2^{s_0}\le3P$ turns per period, so the expected number of descents is at most
$3T/B$.

For the last claim we may assume $\eps<12$, since cutting all of the at most $3T$ edges gives depth~$0$; take
$B=\lceil12/\eps\rceil\ge2$, choose~$\omega$ with cost at most the average, and apply
\cref{lem:dep-label}.
\end{informalproof}

\paragraph{Why they matter.} By \cref{prop:dep-clock} every deep-nesting pair of fixed nesting depth~$m$ is depth-reducible,
with $D(\eps)$ of the form $(c/\eps)^{m+1}$; so depth-robustness can only appear as $m$ grows. The two lemmas show why
the obvious ways of beating the clock fail. For $g\ge2$, a digit that decides one head's hops at several levels exactly,
with slack below $2^g-2$, must descend inside every V of that head (\cref{lem:dep-monotone}; for $g=1$ the lemma is
empty), and the other head's long
hops that span such a descent see it. A digit that decides the other head's hops must instead be non-decreasing across
them. Anchored digits (\cref{lem:dep-visit,lem:dep-band}) and clock digits are the two ends of this spectrum, and
every explicit label we know is built from them.

\paragraph{Packing certificates.} A proof that the pairs are depth-robust as $m$ grows must show that every cut of depth~$D$ is
large. The standard way to certify such lower bounds, linear programming duality or packings of local obstructions,
cannot do it.

\begin{proposition}[No packing certificate]\label{prop:packing}
\notformalized
Let $x_1,\dots,x_k$ be walks of length~$T$, let $H$ be the set of long hops of their walk DAG, and let $D\ge0$. Call a
set~$F$ of edges \emph{$D$-forcing} if no map $\Psi\colon[0,T)\to\{0,\dots,D\}$ satisfies $\Psi(a)\ge\Psi(b)+1$ for
every long hop $a\to b$ in~$F$ and $\Psi(a)\ge\Psi(b)$ for every state edge $a\to b$ in~$F$.
\begin{enumerate}[label=(\alph*)]
\item Every $D$-forcing set contains a path with at least $D+1$ long hops.
\item Let $F_1,\dots,F_n$ be $D$-forcing and $y_1,\dots,y_n\ge0$ with $\sum_{j:e\in F_j}y_j\le1$ for every edge~$e$.
Then every cut of depth at most~$D$ has at least $\sum_jy_j$ edges, and $\sum_jy_j\le\card H/(D+1)\le kT/(D+1)$.
\item The linear program that minimises $\sum_ez_e$ over $z\ge0$ subject to $\sum_{e\in\pi}z_e\ge1$ for every path~$\pi$
with at least $D+1$ long hops has value at most $\card H/(D+1)$.
\end{enumerate}
Hence a lower bound $\eps T$ on the size of all cuts of depth at most~$D$ can be certified in the form~(b) or~(c) only
if $D<k/\eps$; such certificates never show depth-robustness.
\end{proposition}

\begin{informalproof}
(a) If $F$ contains no path with $D+1$ long hops, let $\Psi(a)$ be the largest number of long hops on a path in~$F$
that starts at~$a$. Then $0\le\Psi\le D$, and for every edge $a\to b$ of~$F$, $\Psi(a)\ge\Psi(b)+1$ if it is a long
hop and $\Psi(a)\ge\Psi(b)$ otherwise; so $F$ is not forcing.

(b) Let $F$ be a cut of depth at most~$D$, and let $\Psi$ be the map of the proof of \cref{lem:dep-label}; it satisfies
the constraints of all edges outside~$F$. Since each $F_j$ is forcing, it contains an edge of~$F$, and
$\card F\ge\sum_{e\in F}\sum_{j:e\in F_j}y_j\ge\sum_jy_j$. By~(a), $\card{F_j\cap H}\ge D+1$, so
$(D+1)\sum_jy_j\le\sum_jy_j\card{F_j\cap H}=\sum_{e\in H}\sum_{j:e\in F_j}y_j\le\card H$. Finally $\card H\le kT$,
because each time is the later endpoint of at most one hop of each head.

(c) Put $z_e=1/(D+1)$ on long hops and $z_e=0$ on state edges.

For the last sentence, $\card H/(D+1)\ge\eps T$ forces $D+1\le k/\eps$.
\end{informalproof}

On a cyclic time axis, such as the covers of \cref{def:dnp}, every cyclic hop chain is forcing: a chain of $n$ hops
always contains at least $\lceil n/q\rceil$ violated hops, more than the average $n/q$ unless $q$ divides~$n$. This is
why we use covers on which $q$ divides the visit numbers of the cells that matter; the effect does not arise on the
walk DAG, which is acyclic.

\paragraph{Why it matters.} \Cref{prop:packing} rules out the natural strategy for proving depth-robustness by duality:
packings of forcing structures, of any shape, certify at most what the trivial assignment of weight $1/(D+1)$ to
every long hop certifies. \Cref{lem:dep-euler}(c) below records the same barrier for cycle bounds in the cost
identity. A robustness proof cannot be a fractional packing; it has to exploit integrality, and we expect it to be
multiplicative: it would show that a label with $q$ values cannot serve many levels of both heads at once.

\subsection{Exact identities for labels}\label{sec:dep-identities}

On the cover every label obeys an exact cost identity, which separates the unavoidable floor $1/q$ per hop from the
errors of the hops.

\begin{lemma}[Lemma Zq: the cost identity]\label{lem:Zq}
\formal{DepthLabels.lean}{lemma_Zq, card_violated_eq, card_descents_eq}
Let $\tau$ and $H$ be finite sets, $q\ge1$ an integer, $G\colon\tau\to\{0,\dots,q-1\}$, and $t\mapsto t^-$ a bijection
of~$\tau$. Let $p,u\colon H\to\tau$ be maps and $\mathrm{succ}$ a bijection of~$H$ with
$p(\mathrm{succ}(h))=u(h)$ for all $h\in H$. Put $e_h=(G(u(h))-G(p(h))-1)\bmod q$ and $\chi_t=(G(t)-G(t^-))\bmod q$,
residues in $\{0,\dots,q-1\}$. Then
\[
  q\cdot\#\{h:G(u(h))\le G(p(h))\}=\card H+\sum_he_h,\qquad q\cdot\#\{t:G(t)<G(t^-)\}=\sum_t\chi_t,
\]
and hence $q\,(\#\{h:G(u(h))\le G(p(h))\}+\#\{t:G(t)<G(t^-)\})=\card H+\sum_he_h+\sum_t\chi_t$.
\end{lemma}

\begin{outline}
For $-q\le y<q$ one has $y\bmod q=y+q\cdot[y<0]$ (\lean{emod_of_bounds}). With $y=G(u(h))-G(p(h))-1$, which lies in
$[-q,q-2]$, this gives $e_h=G(u(h))-G(p(h))-1+q\cdot[G(u(h))\le G(p(h))]$. Summing over~$H$ and using
$\sum_hG(u(h))=\sum_hG(p(h))$, obtained by reindexing along the bijection $\mathrm{succ}$, gives the first identity. The
second is the same computation with $y=G(t)-G(t^-)\in(-q,q)$ and $\sum_tG(t^-)=\sum_tG(t)$; the third is their sum.
\end{outline}

On the cover of \cref{def:dnp} take $\tau=\Z_T$, $t^-=t-1$, let $H=\{1,2\}\times\Z_T$ index the hops of the two heads, with
$p,u$ their endpoints, and let $\mathrm{succ}$ map a hop to the next hop of the same head at the same cell; the hypothesis only
says that the hops of each head at each cell form a cycle. Since $\sum_t\chi_t=q\,\mathrm{desc}(G)$, the lemma reads
\begin{equation}\label{eq:dep-Zq}
  \cost(G)\;=\;\mathrm{desc}(G)+\frac1q\sum_{\text{hops }h}(1+e_h).
\end{equation}
Every hop costs at least $1/q$, and beyond that exactly its error divided by~$q$; every descent costs~$1$. For $q=2$,
\eqref{eq:dep-Zq} says that twice the cost is $n_{\mathrm h}$ plus the number of hops with $G(u)=G(p)$ plus the number
of times with $G(t)\ne G(t-1)$, so depth~$1$ is a problem of frustration in a signed graph. The identity also shows
that $\cost(G)$ does not change when $G$ is replaced by $(G+c)\bmod q$.

\begin{corollary}[Per-class split]\label{cor:dep-split}
\notformalized
On the cover of a deep-nesting pair, for every label~$G$,
\[
  \cost(G)=\mathrm{desc}(G)+\sum_{h=1}^{2}\sum_{i=0}^{m}\frac1q\sum_{\substack{\text{hops of head }h\\\text{of level
}i}}
  (1+e).
\]
\end{corollary}

\begin{informalproof}
Group the hops in \eqref{eq:dep-Zq} by head and level; each hop has exactly one of each.
\end{informalproof}

So the classes (head, level) contribute additively, and how well $G$ serves a class is measured by the error mass of
the class: a class on which $G$ is exact costs its size over~$q$.

\paragraph{Lifts.} To state the next lemmas without boundary terms we lift labels, walks and hops to~$\Z$: $G$ and the
walks are $T$-periodic, the hops of head~$h$ are the pairs $(\prev_h(u),u)$, $u\in\Z$, with $\prev_h$ the previous
visit on~$\Z$, and each hop of the cover is the residue of such a pair. For a head~$h$, a cell~$c$ and a visit
$s\in\Z$ of head~$h$ to~$c$, let $\hat\kappa_h(s)$ be the number of visits of head~$h$ to~$c$ at times in $[0,s)$ if
$s\ge0$, and minus the number at times in $[s,0)$ if $s<0$. Thus $\hat\kappa_h(s')-\hat\kappa_h(s)=1$ for consecutive
visits $s<s'$, $\hat\kappa_h(s+T)=\hat\kappa_h(s)+n_h(c)$, where $n_h(c)$ is the number of visits of head~$h$ to~$c$ per
cover, and $\hat\kappa_h=\kappa_h$ on $[0,T)$. Let $t_0(c)$ be the first visit of head~$h$ to~$c$ in $[0,T)$, and let
$\hat E_h(s)$ be the total error of the hops of head~$h$ at~$c$ whose later endpoint lies in $(t_0(c),s]$ if
$s\ge t_0(c)$, and minus the total error of those with later endpoint in $(s,t_0(c)]$ if $s<t_0(c)$. So
$\hat E_h(s')-\hat E_h(s)\ge0$ is the error mass of the hops of~$c$ ending in $(s,s']$ when $s<s'$.

\begin{lemma}[Chain decomposition]\label{lem:chain}
\partformal{DepthLabels.lean}{congr_of_exact}{the general form~(a), with error terms, and its reading~(c) on the cover.}
\begin{enumerate}[label=(\alph*)]
\item For every head~$h$, every cell~$c$ and every visit $s\in\Z$ of head~$h$ to~$c$,
\[
  G(s)\equiv G(t_0(c))+\hat\kappa_h(s)+\hat E_h(s)\pmod q,
\]
and the total error of the $n_h(c)$ hops of head~$h$ at~$c$ on the cover is congruent to $-n_h(c)$ modulo~$q$.
\item (Exact case, abstract form.) Let $\prev\colon\tau\to\tau$ and $\mathrm{cell}\colon\tau\to C$ be maps on a set
$\tau$ with $\mathrm{cell}(\prev(t))=\mathrm{cell}(t)$, such that whenever $\mathrm{cell}(s)=\mathrm{cell}(t)$ some
iterate of $\prev$ maps $t$ to~$s$. Let $S\subseteq C$, and let $G,\kappa\colon\tau\to\Z$ satisfy
$G(t)-G(\prev(t))\equiv1$ and $\kappa(t)-\kappa(\prev(t))\equiv1\pmod q$ whenever $\mathrm{cell}(t)\in S$. Then
$G(s)-\kappa(s)\equiv G(t)-\kappa(t)\pmod q$ whenever $\mathrm{cell}(s)=\mathrm{cell}(t)\in S$.
\item If $G$ is exact on all hops of head~$h$ at a cell~$c$, then $q$ divides $n_h(c)$ and
$G\equiv\kappa_h+G(t_0(c))\pmod q$ at the visits of~$c$ in $[0,T)$. Conversely, if $q$ divides $n_h(c)$ and
$G-\kappa_h$ is constant modulo~$q$ on the visits of~$c$ in $[0,T)$, then $G$ is exact on all hops of head~$h$ at~$c$.
\end{enumerate}
\end{lemma}

\begin{outline}[Outline of the formal proof of~(b)]
Induction on~$n$ shows $G(\prev^n(t))-\kappa(\prev^n(t))\equiv G(t)-\kappa(t)$: the cell stays in~$S$ because
$\mathrm{cell}\circ\prev=\mathrm{cell}$, and at each step the two congruences cancel. The orbit hypothesis supplies~$n$
with $\prev^n(t)=s$.
\end{outline}

\begin{proof}[Proof of~(a) and~(c) (not formalized)]
(a) For consecutive visits $s<s'$ of head~$h$ to~$c$, $G(s')-G(s)\equiv1+e(s,s')$ by the definition of the error,
while $\hat\kappa_h$ rises by~$1$ and $\hat E_h$ by $e(s,s')$. So $G-\hat\kappa_h-\hat E_h$ is constant modulo~$q$ on
the visits of~$c$, and at $t_0(c)$ it equals $G(t_0(c))$. At $s=t_0(c)+T$ we have $G(s)=G(t_0(c))$,
$\hat\kappa_h(s)=n_h(c)$, and $\hat E_h(s)$ is the total error of the $n_h(c)$ hops of~$c$ on the cover; this gives the
second claim.

(c) If all these errors vanish, (a) gives $0\equiv-n_h(c)$ and $G(s)\equiv G(t_0(c))+\kappa_h(s)$ for
$s\in[0,T)$. Conversely, let $t_0<\dots<t_{n-1}$ be the visits of~$c$ in $[0,T)$. For $1\le j<n$ the hop
$(t_{j-1},t_j)$ satisfies $G(t_j)-G(t_{j-1})\equiv\kappa_h(t_j)-\kappa_h(t_{j-1})=1$, and the remaining hop of~$c$ on
the cover, from $t_{n-1}$ to~$t_0$, satisfies $G(t_0)-G(t_{n-1})\equiv0-(n-1)\equiv1$ because $q$ divides~$n$.
\end{proof}

The visit numbers of the tips are $r(m+1)$ and those of interior cells $r(2m+2)$. So if $q$ divides $r(2m+2)$ but not
$r(m+1)$, for instance $r=q/\gcd(q,4)$ with $4\mid q$ at $m=1$, no label is exact on all head-$1$ hops; the natural
class is that of labels exact on all head-$1$ hops \emph{at interior cells}, which by~(c) are the labels
$G\equiv\kappa_1+\Phi(x_1)\pmod q$ at interior cells, for arbitrary functions~$\Phi$ of the cell. It contains the
visit-index labels of \cref{lem:dep-visit}.

\begin{lemma}[Transfer identity]\label{lem:transfer}
\partformal{DepthLabels.lean}{cycleDefect, cycle_errors_congr}{the general form~(a), with head-$1$ error mass, and the
absorption reading~(c).}
\begin{enumerate}[label=(\alph*)]
\item Let $h_j=(p_j,u_j)$, $j\in\Z_\ell$, be hops of head~$2$ on~$\Z$ with $x_1(u_j)=x_1(p_{j+1})$ for every~$j$. Put
$\Delta\hat\kappa_1(h_j)=\hat\kappa_1(u_j)-\hat\kappa_1(p_j)$ and define the \emph{defect}
$\delta=\sum_j(\Delta\hat\kappa_1(h_j)-1)$. Then
\[
  \sum_je(h_j)\;\equiv\;\delta-\sum_j\bigl(\hat E_1(p_{j+1})-\hat E_1(u_j)\bigr)\pmod q.
\]
\item (Exact case, abstract form.) Let $\prev_2\colon\tau\to\tau$, $\mathrm{cell}\colon\tau\to C$, $S\subseteq C$ and
$G,\kappa\colon\tau\to\Z$ be such that $G-\kappa$ is constant modulo~$q$ on the times of each cell of~$S$. Let
$b\colon\Z_\ell\to\tau$ satisfy $\mathrm{cell}(b(j))=\mathrm{cell}(\prev_2(b(j+1)))$ and $\mathrm{cell}(b(j))\in S$ for
all~$j$. Then
\[
  \sum_j\bigl(G(b(j))-G(\prev_2(b(j)))-1\bigr)\;\equiv\;\sum_j\bigl(\kappa(b(j))-\kappa(\prev_2(b(j)))-1\bigr)
  \pmod q;
\]
the right side is the defect \lean{cycleDefect} of $\kappa$, $\prev_2$ and~$b$.
\item (Absorption.) In~(a), call the head-$1$ stretch from $u_j$ to $p_{j+1}$ \emph{forward} if $u_j<p_{j+1}$ and
\emph{backward} otherwise, and let $A_j=\abs{\hat E_1(p_{j+1})-\hat E_1(u_j)}$ be its error mass. If $-q<\delta<0$ and
$\sum_je(h_j)<q-\abs\delta$, then either the backward stretches carry error mass at least
$\abs\delta+\sum_je(h_j)$, or the forward stretches carry error mass at least
$q-\abs\delta-\sum_je(h_j)+\sum_{\text{backward}}A_j$.
\end{enumerate}
\end{lemma}

\begin{outline}[Outline of the formal proof of~(b)]
After reindexing one of the sums by the rotation $j\mapsto j+1$ of $\Z_\ell$, the difference of the
two sides becomes $\sum_j[(G-\kappa)(b(j))-(G-\kappa)(\prev_2(b(j+1)))]$. Each term is divisible by~$q$, because
$b(j)$ and $\prev_2(b(j+1))$ are times of the same cell of~$S$.
\end{outline}

\begin{proof}[Proof of~(a) and~(c) (not formalized)]
(a) By the definition of the error, $G(u_j)-G(p_j)\equiv1+e(h_j)$. By \cref{lem:chain}(a) for head~$1$ at the cell
$x_1(u_j)=x_1(p_{j+1})$, whose base term cancels,
$G(p_{j+1})-G(u_j)\equiv\hat\kappa_1(p_{j+1})-\hat\kappa_1(u_j)+\hat E_1(p_{j+1})-\hat E_1(u_j)$. Summing both
congruences over $j\in\Z_\ell$, the values of~$G$ cancel, and $\sum_j(\hat\kappa_1(p_{j+1})-\hat\kappa_1(u_j))
=-\sum_j\Delta\hat\kappa_1(h_j)$. Rearranging gives~(a).

(c) Write $\sigma_e=\sum_je(h_j)$, $\delta'=-\delta\in(0,q)$, and let $F_A$ and $B_A$ be the error masses of the forward
and of the backward stretches. The term $\hat E_1(p_{j+1})-\hat E_1(u_j)$ is $A_j$ for a forward and $-A_j$ for a
backward stretch, so~(a) says $B_A-F_A\equiv\sigma_e+\delta'\pmod q$, where $0<\sigma_e+\delta'<q$. Write
$B_A-F_A=\sigma_e+\delta'+qz$ with $z\in\Z$. If $z\ge0$, then $B_A\ge\sigma_e+\delta'$. If $z\le-1$, then
$F_A=B_A-\sigma_e-\delta'-qz\ge q-\delta'-\sigma_e+B_A$.
\end{proof}

For labels exact on head~$1$ at interior cells, all errors of head-$1$ hops at interior cells vanish, so by~(a) or~(b)
every closed sequence of head-$2$ hops through interior head-$1$ cells has $\sum_je(h_j)\equiv\delta$: a cell
correction~$\Phi$ telescopes and absorbs nothing. In general, by~(c), a defect $-q<\delta<0$ can be absorbed cheaply,
at a cost of order $\abs\delta/q$, only by errors of head-$1$ hops on backward stretches, which lie inside the overlaps of
consecutive head-$2$ hops; errors on forward stretches enter with the wrong sign and help only after a wrap, at a cost of nearly a whole
unit.

\begin{remark}[Numerical: who absorbs in optimised labels]\label{rem:dep-absorb}
For the deep-nesting pair with $m=3$, $g=1$ and $s_0=6$, we examined the best labels found by optimisation at
$D=7,15,31,63$; they cost about four to five times less than clock labels with the same number of values. For the
pairs of head-$2$ hops $a\to b$, $b\to a$ in head-$1$ cell space whose defect a label exact on head~$1$ could not
absorb, head-$1$ errors on backward stretches did the absorbing in $55\%$, $83\%$, $88\%$ and $96\%$ of the cases,
and head~$2$ paid (error sum at least $q/2$) in $28\%$, $9\%$, $9\%$ and $4\%$. The transfer identity held on all
$6925$ pairs examined.
\end{remark}

\begin{lemma}[Euler: balanced degrees, chain defects, and the barrier in $\Z_q$ form]\label{lem:dep-euler}
\notformalized
On the $r$-fold cover of a deep-nesting pair:
\begin{enumerate}[label=(\alph*)]
\item The multigraph $M_2$ on the cells of head~$1$ with one edge $x_1(p)\to x_1(u)$ for every hop $(p,u)$ of head~$2$
is balanced: the in-degree and the out-degree of a cell~$c$ both equal the number of visits of head~$1$ to~$c$.
\item Let head~$2$ visit the cell~$c'$ at the times $t_0<\dots<t_{n-1}$ of $[0,T)$. Its hops at~$c'$, lifted to
$(t_{j-1},t_j)$ for $1\le j<n$ and $(t_{n-1},t_0+T)$, satisfy the hypothesis of \cref{lem:transfer}(a), and their
defect is $n_1(x_1(t_0))-n_2(c')$. It is $0$ if $c'$ and $x_1(t_0)$ are both interior cells or both tips.
\item In the setting of \cref{lem:transfer}(a), with $A_j$ as in \cref{lem:transfer}(c),
$\sum_je(h_j)+\sum_jA_j\ge\lVert \delta\rVert_q$, where $\lVert
\cdot\rVert_q$
is the distance to $q\Z$, and $\abs\delta$ is at most the number of hops on the closed walk formed by the~$h_j$ and the
head-$1$ stretches, counted with multiplicity. Consequently, if such closed walks $C_1,\dots,C_n$ use every hop of
either head at most once in all, the bound they yield through \eqref{eq:dep-Zq},
$\cost(G)\ge(n_{\mathrm h}+\sum_i\lVert \delta(C_i)\rVert_q)/q$, is at most $2n_{\mathrm h}/q$.
\end{enumerate}
\end{lemma}

\begin{informalproof}
(a) Since $\prev_2$ is a bijection of $\Z_T$, every time~$t$ is the later endpoint of exactly one hop of head~$2$ and
the earlier endpoint of exactly one; the first enters $x_1(t)$ in~$M_2$ and the second leaves it.

(b) Consecutive hops of the list share an endpoint in time, so the head-$1$ cells match, and the last hop ends at
$t_0+T$, where head~$1$ is at $x_1(t_0)$. The sum of the $\Delta\hat\kappa_1$ telescopes to
$\hat\kappa_1(t_0+T)-\hat\kappa_1(t_0)=n_1(x_1(t_0))$, and there are $n=n_2(c')$ hops. As $x_2$ is a time shift
of~$x_1$, $n_2(c')=n_1(c')$, which is $r(2m+2)$ for interior cells and $r(m+1)$ for tips.

(c) By \cref{lem:transfer}(a), $\delta$ is congruent to $\sum_je(h_j)\pm A_1\pm\dots\pm A_\ell$, whose absolute value
is at most $\sum_je(h_j)+\sum_jA_j$. Next, $\hat\kappa_1(p_{j+1})-\hat\kappa_1(u_j)$ is plus or minus the number of
head-$1$ hops on the $j$-th stretch, so $\delta=-\ell-\sum_j(\hat\kappa_1(p_{j+1})-\hat\kappa_1(u_j))$ and $\abs\delta$
is at most the number of hops on the closed walk. If the walks $C_i$ use every hop at most once, the errors they involve
are distinct terms of $\sum_he_h$, so \eqref{eq:dep-Zq} gives $q\cost(G)\ge n_{\mathrm h}+\sum_he_h\ge
n_{\mathrm h}+\sum_i\lVert \delta(C_i)\rVert_q$; and $\sum_i\lVert \delta(C_i)\rVert_q\le\sum_i\abs{\delta(C_i)}\le
n_{\mathrm h}$.
\end{informalproof}

\paragraph{Why they matter.} \Cref{lem:Zq} makes the depth problem on the cover an exact accounting problem: a label
pays $1/q$ per hop, its errors, and its descents, and the error of a hop is a residue modulo~$q$. \Cref{lem:chain}
says that a label is the visit index of either head plus a cell correction plus that head's integrated errors, and
\cref{lem:transfer} says how the two descriptions interact: around a closed walk of head-$2$ hops, the defect of
head~$1$'s visit index must be paid by head~$2$ or absorbed by head~$1$'s errors on backward stretches. By
\cref{lem:dep-euler}(b), the chain of a head-$2$ cell has zero defect unless exactly one of that cell and the head-$1$
cell of its first visit is a tip, and by \cref{lem:dep-euler}(c) lower bounds obtained from hop-disjoint cycles in
this way are bounded by twice the floor: the identity alone does not give robustness, but it says exactly
what a proof has to count.

\subsection{Labels exact on one head}\label{sec:dep-exact}

The one-head mechanism of \cref{lem:dep-visit} counts exactly along one head; adding an arbitrary correction per cell
gives the natural class of labels exact on that head. For this class the transfer identity has no absorption, and
negative cycles force a price that does not depend on~$q$.

\begin{theorem}[Theorem E: labels exact on one head, abstract form]\label{thm:E}
\formal{DepthLabels.lean}{cost_ge_of_negative_cycles, le_of_modEq_neg}
\begin{enumerate}[label=(\alph*)]
\item If $q\ge1$, $E\ge0$ and $\delta<0$ are integers with $E\equiv\delta\pmod q$, then $E\ge q+\delta$.
\item Let $\tau$ be a finite set, $q\ge1$, $G\colon\tau\to\{0,\dots,q-1\}$, and let $t\mapsto t^-$, $\prev_1$ and
$\prev_2$ be bijections of~$\tau$. Let $\mathrm{cell}\colon\tau\to C$, $S\subseteq C$ and $\kappa\colon\tau\to\Z$ be
such that $G-\kappa$ is constant modulo~$q$ on the times of each cell of~$S$. Let $I$ be a finite set and, for $i\in I$,
let $b_i\colon\Z_{\ell_i}\to\tau$ be maps such that $(i,j)\mapsto b_i(j)$ is injective,
$\mathrm{cell}(b_i(j))=\mathrm{cell}(\prev_2(b_i(j+1)))$ and $\mathrm{cell}(b_i(j))\in S$ for all $i,j$, and
\[
  \delta_i:=\sum_{j\in\Z_{\ell_i}}\bigl(\kappa(b_i(j))-\kappa(\prev_2(b_i(j)))-1\bigr)<0\qquad(i\in I).
\]
Then
\begin{multline*}
  2\card\tau+\sum_{i\in I}(q+\delta_i)\\
  \le\;q\Bigl(\#\{t:G(t)\le G(\prev_1t)\}+\#\{t:G(t)\le G(\prev_2t)\}+\#\{t:G(t)<G(t^-)\}\Bigr).
\end{multline*}
\end{enumerate}
\end{theorem}

\begin{outline}
(a) Write $E-\delta=qz$; then $z\ge1$ because $E\ge0>\delta$, so $E\ge q+\delta$. (b) Apply \lean{card_violated_eq} to
each head, with the hops $(\prev_ht,t)$, $t\in\tau$, and the successor bijection $\prev_h^{-1}$, and apply
\lean{card_descents_eq}; all reduced errors are nonnegative. For each~$i$, the reduced errors of the hops
$(\prev_2(b_i(j)),b_i(j))$ have a sum congruent to the unreduced sum, which is congruent to $\delta_i$ by
\lean{cycle_errors_congr}; so by~(a) it is at least $q+\delta_i$. By injectivity these hops are distinct, so the sum
over~$i$ is at most the total reduced error of head~$2$. Linear arithmetic combines the three identities.
\end{outline}

On the cover of a deep-nesting pair take $\tau=\Z_T$, $t^-=t-1$, the hop maps $\prev_1,\prev_2$, $\mathrm{cell}=x_1$,
$S$ the set of interior cells and $\kappa=\kappa_1$; the right side is then $q\cost(G)$, and $2\card\tau=n_{\mathrm h}$.
If $q$ divides $r(2m+2)$, every label exact on all head-$1$ hops at interior cells satisfies the hypothesis on
$G-\kappa_1$, by \cref{lem:chain}(c). A family $(b_i)$ is a set of hop-disjoint closed walks of head-$2$ hops in~$M_2$
through interior head-$1$ cells, each with negative defect. A closed walk with $-q<\delta_i<0$ forces head-$2$ error at
least $q+\delta_i$, nearly~$q$, on its hops, which by \eqref{eq:dep-Zq} costs nearly one unit, whatever~$q$ is. The
formal statement needs only $\delta_i<0$; for $\delta_i\le-q$ the contribution $q+\delta_i$ is not positive and lowers
the bound, so such walks should be left out of the family. Here $\delta_i$ is computed from the residues~$\kappa$ on the
cover, as in the Lean statement; it differs from the lifted defect of \cref{lem:transfer}(a) by $n_1(\cdot)$ on every
hop that wraps around the cover, and in Theorem~E$_1$ no hop wraps.

\paragraph{Why it matters.} \Cref{prop:packing} and \cref{lem:dep-euler}(c) bound every packing certificate for all
labels by about twice the floor. Theorem~E escapes this because the class is restricted: exactness on one head removes
all absorption (\cref{lem:transfer}), so each negative cycle costs about one unit at every depth. A packing of such
cycles with positive density therefore gives a price per hop independent of~$D$, which rules out the whole class as a
depth-reduction mechanism for the instance; the next theorem supplies such a packing in closed form.

\begin{theorem}[Theorem E$_1$]\label{thm:E1}
\notformalized
Let $s_0\ge4$, $\Lambda=2^{s_0}$, let $x_1=U(s_0,1,1)$, of period $P=8\Lambda$, and let $x_2(t)=x_1(t+\xi)$ with
$\xi=3\Lambda+\eta$ for an integer $\eta$ with $0<8\eta<\Lambda$. Let $r,q\ge1$ with $q\mid4r$. Then every label~$G$ on
the $r$-fold cover (defined as in \cref{def:dnp} for this shift) that is exact on all hops of head~$1$ at interior
cells satisfies
\[
  \cost(G)\;\ge\;\nu_1n_{\mathrm h},\qquad\nu_1=\frac{\Lambda/2-4\eta-1}{16\Lambda},
\]
and $\cost(G)\ge\nu_1n_{\mathrm h}+(1-4\nu_1)n_{\mathrm h}/q$ if $q>4$. For the golden shift of \cref{def:dnp} the
condition $0<8\eta<\Lambda$ holds for every $s_0\ge4$; for $s_0\ge5$, $\nu_1\ge67/4096\approx0.01636$, with equality
exactly at $s_0=8$, and for $s_0=4$, $\nu_1=3/256$.
\end{theorem}

\begin{informalproof}
\emph{The walk.} By \cref{def:dep-walk} with $m=g=1$, each level-$0$ stroke is one segment of $\Lambda$ steps followed
by a V of two plain strokes of $\Lambda/2$ steps. So on $[0,4\Lambda]$
\[
  x_1(t)=\begin{cases}
  t,&0\le t\le\Lambda,\\
  2\Lambda-t,&\Lambda\le t\le3\Lambda/2,\\
  t-\Lambda,&3\Lambda/2\le t\le2\Lambda,\\
  3\Lambda-t,&2\Lambda\le t\le3\Lambda,\\
  t-3\Lambda,&3\Lambda\le t\le7\Lambda/2,\\
  4\Lambda-t,&7\Lambda/2\le t\le4\Lambda,
  \end{cases}
\]
and $x_1(t+4\Lambda)=-x_1(t)$ for all~$t$. In $[0,8\Lambda)$, a cell~$c$ with $0<c<\Lambda/2$ is visited at the times
$c$, $3\Lambda-c$, $3\Lambda+c$, $4\Lambda-c$, and a cell with $\Lambda/2<c<\Lambda$ at $c$, $2\Lambda-c$,
$\Lambda+c$, $3\Lambda-c$; we call $0,1,2,3$ the \emph{indices} of these visits. A cell~$c$ with $-\Lambda<c<0$ is
visited at the times $4\Lambda+t$, for the visits $t$ of $-c$ just listed, with the same indices. Each of these cells
is interior and has $4$ visits per period; so for a time~$t$ of the $k$-th period copy $[8\Lambda k,8\Lambda(k+1))$ of
the cover at such a cell, $\kappa_1(t)$ is $4k$ plus the index of the visit. A pair $(p,u)$ is a hop of head~$2$ if and
only if $(p+\xi,u+\xi)$ is a hop of head~$1$.

\emph{The family.} For every integer~$a$ with $1\le a\le\Lambda/2-4\eta-1$ consider the four head-$2$ hops of the
following table, where the third column is the head-$1$ hop they come from, with its cell, and the last two columns
give the cells of head~$1$ at the two endpoints and the difference of the indices of these visits.
\begin{center}
\small
\begin{tabular}{@{}llllc@{}}
\toprule
& head-$2$ hop $(p,u)$ & head-$1$ hop $(p+\xi,u+\xi)$; its cell & $x_1(p)\to x_1(u)$ & $\Delta\kappa_1$\\
\midrule
$e_1$ & $(a,\Lambda-a-2\eta)$ & $(3\Lambda+a+\eta,4\Lambda-a-\eta)$; $a+\eta$ & $a\to\Lambda-a-2\eta$ & $0$\\
$e_2$ & $(2\Lambda+a+2\eta,3\Lambda-a-4\eta)$ & $(5\Lambda+a+3\eta,6\Lambda-a-3\eta)$; $a+3\eta-\Lambda$
& $\Lambda-a-2\eta\to a+4\eta$ & $-2$\\
$e_3$ & $(3\Lambda-a-4\eta,3\Lambda+a+2\eta)$ & $(6\Lambda-a-3\eta,6\Lambda+a+3\eta)$; $a+3\eta-\Lambda$
& $a+4\eta\to a+2\eta$ & $1$\\
$e_4$ & $(3\Lambda-a-2\eta,3\Lambda+a)$ & $(6\Lambda-a-\eta,6\Lambda+a+\eta)$; $a+\eta-\Lambda$
& $a+2\eta\to a$ & $1$\\
\bottomrule
\end{tabular}
\end{center}
The bounds on~$a$ give $0<a+\eta<a+2\eta<a+3\eta<a+4\eta<\Lambda/2$, which places every cell below in the stated range.
Row~$e_1$: the cell $a+\eta$ lies in $(0,\Lambda/2)$ and is visited at $3\Lambda+(a+\eta)$ and $4\Lambda-(a+\eta)$ with
indices $2$ and~$3$, consecutively; at the times $a$ and $\Lambda-a-2\eta$ head~$1$ is on its first stroke, at cells
$a\in(0,\Lambda/2)$ and $\Lambda-a-2\eta\in(\Lambda/2,\Lambda)$, both with index~$0$. Row~$e_2$: the cell
$\Lambda-a-3\eta$ lies in $(\Lambda/2,\Lambda)$ and is visited at $\Lambda+a+3\eta$ and $2\Lambda-a-3\eta$ with indices
$1$ and~$2$, so its negative $a+3\eta-\Lambda$ is visited at the times of the table with indices $1$ and~$2$,
consecutively; at $2\Lambda+a+2\eta=3\Lambda-(\Lambda-a-2\eta)$ head~$1$ visits
$\Lambda-a-2\eta\in(\Lambda/2,\Lambda)$ with index~$3$, and
at $3\Lambda-(a+4\eta)$ it visits $a+4\eta\in(0,\Lambda/2)$ with index~$1$. Row~$e_3$: the visits of
$a+3\eta-\Lambda$ at $6\Lambda-a-3\eta$ and $6\Lambda+a+3\eta=4\Lambda+3\Lambda-(\Lambda-a-3\eta)$ have indices
$2$ and~$3$; at $3\Lambda+(a+2\eta)$ head~$1$ visits $a+2\eta\in(0,\Lambda/2)$ with index~$2$. Row~$e_4$: the cell
$\Lambda-a-\eta$ lies in $(\Lambda/2,\Lambda)$ and is visited at $2\Lambda-a-\eta$ and $2\Lambda+a+\eta$ with indices
$2$ and~$3$, so its negative $a+\eta-\Lambda$ is visited at the times of the table with indices $2$ and~$3$; at
$3\Lambda-(a+2\eta)$ head~$1$ visits $a+2\eta$ with index~$1$, and at $3\Lambda+a$ it visits $a$ with index~$2$. In
each row $p<u$, and all eight times lie in $[0,4\Lambda)$, so the
differences of~$\kappa_1$ in any period copy are the differences of the indices.

\emph{Closure and defect.} The head-$1$ cells run $a\to\Lambda-a-2\eta\to a+4\eta\to a+2\eta\to a$, so
$x_1(u_j)=x_1(p_{j+1})$ for $j\in\Z_4$, identically in $a$, $\Lambda$ and~$\eta$, and all four cells are interior. The
defect is $(0-1)+(-2-1)+(1-1)+(1-1)=-4$. (At $a=0$ and at $a=\Lambda/2-4\eta$ other negative cycles arise, with other
index patterns; we do not use them.)

\emph{The mirror family.} Since $x_1(t+4\Lambda)=-x_1(t)$ for all~$t$, also $x_2(t+4\Lambda)=-x_2(t)$, so adding
$4\Lambda$ to all times of a member gives four head-$2$ hops whose head-$1$ cells are negated. The cycle condition is
preserved, the indices are the same by the description of the visits of negative cells, and the defect is again~$-4$.

\emph{Disjointness.} A hop of head~$2$ on the cover is determined by its later endpoint. In $[0,8\Lambda)$ the later
endpoints of $e_1$, $e_2$, $e_3$, $e_4$ lie in $(\Lambda/2,\Lambda)$, $(5\Lambda/2,3\Lambda)$, $(3\Lambda,7\Lambda/2)$
and $(3\Lambda,7\Lambda/2)$ respectively, and those of the mirror family in the same intervals shifted by $4\Lambda$;
these sets are disjoint except for the last two. For different values of~$a$ the endpoints of $e_1$ differ, as do those
of $e_2$, of~$e_3$ and of~$e_4$, and $3\Lambda+a+2\eta=3\Lambda+a'$ exactly when $a'=a+2\eta$. So two members share a
hop only if their parameters differ by~$2\eta$, the shared hop being $e_3$ of the smaller and $e_4$ of the larger, and
within a member the four hops are distinct. Let $n=\Lambda/2-4\eta-1$ be the number of members of each family. The
members split into chains $a,a+2\eta,a+4\eta,\dots$; keeping every other member of each chain, starting with the
first, keeps $\lceil\ell/2\rceil$ members of a chain of~$\ell$, hence at least $n/2$ members of each family, and the
kept members are pairwise hop-disjoint. With the mirror family this gives at least~$n$ hop-disjoint cycles per period.

\emph{Counting.} On the $r$-fold cover, take these cycles in each of the $r$ period copies, translated by multiples of
$8\Lambda$; they remain hop-disjoint, each has defect $-4$ computed with~$\kappa_1$, and their number $n_{\mathrm c}$ is
at least $rn$. Interior cells have $4r$ visits per cover and $q\mid4r$, so by \cref{lem:chain}(c) the label~$G$
satisfies the hypothesis of \cref{thm:E}(b) with $\kappa=\kappa_1$. That theorem gives
$q\cost(G)\ge2T+n_{\mathrm c}(q-4)$, that is, $\cost(G)\ge n_{\mathrm c}+(2T-4n_{\mathrm c})/q\ge n_{\mathrm c}$,
because the $n_{\mathrm c}$ cycles use $4n_{\mathrm c}$ distinct hops among the $T$ hops of head~$2$. Since
$n_{\mathrm h}=2T=16r\Lambda$ and $n_{\mathrm c}\ge rn$, this is $\cost(G)\ge\nu_1n_{\mathrm h}$. If $q>4$, the right
side $n_{\mathrm c}(1-4/q)+2T/q$ increases with~$n_{\mathrm c}$, and $n_{\mathrm c}\ge2T\nu_1$ gives
$\cost(G)\ge\nu_1n_{\mathrm h}+(1-4\nu_1)n_{\mathrm h}/q$.

\emph{The golden shift.} Here $P=8\Lambda$ and $0.381966\,P=3.055728\Lambda$. For $s_0=5,6,7,8$ one finds
$\lfloor3.055728\Lambda\rfloor=97,195,391,782$, so $\xi=97,195,391,783$, $\eta=1,3,7,15$, $8\eta<\Lambda$, and
$\nu_1=11/512$, $19/1024$, $35/2048$, $67/4096$. For $s_0\ge9$, $3.055728\Lambda-1<\xi\le3.055728\Lambda+1$, so
$0.055728\Lambda-1<\eta\le0.055728\Lambda+1$; hence $0<8\eta<\Lambda$ and
$\nu_1\ge(0.277088-5/\Lambda)/16\ge0.0167>67/4096$, as $\Lambda\ge512$.
\end{informalproof}

\paragraph{Why it matters.} Theorem~E$_1$ is the first price per hop that does not depend on the depth, on the
sharpest open instance, with one level of nesting, for a natural class of labels: the class that extends the mechanism solving one walk by an
arbitrary correction per cell. On the golden pair with $m=1$ and $s_0\ge5$, such labels cost at least $0.016$ per hop on the cover, about
$0.033\,T$, at every depth~$D$, so they cannot give $\cost\le\eps T$ for $\eps<0.03$; for the aligned shift $\xi=1$ the
analogous certificates are close to zero numerically. The theorem says nothing about labels with small but nonzero
errors on
head~$1$, and these are exactly what optimised labels use (\cref{rem:dep-absorb}).

\begin{remark}[From the cover to the walk DAG]\label{rem:dep-linear}
The same argument works on the walk DAG of the pair of \cref{thm:E1} on $[0,NP)$, $N\ge1$, for labels
$\Psi\colon[0,NP)\to\{0,\dots,q-1\}$ that are exact modulo~$q$ on all head-$1$ hops at interior cells. The chains of
hops no longer close up, so summing $q[\Psi(u)\le\Psi(p)]=1+e-(\Psi(u)-\Psi(p))$ over the hops of a chain gives
$q\cdot\#\{\text{violated hops}\}\ge\#\{\text{hops}\}+\sum e-(q-1)\#\{\text{chains}\}$; there are at most
$2(2\Lambda+1)$ chains. Along each interior head-$1$ chain
$\Psi-\kappa_1$ is constant modulo~$q$, with no divisibility condition, and the cycles of the family that lie in
$[0,NP)$, at least $Nn$ after selection, force head-$2$ error at least $q-4$ each. Hence, counting hops per head,
$\cost(\Psi)\ge2\nu_1NP-O(\Lambda)$ for $q>4$, and $\cost(\Psi)\ge NP/2-O(\Lambda)$ for $q\le4$ from the floor. A
pair $(p,u)$ that is a hop of both heads would be a single edge of the walk DAG, but there is none. Indeed, all turn
times and turn positions of $U(s_0,m,g)$ are multiples of $2^{s_m}$, so on every maximal monotone run $x(t)=\pm t+c$
with $2^{s_m}\mid c$. A hop $(p,u)$ leaves $x(p)$ at time~$p$ and returns at time~$u$ along runs of opposite slopes,
so $p+u\equiv0\pmod{2^{s_m}}$, while a hop of $x_2$ satisfies $p+u\equiv-2\xi$; and $2\xi\not\equiv0\pmod{2^{s_m}}$ for
odd~$\xi$ when $s_m\ge2$; for the pair of \cref{thm:E1}, where $\xi=3\Lambda+\eta$ need not be odd, $2^{s_1}=\Lambda/2$
and $2\xi\equiv2\eta\not\equiv0\pmod{\Lambda/2}$ because $0<2\eta<\Lambda/4$.
\end{remark}

\begin{remark}[Numerical: nesting depths $m\le6$]\label{rem:dep-Em}
Theorem~E$_1$ is a closed form of computer-found certificates. For the golden pair $\mathrm{DNP}(s_0,m,1)$ with $1\le
m\le6$, explicit one-parameter families of negative cycles in~$M_2$, fitted as affine functions of~$a$, $\Lambda$
and~$\xi$ on ranges of~$s_0$ between $9$ and~$14$ and checked member by member on instances outside the fitting range
(up to $s_0=16$), certify for labels exact on all head-$1$ hops at interior cells, at every~$q$, prices per hop of
$0.0509$, $0.0554$, $0.0527$, $0.0519$, $0.0526$, $0.0480$ for $m=1,\dots,6$ (from $12,16,95,208,245,531$ families,
with the fractional packing that gives each member $C$ the weight $1/\max_{h\in C}\mathrm{load}(h)$, where the load of a
hop is the number of members through it). These values were stable in~$s_0$; they are per-instance certificates, not proofs
for all~$s_0$, and they have not been checked independently. For $m=1$ a greedy selection of hop-disjoint
members does better: it finds exactly $\Lambda-4$ negative $4$-cycles per period on all instances checked ($7\le s_0\le14$
and $s_0=16$), that is, a price $1/16-1/(4\Lambda)$ per hop at every~$q$; a proof would need the coincidences between the
twelve families. For labels exact only on
the head-$1$ hops of some levels, the certified price falls as levels are released: at $m=1$ exactness on the finest
level alone admits an explicit label of cost at most about $4/q$ per hop, while at $m=3$ and $s_0=7$, exactness on the
finest one, two or three levels or on all four carries certified prices of about $0.003$, $0.026$, $0.066$ and
$0.107$. Why ``every~$m$'' is open: the families show no recursion in~$m$, and the phases of the shift relative to the
level grids, $\psi_i=\mathrm{frac}(2^i\xi/2^{s_0})$, which are the same for all strokes of level~$i$, degenerate for
special~$m$; for instance $\psi_2,\dots,\psi_6$ lie between $0.003$ and $0.06$ at $m=8$, and when $15625$ divides
$m+1$ exact arithmetic gives $\psi_i=2^{i-s_0}$ for every $i\ge3$. A proof for all~$m$ therefore needs a rule
generating level-$i$ families from local configurations, a bound on their overlaps independent of~$m$, and a positive
share of good levels, not equidistribution of phases.
\end{remark}

\subsection{The remaining gap}\label{sec:dep-gap}

Labels found by optimisation are not exact on either head. They trade errors of head~$1$ on backward stretches
against errors of head~$2$ (\cref{rem:dep-absorb}), and by \cref{lem:transfer}(c) such an absorption costs only about
$\abs\delta/q$ per cycle. A lower bound for all labels must show that absorptions cannot be supplied consistently: an
error on a head-$1$ hop lies on backward stretches of some head-$2$ cycles, where it helps, and on forward stretches of
others, where it hurts, and symmetrically for head~$2$. We expect that this forces a multiplicative loss at every level
that both heads use, and we state the expected outcome as a count of served levels.

Let the pair of \cref{def:dnp} be defined, verbatim, for any odd shift~$\xi$. For a label~$G$ on its cover and
$0<\theta<1$, say that level~$i$ is \emph{$\theta$-served for head~$h$} if the hops of head~$h$ of level~$i$ have total error at most $\theta qn_{\mathrm h}/(2(m+1))$, and
\emph{doubly $\theta$-served} if it is $\theta$-served for both heads. When the $n_{\mathrm h}/2$ hops of each head are
spread evenly over the $m+1$ levels, as numerical counts indicate, this says that the mean error per hop of the class
is at most~$\theta q$. Two kinds of labels serve every level for trivial reasons; the conjecture excludes the first by
its bound on the cost and the second by its choice of the shift. Labels with few or random values serve every level at
large~$\theta$, but cost at least $n_{\mathrm h}/q$, more than $n_{\mathrm h}/8$ if $q<8$, respectively about
$n_{\mathrm h}/2$, by \eqref{eq:dep-Zq}. And for a shift aligned with the dyadic scales, such as the golden shift
when $15625$ divides $m+1$, a label exact on one head is nearly exact on the other.

\begin{conjecture}[Conjecture A$'$: few doubly served levels]\label{conj:Aprime}
There are $\theta_0>0$ and constants $C,C'$, and for all $s_0$ and~$m$ an odd shift~$\xi$, such that for all $r$
and~$q$ every label~$G$ with $\cost(G)\le n_{\mathrm h}/8$ on the $r$-fold cover of the pair $x_1=U(s_0,m,1)$,
$x_2=x_1(\cdot+\xi)$ has at most $C\log_2q+C'$ doubly $\theta_0$-served levels.
\end{conjecture}

We expect shifts misaligned with the dyadic scales at every level to qualify. The golden shift is not always one of
them: even when $15625\nmid m+1$, its phases nearly align at several levels for some~$m$, such as $m=8$
(\cref{rem:dep-Em}).

If the conjecture holds, let $G$ be a label, on a cover of the pair with the shift that it provides, of cost
$\cost(G)\le\eps n_{\mathrm h}$, where $\eps\le1/8$ and
$\eps<\theta_0/2$. By \cref{cor:dep-split}, every level that is not doubly $\theta_0$-served contributes more than
$\theta_0n_{\mathrm h}/(2(m+1))$ to the cost, so fewer than $2\eps(m+1)/\theta_0$ levels are not doubly served, and
$(m+1)(1-2\eps/\theta_0)<C\log_2q+C'$. This gives
\[
  D=q-1\;\ge\;\exp\bigl(c(1-2\eps/\theta_0)(m+1)-c'\bigr)-1
\]
for constants $c,c'>0$, so the depth needed at a fixed cost per hop grows exponentially in the depth of nesting,
matching the clock labels of \cref{prop:dep-clock} up to the constant in the exponent. The same statement for labels
on the walk DAG of many periods would bound the depth of every cut of size $\eps T$ in the same way, because
\eqref{eq:dep-Zq} and \cref{cor:dep-split} hold there up to boundary terms and, when $s_m\ge2$, no pair is a hop of
both heads (both as in \cref{rem:dep-linear}). It would make the family of these pairs with $s_m\ge2$, $m\to\infty$,
depth-robust, and \cref{prob:tcp} for $k=2$ would have a negative answer.

\begin{remark}[Numerical: evidence for Conjecture~A$'$]\label{rem:dep-Aevidence}
The evidence is limited. In the best labels found by optimisation for $m=3$ ($g=1$, $s_0=6$), with mean error at most
$q/10$ per hop counted as served, the number of doubly served levels was $1$, $1$, $2$ and $2$ or~$3$ at
$D=7,15,31,63$, consistent with growth like $\log q$, and level~$0$ of head~$1$ was never served (mean error about
$0.17q$ to $0.20q$ at every~$D$). Labels exact on several levels of one head carry a price that grows with the number
of exact levels (\cref{rem:dep-Em}), a restricted form of the conjecture. Small instances cannot decide asymptotic
questions, and these labels are upper bounds of uncertain quality.
\end{remark}

\paragraph{Status.} We do not know whether the deep-nesting pairs are depth-robust as $m$ grows. No mechanism we know
cuts them at a depth independent of~$m$: per-head band rules appear to fail (proof sketches of \cref{rem:dep-N,rem:dep-Nprime}); for $g\ge2$, anchored digits
cannot be monotone (\cref{lem:dep-monotone}); and explicit combinations of visit indices and band counts kept a cost per
hop bounded away from zero at every depth tried, in computations up to $m=8$. Clock labels cut each pair, but only at depth $(c/\eps)^{m+1}$ (\cref{prop:dep-clock}). On the other
side, the only lower bounds that hold at every depth are for restricted classes
(\cref{thm:E,thm:E1}), and packing certificates cannot do better (\cref{prop:packing}).

\paragraph{Why it matters.} \Cref{prob:tcp} is the walk-level form of the structural hypothesis of \cref{thm:D}, and
\cref{thm:B} solves it for home-based walks. On the sharpest open instance, the evidence just listed leaves a robust
multiplicative count of absorptions as the missing step, of which \cref{conj:Aprime} is one form. Either answer is
informative. If the pairs are depth-robust as $m$ grows, $\TCP(2)$ fails as stated; if moreover oblivious linear-time
machines can generate them as dummy traffic and they are depth-robust at every $\eps$ below the crux rates in question, the structural hypothesis of \cref{thm:D} must be weakened as explained in
\cref{sec:dep-problem}, and these pairs, placed at intermediate scales inside curved motion, are candidates for the depth-robust schedule that
an oblivious linear-time transposer obeying the crux would need. If they are depth-reducible at a depth independent
of~$m$, the mechanism must beat clock
labels by absorbing defects consistently across levels, and it would be the core of a general proof of $\TCP(2)$.

\section{Storage codes on sparse side information}\label{sec:codes}

The second hypothesis of Theorem~D is an all-codes index-coding statement, the crux (\cref{def:crux}): on the side information
produced by a bounded-depth genie, every index code for the transposition demand must broadcast a constant fraction of
the input. An oblivious linear-time transposer whose dependency DAG has cheap genies of bounded depth needs the
opposite, broadcasts of $\eps L$ bits with $\eps$ as small as the trade-off of \cref{prop:tradeoff} allows. Rate~$\tfrac12$ is the natural threshold: when the receivers can be paired so
that each knows its partner's demand, broadcasting the XOR of each pair serves everyone with $L/2$ bits, and no
construction we know does better on bounded-degree side information of growing odd girth.

In the hardest instances we know, coming from schedules whose head motion expands (the ones \cref{sec:motion} leaves
open), the side information that matters is symmetric and sparse: two receivers that know each other's demands form an
edge of a bounded-degree graph that looks locally like a tree. This is a heuristic description, not a theorem, but it
singles out the following question, the undirected form of the crux. By the duality between index codes and storage codes
(\cref{sec:prelim-codes}), it asks whether storage codes can beat rate~$\tfrac12$.

\begin{question}[Undirected crux]\label{prob:undirected-crux}
Let $(G_n)$ be graphs of maximum degree at most $\Delta$ that are \emph{locally bipartite}: for each fixed $\ell$, only
$o(\card{V(G_n)})$ vertices lie on odd cycles of length at most $\ell$. Does every storage code on $G_n$ have rate at most
$\tfrac12+o(1)$, over every alphabet?
\end{question}

A negative answer is necessary for the kind of coding gain a linear-time transposer would exploit on such side
information; a positive answer, in the directed and depth-bounded generality of \cref{def:crux}, would combine with depth
reduction to refute oblivious~(T). This section proves the positive answer for restricted classes of codes and of
graphs, and shows where the remaining difficulty lies:
\begin{itemize}[leftmargin=2em]
\item every storage code, linear or not, on any graph that maps homomorphically to $C_5$, to the circular clique
$K_{14/5}$ or to the Petersen graph has rate at most $\tfrac12$, with no degree assumption (\cref{cor:H2});
\item on triangle-free graphs of maximum degree~$3$, no linear storage code over any field beats rate~$\tfrac12$
(Theorem~S, \cref{thm:S});
\item with bounded complex coefficients, no linear code beats rate $\tfrac12+o(1)$ on locally bipartite graphs of
bounded degree, and over finite fields, with a fixed number of symbols per vertex, a gain must come from $p$-torsion
(Theorem~A, \cref{thm:A,cor:torsion});
\item at degree~$5$ the threshold fails: the Clebsch graph carries a binary linear code of rate $5/8$, and no storage code
does better, so its guessing number is exactly~$10$ (\cref{cor:H4});
\item local arguments inside tree-like regions cannot give rate $\tfrac12+o(1)$ at degree at least~$4$, real eigenvalue
codes cannot reach rate~$\tfrac12$ at girth at least~$7$, and Mycielski's construction does not preserve goodness;
whether goodness bounds the fractional chromatic number is open (\cref{sec:ana-limits,sec:ana-targets}).
\end{itemize}
Every numbered result of \cref{sec:entropic,sec:thmH,sec:symmetrization,sec:targets,sec:clebsch,sec:subcubic} is
formalized, apart from the classical Tutte--Berge formula used in the matching form of Theorem~S (\cref{thm:S-barrier});
\cref{sec:ana-A,sec:ana-fields,sec:ana-limits,sec:ana-targets} are not formalized.

\subsection{Entropic polymatroids}\label{sec:entropic}

\begin{lemma}[Entropy of a recoverable word]\label{lem:entropic}
\formal{StorageCode.lean}{codeEntropy, codeEntropy_isPolymatroid, codeEntropy_isValid}
Let $X\colon\Omega\to\alpha^V$ be a random word on a finite probability space. Then $S\mapsto H(X_S)$ is a polymatroid on
$V$. If $X$ is recoverable for $N$, this polymatroid is valid for~$N$.
\end{lemma}

\begin{outline}
$H(X_S)$ is the entropy of the observation $\omega\mapsto X(\omega)|_S$, written as a word in $(\alpha\cup\{\bot\})^V$ that
is $\bot$ outside $S$, in the second paper's entropy library~\cite{SecondPaper2026}. The empty observation is constant,
so it has entropy~$0$. Monotonicity holds because $X_S$ is a function of $X_T$ for $S\subseteq T$. Submodularity is that
library's strong subadditivity for the observations $X_S$, $X_T$ and $X_{S\cap T}$. Validity: recoverability makes
$X_{N(v)\cup\{v\}}$ a function of $X_{N(v)}$, so the two entropies agree.
\end{outline}

\begin{theorem}[Rate bounds bound codes]\label{thm:entropic-rate}
\formal{StorageCode.lean}{codeEntropy_univ_le_of_rateBound, codeEntropy_univ_uniform, log_card_le_of_rateBound}
Suppose $N$ satisfies the rate bound $c\ge0$. Then every random word $X$ recoverable for $N$ satisfies
$H(X_V)\le c\sum_vH(X_v)$, and every storage code $C$ for $N$ over a nonempty alphabet $\alpha$ satisfies
$\log\card C\le c\,\card V\log\card\alpha$.
\end{theorem}

\begin{outline}
The first claim applies the rate bound to the polymatroid of \cref{lem:entropic}. For the second, take $X$ uniform on
$C$; then $H(X_V)=\log\card C$, and each $H(X_v)\le\log\card\alpha$.
\end{outline}

\paragraph{Why it matters.} This is the bridge from polymatroids to codes: a rate bound proved for valid polymatroids,
by a certificate or by Theorem~H, bounds every storage code, linear or not, over every alphabet. For~(T), a rate
bound~$c$ on the side information of the cones gives the crux at rate~$1-c$ (\cref{rem:dec-crux}(b)).

Valid polymatroids inherit the recovery structure as a closure operation, which every certificate below uses.

\begin{lemma}[Closure]\label{lem:closure}
\formal{Polymatroid.lean}{IsPolymatroid.eq_of_closure, IsPolymatroid.insert_eq_of_subset}
Let $h$ be a polymatroid valid for $N$, and let $S\subseteq T$ be such that $N(v)\subseteq S$ for every
$v\in T\setminus S$. Then $h(T)=h(S)$.
\end{lemma}

\begin{outline}
Submodularity gives diminishing returns: if adding $v$ to $A$ does not increase $h$, then adding $v$ to any $B\supseteq A$
does not either. Apply this with $A=N(v)$ and $B\supseteq S$, adding the vertices of $T\setminus S$ one at a time.
\end{outline}

\subsection{Theorem H: rate bounds pull back along homomorphisms}\label{sec:thmH}

\begin{theorem}[Theorem H]\label{thm:H}
\formal{Polymatroid.lean}{pushRank, IsPolymatroid.pushRank, IsValid.pushRank, pushRank_univ, sum_pushRank_singleton_le, RateBound.of_hom, RateBound.of_graphHom, RateBound.of_digraphHom}
Let $N_G$ and $N_T$ be side information on $V$ and $W$, and let $\varphi\colon V\to W$ map $N_G(u)$ into $N_T(\varphi(u))$
for every $u$. For a polymatroid $h$ on $V$ put $h'(Z)=h(\varphi^{-1}(Z))$. Then $h'$ is a polymatroid on $W$, valid for
$N_T$ if $h$ is valid for $N_G$, with $h'(W)=h(V)$ and $\sum_th'(\{t\})\le\sum_vh(\{v\})$. Consequently, if $N_T$ satisfies
the rate bound $c\ge0$, so does $N_G$. In particular this holds for every graph homomorphism $G\to T$ and every digraph
homomorphism (arcs to arcs).
\end{theorem}

\begin{outline}
Preimages preserve unions, intersections and inclusions, so $h'$ is a polymatroid. Validity: if $u$ lies in the fibre of
$t$, then $N_G(u)\subseteq\varphi^{-1}(N_T(t))$, so by \cref{lem:closure} adding the whole fibre of $t$ to
$\varphi^{-1}(N_T(t))$ does not increase~$h$. No looplessness or injectivity is needed. The fibres partition $V$, and
subadditivity bounds $h$ of each fibre by the sum over its vertices. Finally
$h(V)=h'(W)\le c\sum_th'(\{t\})\le c\sum_vh(\{v\})$.
\end{outline}

The formal theorem is about all polymatroids, so it concerns the Shannon bound. The same argument shows that the
push-forward of an entropic polymatroid is entropic: if $h(S)=H(X_S)$, then $h'(Z)=H(Y_Z)$ for the random word
$Y_t=X_{\varphi^{-1}(t)}$. We do not need this, because every target below is certified by Shannon inequalities.

\paragraph{Why it matters.} Theorem~H turns one finite certificate into a statement about infinitely many graphs: a rate
bound for a fixed target~$T$ holds for every side information that maps to~$T$, whatever its size or degree. This is the
mechanism behind every all-codes result of this section about infinitely many graphs, and
the reason a target's structure, not the size of the side
information, decides goodness. For~(T), it underlies the homomorphism route to the crux (\cref{sec:ana-targets}). The
theorem is folklore-level (the proof is one line); related monotonicity statements
for index coding appear in~\cite{EbrahimiSiavoshani2014}, and blow-up lemmas in~\cite{Baber2017,CameronDangRiis2016}.

\subsection{Symmetrization}\label{sec:symmetrization}

\begin{theorem}[Symmetrization]\label{thm:symmetrization}
\formal{Symmetry.lean}{IsAut, symmetrize, IsPolymatroid.symmetrize, IsValid.symmetrize, symmetrize_invariant, symmetrize_univ, sum_symmetrize_singleton, RateBound.of_invariant}
Let $N$ be side information on $V$. If every valid polymatroid that is invariant under all automorphisms of $N$ satisfies
$h(V)\le c\sum_vh(\{v\})$, then $N$ satisfies the rate bound~$c$.
\end{theorem}

\begin{outline}
Average a valid polymatroid over the finite group of permutations $\sigma$ of $V$ with $\sigma(N(v))=N(\sigma(v))$. Each
relabelled function $S\mapsto h(\sigma S)$ is a valid polymatroid, so the average is one too. It is invariant, because
right multiplication by an automorphism permutes the group, and it has the same value on $V$ and the same sum over
singletons as~$h$.
\end{outline}

\paragraph{Why it matters.} With symmetrization, a certificate (\cref{sec:targets}) needs one elemental inequality per orbit, and each set it
mentions can be replaced by an
orbit representative through an explicit automorphism. This is what makes the Clebsch certificate below small enough to
check.

\subsection{Good targets}\label{sec:targets}

A certificate for a rate bound $c$ is a nonnegative combination of \emph{elemental inequalities}
$h(K\cup\{i\})+h(K\cup\{j\})\ge h(K\cup\{i,j\})+h(K)$ that equals $c\sum_vh(\{v\})-h(V)$ once every set is replaced by its
closure (\cref{lem:closure}). We found the certificates by linear programming, re-solved them exactly in rational
arithmetic on their supports, and checked the Petersen, $K_{14/5}$ and Clebsch certificates in Lean with
\texttt{linear\_combination} and exact coefficients (the five-inequality pentagon certificate is closed by
\texttt{linarith}); the generator is \leanfile{tools/certificates.py}. The smallest example shows the mechanism.

\begin{proposition}[The pentagon is good]\label{prop:C5}
\formal{Pentagon.lean}{pentagonNbhd, nbhd_cycleGraph_five, rateBound_pentagon, rateBound_cycleGraph_five}
Every valid polymatroid on the $5$-cycle $0\text{-}1\text{-}2\text{-}3\text{-}4\text{-}0$ satisfies $2h(V)\le\sum_vh(\{v\})$.
\end{proposition}

\begin{outline}
Writing $I(i;j\mid K)$ for the elemental quantity above, the five non-adjacent inequalities
\[
  I(0;2)+I(0;3\mid1,2)+I(1;3)+I(1;4)+I(2;4\mid1)\;\ge\;0
\]
sum to $\sum_vh(\{v\})-2h(V)$ after the closures $\{0,1,2\}\to\{0,2\}$, $\{1,2,3\}\to\{1,3\}$,
$\{0,1,2,3\}\to V$ and $\{1,2,4\}\to V$. The Lean proof is exactly this computation.
\end{outline}

The \emph{circular clique} $K_{p/q}$ has vertices $\Z/p$, with $i\sim j$ if and only if $q\le(j-i)\bmod p\le p-q$.

\begin{proposition}[Further good targets]\label{prop:targets}
\formal{Petersen.lean}{petersenGraph, petersenGraph_adj_iff, petersenKneser_card, petersenKneser_injective, rateBound_petersenGraph}
\formal{CircularClique.lean}{circularClique, rateBound_circularClique_14_5, circularClique_adj_of_lt, circularCliqueHom, rateBound_circularClique_of_le, cycleToCircularClique, rateBound_cycleGraph_odd}
The Petersen graph (the Kneser graph $K(5,2)$) and the circular clique $K_{14/5}$ (the Andrásfai graph
$\mathrm{And}(5)$) are good. If $p/q\le p'/q'$ with $q'\ge1$, the map $i\mapsto\lfloor p'i/p\rfloor$ is a homomorphism
$K_{p/q}\to K_{p'/q'}$ (Bondy and Hell~\cite{BondyHell1990}). Hence every circular clique $K_{p/q}$ with $p/q\le14/5$ is
good, for example $K_{8/3}$ (the Möbius ladder $\mathrm{And}(3)$) and $K_{11/4}$ ($\mathrm{And}(4)$), and so is every
odd cycle $C_{2k+1}$ with $k\ge2$.
\end{proposition}

\begin{outline}
The Petersen certificate is a
combination of $58$ elemental inequalities with weights $\tfrac14$ and~$\tfrac12$ and $25$ closure equalities (of the $40$ that the
file proves), and the labelling is identified with $K(5,2)$ by an explicit bijection. For $K_{14/5}$ the polymatroid is first symmetrized over the
dihedral group (\cref{thm:symmetrization}); the certificate then uses $32$ orbit rows, $21$ explicit automorphisms, $116$
invariance equalities and $13$ closure equalities. For the homomorphism, if $i<j$ and $q\le j-i\le p-q$,
then $\lfloor p'i/p\rfloor+q'\le\lfloor p'j/p\rfloor\le\lfloor p'i/p\rfloor+p'-q'$ because $pq'\le p'q$, and adjacency is
symmetric. The odd cycle $C_{2k+1}$ maps to $K_{(2k+1)/k}$ by $i\mapsto ki$, which turns the differences $\pm1$ into
$\pm k$, and $(2k+1)/k\le14/5$ for $k\ge2$. Theorem~H applies.
\end{outline}

\begin{corollary}[Corollary H2: rate $\tfrac12$ for every code]\label{cor:H2}
\formal{RateHalf.lean}{card_sq_le_of_rateBound_half, IsGood.card_sq_le, IsGood.codeEntropy_le}
\formal{Polymatroid.lean}{IsGood.of_hom}
Let $G$ be a good graph, for instance a graph with a homomorphism to $C_5$, to the Petersen graph or to a circular clique
$K_{p/q}$ with $p/q\le14/5$. Then every storage code $C\subseteq\alpha^{V}$ on $G$ over a nonempty alphabet, linear or
not, satisfies $\card C^2\le\card\alpha^{\card V}$, and every random word $X$ recoverable for $G$ satisfies
$2H(X_V)\le\sum_vH(X_v)$.
\end{corollary}

\begin{outline}
Goodness pulls back along homomorphisms (Theorem~H); combine with \cref{prop:C5,prop:targets} and
\cref{thm:entropic-rate}. The integer form takes logarithms.
\end{outline}

The graphs with a homomorphism to some $K_{p/q}$ with $p/q\le14/5$ are exactly the graphs of circular chromatic number at
most $14/5$, since for a finite graph the circular chromatic number is the least $p/q$ such that it maps to
$K_{p/q}$~\cite{Vince1988,Zhu2001}.

A graph with a perfect matching carries the pairing code (copy one symbol across each matched edge), of rate exactly
$\tfrac12$, so for good graphs with a perfect matching the corollary determines the guessing number:
$\gn(G)=\card V/2$. The corollary strictly
extends the guessing-number results of Mazumdar, McGregor and Vorotnikova~\cite[Thms~9--11]{MMV2019} (their graphs map to
$K_{8/3}$ or to $C_5$), and contains the odd-cycle theorem of Christofides and Markström~\cite{ChristofidesMarkstrom2011}.
For Petersen it makes rigorous the numerical values of Blasiak, Kleinberg and Lubetzky~\cite{BKL2013}.

\paragraph{Why it matters.} \Cref{cor:H2} is an all-codes answer to the undirected crux, \cref{prob:undirected-crux}, for
every graph that maps to a good target: the graphs need not have bounded degree, and they include graphs of arbitrarily
large odd girth, for example long odd cycles. The question is
whether high-girth bounded-degree graphs map to good targets. For cubic graphs this is a known open problem:

\begin{corollary}[Corollary H3: a conditional degree-$3$ result]\label{cor:H3}
\formal{RateHalf.lean}{IsGood.card_sq_le}
\formal{Polymatroid.lean}{IsGood.of_hom}
\formal{Pentagon.lean}{rateBound_cycleGraph_five}
\formal{CircularClique.lean}{rateBound_circularClique_14_5}
If every cubic graph of sufficiently large girth maps to $C_5$ (Nešetřil's Pentagon Problem~\cite{Nesetril1999}), or to
$K_{14/5}$, then every storage code on such a graph, nonlinear included, has rate at most $\tfrac12$.
\end{corollary}

\begin{outline}
This is \cref{cor:H2} applied under the stated hypothesis.
\end{outline}

Whether high-girth cubic graphs have circular chromatic number below~$3$ is itself open~\cite{DaneshgarMadani2015}. Some
targets are known to fail: random cubic graphs do not map to $C_7$~\cite{Hatami2005}, and the Clebsch graph, to which every
subcubic graph of girth at least~$17$ maps by a theorem of DeVos and Šámal~\cite{DeVosSamal2006}, is not good
(\cref{cor:H4}). For linear codes, Theorem~S below settles degree~$3$ unconditionally.

\subsection{The Clebsch graph: guessing number ten}\label{sec:clebsch}

The Clebsch graph is the folded $5$-cube: vertices $0,\dots,15$, with $i\sim j$ if and only if $i\oplus j$ (bitwise
exclusive or) is a power of two or equals $15$. It is triangle-free and $5$-regular.

\begin{corollary}[Corollary H4: $\gn(\text{Clebsch})=10$]\label{cor:H4}
\formal{ClebschCertificate.lean}{rateBound_clebsch, rateBound_clebschGraph, clebsch_card_le, clebsch_guessingNumber, clebschCode_eq, mem_ker_clebsch_iff, finrank_ker_clebsch, rank_clebsch}
\formal{Clebsch.lean}{clebschGraph, clebschGraph_adj_iff, clebschCode, clebschCode_card, clebschCode_isStorageCode, exists_clebsch_storageCode}
Every valid polymatroid on the Clebsch graph satisfies $h(V)\le\tfrac58\sum_vh(\{v\})$. Hence every storage code on the
Clebsch graph over an alphabet of size $q$ has at most $q^{10}$ codewords, while the binary linear storage code
$\ker(I+A)$ over $\F_2$, where $A$ is the adjacency matrix, has exactly $2^{10}$. So $\gn(\text{Clebsch})=10$, and
$I+A$ has rank~$6$ over $\F_2$.
\end{corollary}

\begin{outline}
Symmetrize over the automorphism group, of order $1920$ (\cref{thm:symmetrization}). The certificate has $54$ orbit rows;
unlike the certificates for $C_5$, the Petersen graph and $K_{14/5}$, it uses elemental inequalities between adjacent
vertices as well. It uses $126$ explicit
automorphisms, $225$ invariance equalities and $54$ closure equalities, and is checked with exact rational coefficients.
For the lower bound, the code is the span of ten explicit vectors, each satisfying $x_v=\sum_{u\sim v}x_u$ over
$\F_2$, in systematic form on ten coordinates; injectivity of the encoding gives $2^{10}$ codewords. Since $\ker(I+A)$ is itself a storage code, the upper bound shows that
it has at most $2^{10}$ words, so it equals the code; its dimension is $10$, and the rank of $I+A$ is $16-10=6$. The
identification with the folded $5$-cube is checked by decision procedure.
\end{outline}

\begin{proposition}[The Shannon bound of the Clebsch graph]\label{prop:clebsch-exact}
\formal{ClebschCertificate.lean}{clebsch_rateBound_iff}
\formal{Polymatroid.lean}{RateBound.mono}
The Clebsch graph satisfies the rate bound $c$ if and only if $c\ge\tfrac58$.
\end{proposition}

\begin{outline}
The \emph{if} direction is \cref{cor:H4} with monotonicity in $c$. For the \emph{only if} direction, the entropic
polymatroid of the uniform distribution on the binary code has $h(V)=10\log2$ and $h(\{v\})\le\log2$, which violates
every rate bound below~$\tfrac58$.
\end{outline}

Cameron, Dang and Riis proved $10\le\gn(\text{Clebsch})\le11$~\cite{CameronDangRiis2016} and, in the preprint
version, asked for the exact guessing numbers of the strongly regular graphs they consider~\cite[Problem~5.11]{CameronDangRiis2014};
\cref{cor:H4} answers this for the Clebsch graph. The binary code $\ker(I+A)$ goes back to Blasiak,
Kleinberg and Lubetzky~\cite{BKL2013}.

\paragraph{Why it matters.} The Clebsch graph is triangle-free with odd girth~$5$ and degree~$5$, and its codes beat the
pairing by a constant factor. So bounded degree alone does not force rate~$\tfrac12$: any proof of the undirected crux must
use local bipartiteness (few short odd cycles), not only sparsity. It also rules out
targets: it is not good, although every subcubic graph of girth at least~$17$ maps to it~\cite{DeVosSamal2006}, and
by Theorem~H no graph that receives it is good, which makes iterated Mycielski graphs bad (\cref{prop:M}). On the
construction side it shows what a coding gain can look like
at small girth: the binary code is the reduction modulo~$2$ of a rational eigenvalue code with coefficients~$1$
(\cref{rem:ana-clebsch}), a mechanism that Theorem~A and \cref{cor:torsion} rule out on locally bipartite families of
bounded degree.

\subsection{Theorem S: linear codes on subcubic triangle-free graphs}\label{sec:subcubic}

Let $F$ be a field and $M\in F^{V\times V}$ a matrix with unit diagonal. The \emph{recovery set} of $v$ is
$S_v=\{u\ne v: M_{vu}\ne0\}$. The matrix has \emph{no recovery triangle} if $q\notin S_p$ whenever $p,q\in S_r$; this
holds whenever the off-diagonal support of $M$ lies in the edges of a triangle-free graph. A vertex satisfies
\emph{reciprocity}~$(\star\star)$ if $\card{S_v}\le2$, or $\card{S_v}=3$ and $r\in S_v$ whenever $v\in S_r$.

\begin{theorem}[Theorem S]\label{thm:S}
\formal{Subcubic.lean}{Star, RecTriangleFree, ThmSHyp, finrank_subCode_le_half, thmSHyp_of_cliqueFree, star_of_degree_le_three, finrank_ker_le_half_of_degree_le_three, rank_ge_half_of_degree_le_three}
Let $M\in F^{V\times V}$ have unit diagonal and no recovery triangle. For every vertex set $W$ on which every vertex
satisfies $(\star\star)$, the kernel of the principal submatrix $M[W,W]$ has dimension at most $\lfloor\card W/2\rfloor$.
In particular, if $G$ is a triangle-free graph of maximum degree at most~$3$ and the off-diagonal support of $M$ lies in
the edges of $G$, then $\dim\ker M\le\lfloor\card V/2\rfloor$ and $\rank M\ge\lceil\card V/2\rceil$, over every field.
\end{theorem}

\begin{outline}
Codes on vertex sets are subspaces of $F^V$: the code on $W$ is the set of vectors supported on $W$ that satisfy the rows of
$W$, so deleting vertices is an inclusion of subspaces. The proof is a strong induction over $W$; suppose $W$ is a minimal
counterexample.
\begin{enumerate}[leftmargin=2em]
\item \emph{Start.} Deleting one vertex costs at most one dimension, so $\card W=2m+1$ is odd, the code has dimension
exactly $m+1$, and every singleton deletion is \emph{tight}: it leaves the largest dimension the theorem allows.
\item \emph{Tightness.} On a tight deletion no coordinate vanishes identically, since otherwise a smaller set would be a
counterexample. If two coordinates $p,q$ vanish together on a tight code, deleting both is again tight, and the deleted
code is the subspace $\{y_p=0\}$ of the old one.
\item \emph{Maximal chains.} A \emph{good chain} starts at a vertex $v_0$ with at most two recovery neighbours in $W$, which
exists by a parity argument (otherwise the recovery relation would be a symmetric $3$-regular graph on an odd set). It
repeatedly adds a pair $\{p,q\}$ that is exactly the part of a chain vertex's recovery set outside the chain. Along a good
chain, every deletion is tight and its code lies inside the code of~$W$, so every row of a chain vertex restricted to
the rest is a relation on the remaining code. If some chain vertex has exactly one recovery neighbour outside the chain,
that coordinate vanishes, contradicting tightness. If one has exactly two, the chain extends. Reciprocity rules out
three, because the vertex that added a chain vertex lies in its recovery set.
\item \emph{Final count.} A maximal chain therefore leaves every chain vertex with no recovery neighbours outside it. Call
\emph{free} the vertex $v_0$ and the first vertex of each pair. A codeword that vanishes on the free vertices vanishes on
the whole chain, because the row of each pair's parent then forces the second vertex of the pair to vanish. For the last
pair $\{p,q\}$ even $p$ is forced: its row involves only earlier chain vertices, since $q\notin S_p$ ($p,q\in S_r$ for the
parent $r$, and $M$ has no recovery triangle). So a codeword is determined, on the chain, by $k$ coordinates for a chain
with $k$ pairs, and $m+1\le k+\lfloor(\card W-2k-1)/2\rfloor=m$, a contradiction.
\end{enumerate}
\end{outline}

The proof needs neither Gallai's lemma nor connectivity, and the first part of \cref{thm:S}, stated for recovery sets, is slightly stronger
than the graph-theoretic statement. The triangle condition is only needed on the recovery sets, and reciprocity only on the vertex
set being bounded. Over a field, linear storage codes are exactly the subspaces of the kernels of such matrices:

\begin{lemma}[Linear storage codes are subspaces of kernels]\label{lem:linear-pattern}
\formal{Subcubic.lean}{IsLinearStorageCode, IsLinearStorageCode.exists_coeffs, IsLinearStorageCode.exists_pattern, isLinearStorageCode_of_le_ker, isStorageCode_iff_isLinearStorageCode}
Call a subspace $C\subseteq F^V$ a \emph{linear storage code} on $G$ if every $x\in C$ that vanishes on the neighbours
of~$v$ vanishes at~$v$; over a finite field this is equivalent, by linearity, to the defining property of
\cref{def:storage-code}. If $C$ is a linear storage code, each coordinate is a fixed linear combination of the
neighbouring coordinates on~$C$, so $C\subseteq\ker M$ for a matrix $M$ with unit diagonal and off-diagonal support in the
edges of~$G$. Conversely, every subspace of such a kernel is a linear storage code.
\end{lemma}

\begin{outline}
On $C$, the evaluation at $v$ vanishes on the common kernel of the evaluations at the neighbours of $v$, so it lies in
their span (Mathlib's \texttt{mem\_span\_of\_iInf\_ker\_le\_ker}). Conversely, on $\ker M$ the row of~$v$ gives
$x_v=-\sum_{u\sim v}M_{vu}x_u$. For the equivalence, $x,y\in C$ agree on a set exactly when $x-y\in C$ vanishes there.
\end{outline}

\begin{corollary}[Theorem S for linear storage codes]\label{cor:S-linear}
\formal{Subcubic.lean}{IsLinearStorageCode.finrank_le_half, IsLinearStorageCode.card_sq_le}
On a triangle-free graph of maximum degree at most~$3$, every linear storage code over any field has dimension at most
$\lfloor\card V/2\rfloor$; over a finite field $F$ it has $\card C^2\le\card F^{\card V}$.
\end{corollary}

\begin{outline}
By \cref{lem:linear-pattern}, $C\subseteq\ker M$ for a matrix $M$ as in \cref{thm:S}; apply \cref{thm:S}. Over a finite
field, $\card C=\card F^{\dim C}$.
\end{outline}

The exact minrank follows by splitting along separated vertex sets.

\begin{theorem}[Theorem S, barrier and matching forms]\label{thm:S-barrier}
\partformal{Subcubic.lean}{finrank_subCode_union_le, finrank_subCode_biUnion_le, finrank_ker_le_barrier, finrank_ker_le_heavy, exists_linearStorageCode_of_matching, exists_pattern_finrank_ker_ge_of_matching}{the identification of the minimum of the bound with the matching number, which is the Tutte--Berge formula~\cite{Berge1958}.}
Let $M$ have unit diagonal, no recovery triangle, and off-diagonal support in the edges of a graph~$G$. Let $U\subseteq V$ be
a set outside which every vertex satisfies $(\star\star)$, and let the vertices outside $U$ be partitioned into blocks with
no edges between different blocks (for instance the components of $G-U$). Then
\[
  \dim\ker M\;\le\;\card U+\sum_{B}\lfloor\card B/2\rfloor .
\]
Taking for $U$ the vertices violating $(\star\star)$ gives $\dim\ker M\le\card U+\lfloor(\card V-\card U)/2\rfloor$. By the
Tutte--Berge formula, the minimum of the right side over all $U$, with the components of $G-U$ as blocks, is the matching
number $\nu(G)$. Conversely, over every field and for every graph, a matching of $m$ edges gives such a matrix with
$\dim\ker M\ge m$. Hence, for triangle-free $G$ of maximum degree at most~$3$, $\minrk_F(G)=\card V-\nu(G)$ over every
field.
\end{theorem}

\begin{outline}
Deleting $U$ costs at most $\card U$ dimensions. A code on a union of blocks with no edges between them splits as a sum
of the codes on the blocks, because the rows of one block vanish on the others; each block's code obeys \cref{thm:S}. The
matching code puts one free symbol on each edge of the matching, constant along the edge, and zero elsewhere; it is a
linear storage code of dimension~$m$, and \cref{lem:linear-pattern} turns it into a matrix with $\dim\ker M\ge m$. The
Tutte--Berge formula is classical and not formalized.
\end{outline}

The statement appears to be new: we found no statement that the minrank of triangle-free subcubic graphs equals the
clique-cover number $\card V-\nu(G)$ over every field, in the minrank, index-coding, storage-code or guessing-game
literature. Special cases are known. Tims~\cite{Tims2013} proves that the minimum rank over
every field equals the clique-cover number for all graphs of order at most~$10$ (by computer), the Petersen graph
included, and for the Möbius ladders, cubic triangle-free graphs for which Blasiak, Kleinberg and
Lubetzky~\cite{BKL2013} had shown broadcast rate $\card V/2$ for all codes; Berliner and
Langberg~\cite{BerlinerLangberg2011} prove the same equality over every finite field for outerplanar graphs. None of the hypotheses can be dropped.
The Clebsch graph (degree~$5$, \cref{cor:H4}) has binary minrank $6<8=\card V-\nu$. The complete graph $K_4$ shows that
triangles must be excluded. On $C_5$, vector-linear codes with two symbols per vertex reach dimension~$5$, more than
$2\nu(C_5)=4$, so the theorem is genuinely scalar (although rate~$\tfrac12$ still holds for $C_5$ by \cref{prop:C5}).

\paragraph{Why it matters.} Theorem~S settles, for scalar linear codes, the first whole class of bounded
degree beyond degree~$2$, where earlier results covered particular families, and it is exact: in degree~$3$ no scalar
linear code over any field, characteristic~$2$ included, beats the pairing on any triangle-free graph, locally
tree-like or not. With Theorem~A it narrows where a coding gain could live: for scalar linear codes, at degree at
least~$4$ over a finite field through $p$-torsion; for vector-linear codes over finite fields and for nonlinear codes,
already at degree~$3$.

% Continuation of the codes section (\cref{sec:codes}): subsections only.

\subsection{Bounded complex coefficients: Theorem A}\label{sec:ana-A}

The rate bounds proved so far in this section hold either for all codes on special graphs (\cref{thm:H} and its
corollaries) or for all linear codes on triangle-free graphs of maximum degree~$3$ (\cref{thm:S}). The undirected crux
(\cref{prob:undirected-crux}) concerns locally bipartite side information of bounded degree. In this subsection we
prove the bound~$\frac12+o(1)$ on the rate in that setting for one class of codes: linear codes over~$\R$ or~$\C$
whose recovery coefficients are bounded.

Let $G$ be a finite graph with vertex set $V$ and $N=\card V$. A matrix $M\in F^{V\times V}$ over a field $F$
\emph{fits} $G$ if $M_{vv}=1$ for every $v$ and $M_{vu}=0$ whenever $u\ne v$ and $uv\notin E(G)$. By
\cref{lem:linear-pattern}, the linear storage codes on $G$ over $F$ are the subspaces of kernels of matrices that fit
$G$; the coefficients $R_{vu}$ of a recovery rule $x_v=\sum_{u\sim v}R_{vu}x_u$ are the negated off-diagonal entries of
such a matrix, which we call a \emph{check matrix} of the code.
A complex matrix is \emph{$B$-bounded} if all its entries have modulus at most $B$; a matrix that fits a
graph can be $B$-bounded only if $B\ge1$. For $\ell\ge1$ let $\mathrm{odd}_\ell(G)$ be the set of vertices of $G$ that
lie on an odd cycle of length at most $\ell$. As in \cref{prob:undirected-crux}, a sequence of graphs $G_n$ on
$N_n\to\infty$ vertices is \emph{locally bipartite} if $\card{\mathrm{odd}_\ell(G_n)}=o(N_n)$ for every fixed $\ell$.
Sequences whose girth tends to infinity
are locally bipartite, and so are \emph{locally tree-like} sequences, in which, for every fixed $\ell$, only $o(N_n)$
vertices lie on cycles of length at most $\ell$.

\begin{theorem}[Theorem A: bounded complex coefficients]\label{thm:A}
\notformalized
Let $D\ge1$ and $B\ge1$.
\begin{enumerate}[label=(\roman*)]
\item Let $(G_n)$ be a locally bipartite sequence of graphs of maximum degree at most $D$ on $N_n$ vertices, and let
  $M_n\in\C^{V(G_n)\times V(G_n)}$ be $B$-bounded matrices that fit $G_n$. Then
  $\liminf_n\rank(M_n)/N_n\ge\frac12$.
\item \emph{(Uniform form.)} For every $\eps>0$ there are $\delta>0$ and $N_0$ such that $\rank M\ge(\frac12-\eps)N$
  for every graph $G$ on $N\ge N_0$ vertices of maximum degree at most $D$ with
  $\card{\mathrm{odd}_{\lfloor1/\delta\rfloor}(G)}\le\delta N$ and every $B$-bounded complex matrix $M$ that fits $G$.
\item \emph{(Random regular graphs.)} Let $d\ge3$ and $\eps>0$, and let $G_{n,d}$ be a uniformly random $d$-regular
  graph on $n$ vertices ($nd$ even). With probability tending to $1$ as $n\to\infty$, every $B$-bounded complex matrix
  that fits $G_{n,d}$ has rank at least $(\frac12-\eps)n$.
\end{enumerate}
\end{theorem}

In terms of codes (\cref{cor:ana-codes}): on a locally bipartite sequence of graphs of bounded degree, a linear storage
code over $\R$ or $\C$ whose recovery coefficients have modulus at most $B$ has dimension at most $(\frac12+o(1))N_n$.

\paragraph{Why it matters.} For the impossibility side, Theorem~A proves the analogue of the crux's rate bound
(\cref{def:crux}) for real and complex linear codes with bounded recovery coefficients, scalar or with a fixed number of
symbols per vertex (\cref{prop:ana-A-ext}), on side information of bounded total degree whose underlying graphs are
locally bipartite. The crux itself concerns binary codes; Theorem~A bears on them through \cref{cor:torsion}. Only the support of the check matrix enters, so directed side information is
covered through its underlying graph, and locally tree-like side information is covered. For the construction side it
is a design rule: on such side information, a code that beats rate~$\frac12$ by a constant is nonlinear, or vector-linear
with a number of symbols per vertex that grows with $N$, or linear with coefficients that no diagonal rescaling makes
bounded (\cref{prop:ana-A-ext}(a), for $t=1$), or linear over a finite field, and then every bounded integer lift of its
check matrix has $p$-torsion of linear rank (\cref{cor:torsion}). Codes over finite alphabets, linear or not, are not
covered; they are the open part of the crux.

\paragraph{Local limits.} We use the framework of Benjamini and Schramm~\cite{BenjaminiSchramm2001} and of Aldous and
Lyons~\cite{AldousLyons2007}. A \emph{network} is a connected graph of maximum degree at most $D$ with \emph{marks}: a
point of a compact metric space $\Xi_V$ at every vertex and a point of a compact metric space $\Xi_E$ at each end of
every edge. A \emph{rooted} network has a distinguished vertex $o$, a \emph{doubly rooted} network an ordered pair of
distinguished vertices, and we identify networks that are isomorphic by graph isomorphisms preserving the roots and the
marks. Two rooted networks are close if, for a large $R$, the balls of radius $R$ around their roots are isomorphic by a
root-preserving graph isomorphism that changes every mark by less than $1/R$. Since degrees are bounded and the mark
spaces are compact, the space of rooted networks is compact (a diagonal argument). A finite marked graph of maximum
degree at most $D$, connected or not, defines a random rooted network: choose $o$ uniformly among its vertices and keep
the component of $o$. A sequence of finite marked graphs \emph{converges locally} to a probability measure $\mu$ on
rooted networks if these random rooted networks converge to $\mu$ in distribution. A probability measure $\mu$ on
rooted networks is \emph{unimodular} if it satisfies the mass-transport principle
\begin{equation}\label{eq:ana-mtp}
  \mathbb{E}_\mu\sum_{x\in V(G)}f(G,o,x)\;=\;\mathbb{E}_\mu\sum_{x\in V(G)}f(G,x,o)
\end{equation}
for every Borel function $f\ge0$ on doubly rooted networks, where $(G,o)$ denotes a $\mu$-random rooted network. We use
four standard facts.
\begin{enumerate}[label=(L\arabic*)]
\item Every sequence of finite marked graphs as above has a locally convergent subsequence (compactness and
  Prokhorov's theorem).
\item Every local limit of finite marked graphs is unimodular~\cite[Section~2]{AldousLyons2007}. (A finite graph with
  several components is a convex combination of its components, each rooted uniformly.)
\item Let $\mu$ be unimodular. A family $T=(T_G)$ of bounded operators $T_G$ on $\ell^2(V(G))$, indexed by networks,
  is \emph{equivariant} if $\langle T_{G'}\delta_{\varphi x},\delta_{\varphi y}\rangle=\langle T_G\delta_x,\delta_y\rangle$
  for every isomorphism $\varphi\colon G\to G'$ of networks and all vertices $x,y$, and \emph{measurable} if
  $(G,x,y)\mapsto\langle T_G\delta_x,\delta_y\rangle$ is Borel on doubly rooted networks. The $\mu$-classes of
  equivariant, measurable families with $\sup_G\lVert T_G\rVert<\infty$ form a von Neumann algebra $\mathrm{Alg}(\mu)$,
  acting on the direct integral of the spaces $\ell^2(V(G))$, on which
  $\mathrm{Tr}(T)=\mathbb{E}_\mu\langle T_G\delta_o,\delta_o\rangle$ is a positive linear functional with
  $\mathrm{Tr}(ST)=\mathrm{Tr}(TS)$~\cite[Section~5]{AldousLyons2007}. Its elements act fibrewise, and their kernels,
  ranges and spectral projections are computed fibrewise.
\item (Polar decomposition.) For $T$ in a von Neumann algebra $\mathcal A$ there is a partial isometry $U\in\mathcal A$
  such that $U^*U$ is the projection onto $(\ker T)^\perp$ and $UU^*$ is the projection onto the closure of the range
  of $T$~\cite[Theorem~6.1.2 and Proposition~6.1.3]{KadisonRingrose1997}.
\end{enumerate}

\begin{informalproof}
(i) Let $\nu_n=1-\rank(M_n)/N_n$ be the normalized nullity and $\nu=\limsup_n\nu_n$; we show $\nu\le\frac12$. Make
$G_n$ a marked graph by putting the mark $M_n[v,u]$, a point of the disc $\{z\in\C:\abs z\le B\}$, at the end $v$ of
each edge $vu$. Pass to a subsequence along which $\nu_n\to\nu$ and, by (L1), the marked graphs converge locally to some
$\mu$, which is unimodular by (L2). Write $(G,M,o)$ for a $\mu$-random rooted network. The marks define an operator $M_G$
on $\ell^2(V(G))$ by $(M_Gx)_v=x_v+\sum_{u\sim v}M[v,u]\,x_u$. Each row and each column of $M_G$ has at most $D+1$
nonzero entries, of modulus at most $B$, so $\lVert M_G\rVert\le(D+1)B$, and $(M_G)$ is an equivariant measurable
family.

\emph{Step 1: the limit is bipartite.} For $\ell,R\ge1$ let $A_{\ell,R}$ be the set of rooted networks in which some odd
cycle of length at most $\ell$ meets the ball of radius $R$ around the root. Such a cycle lies in the ball of radius
$R+\ell$, so membership in $A_{\ell,R}$ is decided by that ball without its marks, and $A_{\ell,R}$ is open and closed.
A uniform root of $G_n$ lies in $A_{\ell,R}$ only if it is within distance $R$ of $\mathrm{odd}_\ell(G_n)$, which has
probability at most $(1+D+\dots+D^R)\card{\mathrm{odd}_\ell(G_n)}/N_n\to0$. Since $A_{\ell,R}$ is open and closed,
$\mu(A_{\ell,R})$ is the limit of these probabilities, so $\mu(A_{\ell,R})=0$. Every odd cycle of $G$ meets some ball
around $o$, so $\mu$-almost surely $G$ has no odd cycle: $G$ is bipartite.

\emph{Step 2: semicontinuity of the nullity.} Let $H_n=M_n^*M_n$ and let $\sigma_n$ be the empirical eigenvalue
distribution of $H_n$ (eigenvalues counted with multiplicity, each with mass $1/N_n$); it is supported in $[0,c]$ with
$c=((D+1)B)^2$. For $k\ge0$, $\int x^k\,d\sigma_n=N_n^{-1}\operatorname{tr}H_n^k$ is the expectation, over a uniform
vertex $o$ of $G_n$, of $\langle H_n^k\delta_o,\delta_o\rangle$. Expanding the product,
$\langle(M^*M)^k\delta_o,\delta_o\rangle$ is a sum of at most $(D+1)^{2k}$ products of $2k$ entries of $M$ or their
conjugates, taken along walks of length $2k$ from $o$ to $o$ whose steps stay put or move to a neighbour; such walks
stay in the ball of radius $k$. Hence $(G,M,o)\mapsto\langle(M_G^*M_G)^k\delta_o,\delta_o\rangle$ is a bounded continuous
function on rooted networks, and $\int x^k\,d\sigma_n\to\mathbb{E}_\mu\langle H^k\delta_o,\delta_o\rangle=\int
x^k\,d\sigma$, where $H=M_G^*M_G$ and $\sigma=\mathbb{E}_\mu\sigma_{H,o}$ is the average of the spectral measures
$\sigma_{H,o}$ of $H$ at $\delta_o$. All these are probability measures on $[0,c]$, so convergence of moments implies
$\sigma_n\to\sigma$ weakly (Weierstrass approximation), and the portmanteau theorem for the closed set $\{0\}$ gives
$\limsup_n\sigma_n(\{0\})\le\sigma(\{0\})$. Since $\ker(M^*M)=\ker M$ for matrices and for bounded operators,
$\sigma_n(\{0\})=\nu_n$ and $\sigma(\{0\})=\mathbb{E}_\mu\langle P_K\delta_o,\delta_o\rangle$, where $P_K$ is the
orthogonal projection onto $K=\ker M_G$. Therefore
\begin{equation}\label{eq:ana-step2}
  \nu\;\le\;\mathbb{E}_\mu\langle P_K\delta_o,\delta_o\rangle .
\end{equation}
This is the easy half of L\"uck's approximation theorem~\cite{Luck1994}; for adjacency matrices of random graphs the
same argument is used by Bordenave, Lelarge and Salez~\cite{BLS2011}. The bound $B$ is used here and in~(L1).

\emph{Step 3: a random colouring.} For $\mu$-almost every $(G,M,o)$ the graph $G$ is connected and bipartite, so it has
exactly two proper colourings with colours $\pm1$, each the negative of the other. Let $\hat\mu$ be the law of
$(G,M,\beta,o)$, where $(G,M,o)$ has law $\mu$ and $\beta$, a vertex mark, is one of the two colourings, chosen by a fair
coin independent of everything else. Then $\hat\mu$ is unimodular. Indeed, let $f\ge0$ be Borel on doubly rooted
coloured networks. For a doubly rooted network $(G,M,x,y)$ with $G$ bipartite let $\pi_x(w)=(-1)^{\dist(x,w)}$, and put
$\bar f(G,M,x,y)=\frac12\bigl(f(G,M,\pi_x,x,y)+f(G,M,-\pi_x,x,y)\bigr)$ (and $\bar f=0$ if $G$ is not bipartite).
Since $\{\pi_x,-\pi_x\}=\{\pi_y,-\pi_y\}$ is the set of both colourings, $\bar f$ is the average of $f$ over the two
colourings whichever root is used to define it; it is Borel and invariant under isomorphisms. Hence
\[
  \mathbb{E}_{\hat\mu}\sum_xf(G,M,\beta,o,x)=\mathbb{E}_\mu\sum_x\bar f(G,M,o,x)
  =\mathbb{E}_\mu\sum_x\bar f(G,M,x,o)=\mathbb{E}_{\hat\mu}\sum_xf(G,M,\beta,x,o),
\]
the middle equality being \eqref{eq:ana-mtp} for $\mu$. Moreover $\hat\mu(\beta(o)=-1)=\frac12$. (Without the coin,
colouring by the parity of the distance to the root is not unimodular in general: the root would always have colour
$+1$ and its neighbours colour $-1$.)

\emph{Step 4: the trace.} By (L3) the algebra $\mathrm{Alg}(\hat\mu)$ contains $M=(M_G)$, hence $H=M^*M$ and its spectral
projection $P_K=\mathbf 1_{\{0\}}(H)$, and the projection $P_J$ onto $\ell^2(J)$, where $J=\beta^{-1}(-1)$; let
$V_+=\beta^{-1}(+1)$. Let $T=P_JP_K\in\mathrm{Alg}(\hat\mu)$. If $Tx=0$, write $x=k+k'$ with $k\in K$ and
$k'\in K^\perp$; then $P_Jk=Tx=0$, so $k$ vanishes on $J$. Every neighbour of a vertex $v\in V_+$ lies in $J$, so
$0=(M_Gk)_v=k_v+\sum_{u\sim v}M[v,u]\,k_u=k_v$. Hence $k=0$, and $\ker T=K^\perp$. The range of $T$ lies in
$\ell^2(J)$. By (L4) there is a partial isometry $U\in\mathrm{Alg}(\hat\mu)$ with $U^*U=P_K$ and $UU^*\le P_J$, and by
(L3)
\[
  \mathbb{E}_\mu\langle P_K\delta_o,\delta_o\rangle=\mathrm{Tr}(P_K)=\mathrm{Tr}(U^*U)=\mathrm{Tr}(UU^*)
  \le\mathrm{Tr}(P_J)=\hat\mu(\beta(o)=-1)=\tfrac12 .
\]
The first equality holds because $\langle P_K\delta_o,\delta_o\rangle$ does not depend on $\beta$ and the law of
$(G,M,o)$ under $\hat\mu$ is $\mu$. With \eqref{eq:ana-step2}, $\nu\le\frac12$.

(ii) If (ii) failed for some $\eps$, then for every $m$ there would be a graph $G_m$ on $N_m\ge m$ vertices of maximum
degree at most $D$ with $\card{\mathrm{odd}_m(G_m)}\le N_m/m$, and a $B$-bounded matrix $M_m$ fitting $G_m$ with
$\rank M_m<(\frac12-\eps)N_m$. For fixed $\ell$ and $m\ge\ell$ we have
$\card{\mathrm{odd}_\ell(G_m)}\le\card{\mathrm{odd}_m(G_m)}\le N_m/m$, so $(G_m)$ is locally bipartite, and (i)
gives a contradiction.

(iii) By a theorem of Wormald~\cite{Wormald1981}, for each fixed $\ell\ge3$ the number $X_\ell$ of cycles of length
$\ell$ in $G_{n,d}$ converges in distribution as $n\to\infty$, to a Poisson law with mean $(d-1)^\ell/(2\ell)$; in
particular $X_\ell$ is bounded in probability. Let $\delta$ and $N_0$ be given by (ii) for $D=d$. Then
$\card{\mathrm{odd}_{\lfloor1/\delta\rfloor}(G_{n,d})}\le\sum_{\ell\le1/\delta}\ell X_\ell$ is bounded in probability,
so it is at most $\delta n$ with probability tending to $1$, and (ii) applies.
\end{informalproof}

\begin{proposition}[Extensions of Theorem A]\label{prop:ana-A-ext}
\notformalized
Let $D\ge1$ and $B\ge1$.
\begin{enumerate}[label=(\alph*)]
\item \emph{(Tight coefficients and gauge.)} For every $\eps>0$ there are $\delta'>0$ and $N_0'$ with the following
  property. Let $G$ be a graph on $N\ge N_0'$ vertices of maximum degree at most $D$ with
  $\card{\mathrm{odd}_{\lfloor1/\delta'\rfloor}(G)}\le\delta'N$, let $M\in\C^{V\times V}$ fit $G$, let $S$ be an
  invertible diagonal matrix, and let $X$ be the set of rows of $S^{-1}MS$ that contain an entry of modulus larger than
  $B$. Then $\rank M\ge(\frac12-\eps)(N-\card X)$.
\item \emph{(A fixed number of symbols per vertex.)} Let $t\ge1$. Parts (i) and (ii) of \cref{thm:A} hold for complex
  matrices $M$ with rows and columns indexed by $V\times[t]$ and entries of modulus at most $B$ whose $t\times t$ blocks
  satisfy $M_{vv}=I_t$ and $M_{vu}=0$ whenever $u\ne v$ and $uv\notin E(G)$, with $N$ replaced by $tN$ in the
  conclusions; the constants $\delta$ and $N_0$ depend on $t$.
\item \emph{(Invariant independent sets.)} Let $B$-bounded matrices $M_n$ fitting graphs $G_n$ of maximum degree at most
  $D$, viewed as marked graphs as in the proof of \cref{thm:A}, converge locally to $\mu$. Let $\alpha_U(\mu)$ be the
  supremum of $\hat\mu(\iota(o)=1)$ over all unimodular laws $\hat\mu$ of networks with an additional vertex mark
  $\iota\in\{0,1\}$ such that forgetting $\iota$ gives $\mu$ and $\iota^{-1}(1)$ is $\hat\mu$-almost surely an
  independent set. Then $\liminf_n\rank(M_n)/N_n\ge\alpha_U(\mu)$.
\item \emph{(Sharpness.)} If $G$ has a perfect matching $P$ with adjacency matrix $A_P$, then $I+A_P$ is $1$-bounded,
  fits $G$, and has rank $N/2$. In particular the constant $\frac12$ of \cref{thm:A} is attained on every regular
  bipartite graph of positive degree.
\end{enumerate}
\end{proposition}

\begin{informalproof}
(a) The matrix $S^{-1}MS$ fits $G$ and has the rank of $M$, so we may assume $S=I$. Let $\delta$ and $N_0$ be given by
\cref{thm:A}(ii) for $D$, $B$ and $\eps$, and put $\delta'=2\delta/(D+1)$ and $N_0'=\lceil(D+1)N_0/2\rceil$. Let
$N'=N-\card X$. If $N'\le2N/(D+1)$, let $Y$ be an independent set of $G$ with $\card Y\ge N/(D+1)$ (a greedy choice);
the principal submatrix $M[Y,Y]$ is the identity, so $\rank M\ge\card Y\ge N'/2$. Otherwise $N'>2N/(D+1)\ge N_0$, and
the principal submatrix $M'=M[V\setminus X,V\setminus X]$ is $B$-bounded and fits $G'=G[V\setminus X]$. Every odd cycle
of $G'$ is an odd cycle of $G$, and $\lfloor1/\delta\rfloor\le\lfloor1/\delta'\rfloor$, so
$\card{\mathrm{odd}_{\lfloor1/\delta\rfloor}(G')}\le\card{\mathrm{odd}_{\lfloor1/\delta'\rfloor}(G)}\le\delta'N<\delta N'$.
\Cref{thm:A}(ii) gives $\rank M\ge\rank M'\ge(\frac12-\eps)N'$.

(b) Let $G[t]$ be the graph on $V\times[t]$ in which $(v,i)\sim(u,j)$ if and only if $uv\in E(G)$. Such a matrix $M$
fits $G[t]$, whose maximum degree is $tD$. An odd cycle of length $\ell'\le\ell$ through $(v,i)$ in $G[t]$ projects to a
closed walk of odd length $\ell'$ through $v$ in $G$; a closed walk of odd length contains an odd cycle of length at
most $\ell'$ formed by vertices of the walk, all within distance $\ell'/2\le\ell$ of $v$. Hence
$\mathrm{odd}_\ell(G[t])\subseteq W\times[t]$, where $W$ is the set of vertices within distance $\ell$ of
$\mathrm{odd}_\ell(G)$, and $\card{\mathrm{odd}_\ell(G[t])}\le t(1+D+\dots+D^\ell)\card{\mathrm{odd}_\ell(G)}$. So
$(G_n[t])$ is locally bipartite when $(G_n)$ is, and \cref{thm:A}(i) applies to $G_n[t]$ with $D$ replaced by $tD$.
For~(ii), take $\delta_t$ and $N_0$ from \cref{thm:A}(ii) for degree $tD$, put $\ell=\lfloor1/\delta_t\rfloor$ and
$\delta=\min\bigl(\delta_t/(1+D+\dots+D^\ell),1/\ell\bigr)$. Then $\card{\mathrm{odd}_{\lfloor1/\delta\rfloor}(G)}\le\delta N$
gives $\card{\mathrm{odd}_\ell(G[t])}\le t(1+D+\dots+D^\ell)\delta N\le\delta_t\,tN$, and $tN\ge N\ge N_0$.

(c) Steps~2 and~4 of the proof of \cref{thm:A} do not use bipartiteness. Step~2, applied to the whole sequence, gives
$\limsup_n(1-\rank(M_n)/N_n)\le\mathbb{E}_\mu\langle P_K\delta_o,\delta_o\rangle$. Step~4, run in
$\mathrm{Alg}(\hat\mu)$ with $\iota^{-1}(1)$ in place of $V_+$ and $J=\iota^{-1}(0)$ (every neighbour of a vertex of
$\iota^{-1}(1)$ lies in $J$, since $\iota^{-1}(1)$ is independent), gives $\mathbb{E}_\mu\langle P_K\delta_o,\delta_o\rangle\le\hat\mu(\iota(o)=0)
=1-\hat\mu(\iota(o)=1)$. Take the supremum over $\hat\mu$.

(d) Ordering the vertices in matched pairs, $I+A_P$ is block diagonal with one block
$\bigl(\begin{smallmatrix}1&1\\1&1\end{smallmatrix}\bigr)$ of rank $1$ per edge of $P$. A regular bipartite graph with
at least one edge has a perfect matching, by Hall's theorem.
\end{informalproof}

\begin{remark}[On the extensions]\label{rem:ana-A-ext}
Part~(a) says that boundedness is needed only up to a diagonal rescaling and outside $o(N)$ rows. A rescaling changes
$M_{vu}$ into $s_v^{-1}M_{vu}s_u$, so it preserves the products of entries around closed walks, for instance
$M_{vu}M_{uv}$; ``bounded in some gauge'' is a condition on these products. In part~(b) the constants depend on $t$
through the degree $tD$ of $G[t]$, so the argument says nothing about the linear capacity, a supremum over~$t$, or
about codes whose number of symbols per vertex grows with $N$. Part~(c) explains why the limit object is needed. A
finite independent set $Y$ gives only $\rank M\ge\card Y$, since $M[Y,Y]$ is the identity, and graphs of large girth
can have small independence ratio: there are cubic graphs of arbitrarily large girth with independence ratio less than
$6/13$~\cite{Bollobas1981}. In a limit concentrated on the $d$-regular tree, by contrast, $\alpha_U=\frac12$: the
coin-flip colouring of Step~3 gives $\alpha_U\ge\frac12$, and the mass transport that sends one unit from each vertex
of $\iota^{-1}(1)$ to each of its $d$ neighbours gives $d\,\hat\mu(\iota(o)=1)\le d\,\hat\mu(\iota(o)=0)$. By part~(d)
the bound is attained along locally tree-like sequences as well, for instance on bipartite double covers of $d$-regular
graphs of large girth.
\end{remark}

\begin{remark}[Literature]\label{rem:ana-A-lit}
We have not found Theorem~A in the literature; its proof combines the easy half of
L\"uck approximation~\cite{Luck1994,BLS2011} with a tracial form of the Haemers bound~\cite{Haemers1979},
$\rank M\ge\card Y$ for an independent set $Y$, applied to an invariant random independent set of the limit; a tracial
version of the Haemers bound for finite graphs was studied by Gao, Gribling and Li~\cite{GaoGriblingLi2022}. The
positive semidefinite analogue is classical and exact: a positive semidefinite matrix with unit diagonal whose
off-diagonal support lies in the edges of a triangle-free graph has rank at least $N/2$ (Rosenfeld~\cite{Rosenfeld1991}
for real matrices, and Deaett~\cite{Deaett2011}). It needs no degree or coefficient bound, but it does not imply
Theorem~A, which concerns matrices that need not be symmetric or positive semidefinite. For such matrices
triangle-freeness alone does not suffice without a degree bound: Codenotti, Pudl\'ak and Resta~\cite{CPR2000} construct
triangle-free graphs for which $A-2I$, a matrix that fits the graph after scaling by $-\frac12$, has rank
$O(N^{3/4})$.
\end{remark}

\subsection{Codes over fields: torsion and odd characteristic}\label{sec:ana-fields}

We record what Theorem~A says about codes, and what it implies for codes over finite fields.

\begin{corollary}[Bounded real and complex linear codes]\label{cor:ana-codes}
\notformalized
Let $F\in\{\R,\C\}$, let $D,t\ge1$ and $B\ge1$ be fixed, and let $(G_n)$ be a locally bipartite sequence of graphs of
maximum degree at most $D$ on $N_n$ vertices. For $x\in F^{V\times[t]}$ write $x_v\in F^t$ for its block at $v$.
\begin{enumerate}[label=(\roman*)]
\item Let $C_n\subseteq F^{V(G_n)\times[t]}$ be a linear storage code with $t$ symbols per vertex and bounded
  coefficients: there are $t\times t$ matrices $R_{vu}$ ($uv\in E(G_n)$) with entries of modulus at most $B$ such that
  $x_v=\sum_{u\sim v}R_{vu}x_u$ for every $x\in C_n$ and every $v$. Then $\dim C_n\le(\frac12+o(1))\,tN_n$.
\item Let every vertex $v$ of $G_n$ be a receiver that wants $x_v$ and knows the blocks $x_u$ for $u$ in a set
  $I_n(v)\subseteq N_{G_n}(v)$ (directed side information). Let a linear map $E\colon F^{V(G_n)\times[t]}\to F^r$ be a
  linear index code with decoders $x_v=\Lambda_v(Ex)+\sum_{u\in I_n(v)}R_{vu}x_u$, valid for all $x$, where the maps
  $\Lambda_v$ are linear and the matrices $R_{vu}$ have entries of modulus at most $B$. Then $r\ge(\frac12-o(1))\,tN_n$.
\end{enumerate}
\end{corollary}

\begin{informalproof}
In both cases let $M_n$ be the matrix with blocks $(M_n)_{vv}=I_t$, $(M_n)_{vu}=-R_{vu}$ for the pairs $(v,u)$ that carry
a coefficient, and zero blocks elsewhere; it satisfies the hypotheses of \cref{prop:ana-A-ext}(b), so
$\rank M_n\ge(\frac12-o(1))\,tN_n$. (i) $C_n\subseteq\ker M_n$. (ii) For every $x$, the block of $M_nx$ at $v$ is
$\Lambda_v(Ex)$; so $M_n=\Lambda E$, where $\Lambda$ stacks the maps $\Lambda_v$, and $\rank M_n\le r$.
\end{informalproof}

By \cref{prop:ana-A-ext}(a), for $t=1$ the coefficients need to be bounded only after a diagonal rescaling and outside
$o(N_n)$ vertices. The statement does not cover coefficients that no rescaling makes bounded, numbers of symbols per
vertex that grow with $N_n$, codes over finite fields, or nonlinear codes. For finite fields Theorem~A still says
something: a gain must be a torsion effect.

\begin{corollary}[Torsion-only principle]\label{cor:torsion}
\notformalized
Fix a prime $p$, integers $D,t\ge1$, a real $B\ge1$ and $c>0$. Let $(G_n)$ be a locally bipartite sequence of graphs
of maximum degree at most $D$ on $N_n$ vertices. Let $M_n$ be matrices over $\F_p$ with rows and columns indexed by
$V(G_n)\times[t]$ whose $t\times t$ blocks satisfy $(M_n)_{vv}=I_t$ and $(M_n)_{vu}=0$ whenever $u\ne v$ and
$uv\notin E(G_n)$, and suppose that $\dim\ker M_n\ge(\frac12+c)\,tN_n$ for all $n$. Let $\tilde M_n$ be integer
matrices whose diagonal blocks are $I_t$, whose off-diagonal blocks vanish at non-edges and whose entries have modulus at
most $B$, that reduce to $M_n$ modulo $p$ (for instance the lifts with these diagonal blocks and entries in $[-p/2,p/2]$,
with $B=p/2$). Then $p$ divides at least $(c-o(1))\,tN_n$ of the invariant
factors of $\tilde M_n$; equivalently, the cokernel $\Z^{tN_n}/\tilde M_n\Z^{tN_n}$ contains a subgroup isomorphic to
$(\Z/p\Z)^{r_n}$ with $r_n\ge(c-o(1))\,tN_n$.
\end{corollary}

\begin{informalproof}
Write $\tilde M_n=U\Sigma W$ in Smith normal form, with $U,W$ invertible over $\Z$ and
$\Sigma=\mathrm{diag}(s_1,\dots,s_\rho,0,\dots,0)$, where $s_i\ge1$ and $\rho=\rank_\Q\tilde M_n$. Reducing modulo $p$,
where $U$ and $W$ stay invertible, gives $\rank_{\F_p}M_n=\card{\{i\le\rho:p\nmid s_i\}}$, so the number of invariant
factors divisible by $p$ is $\rho-\rank_{\F_p}M_n$. As a complex matrix, $\tilde M_n$ satisfies the hypotheses of
\cref{prop:ana-A-ext}(b), so $\rho\ge(\frac12-o(1))\,tN_n$, while $\rank_{\F_p}M_n=tN_n-\dim\ker M_n\le(\frac12-c)\,tN_n$.
The cokernel is isomorphic to $\Z^{tN_n-\rho}\oplus\bigoplus_i\Z/s_i\Z$, and its elements of order dividing $p$ form a
group $(\Z/p\Z)^{r_n}$ with $r_n=\card{\{i:p\mid s_i\}}$.
\end{informalproof}

An $\F_{p^e}$-linear storage code becomes, after the choice of a basis of $\F_{p^e}$ over $\F_p$, an $\F_p$-linear code
with $t=e$ symbols per vertex and the same rate, because multiplication by an element of $\F_{p^e}$ is $\F_p$-linear; so
\cref{cor:torsion} covers $\F_{p^e}$ for fixed $e$.

\begin{remark}[The Clebsch code is inherited from $\Q$]\label{rem:ana-clebsch}
The Clebsch graph is strongly regular with parameters $(16,5,0,2)$, so its adjacency matrix satisfies
$A^2=5I+2(J-I-A)=3I-2A+2J$, where $J$ is the all-ones matrix. Hence $(I+A)^2\equiv0\pmod 2$, and over $\F_2$ the image
of $I+A$ lies in its kernel. In fact $\rank_{\F_2}(I+A)=6$, which \cref{cor:H4} establishes formally; directly, the Clebsch graph is the Cayley
graph of
$\Z_2^4$ with connection set $\{e_1,\dots,e_4,e_1+\dots+e_4\}$, and in the group algebra
$\F_2[\Z_2^4]=\F_2[y_1,\dots,y_4]/(y_1^2,\dots,y_4^2)$, with $e_i\mapsto1+y_i$, the matrix $I+A$ acts as multiplication
by $a=\sigma_2+\sigma_3+\sigma_4$, where $\sigma_k$ is the $k$-th elementary symmetric polynomial in the $y_i$. The
product of $a$ with a monomial of degree $2$ is $y_1y_2y_3y_4$, and with a monomial of higher degree it is $0$, so the
image is spanned by $a$, $y_1a,\dots,y_4a$ and $y_1y_2y_3y_4$, which are linearly independent. The kernel of $I+A$, of
dimension $16-6=10>8$, is the code behind \cref{cor:H4}, and it is not a torsion effect. Over $\Q$ the eigenvalues of $A$ are $5$, $1$
and $-3$, with multiplicities $1$, $10$ and $5$, so $\rank_\Q(I-A)=6$. The integer kernel of $I-A$ is a direct summand
of $\Z^{16}$ of rank $10$, so its reduction modulo~$2$ is a $10$-dimensional subspace of $\ker_{\F_2}(I-A)=\ker_{\F_2}(I+A)$,
hence all of it. The binary code is therefore the reduction of the rational eigenvalue code of the eigenvalue~$1$
(\cref{prop:E}), whose recovery coefficients are $1$; \cref{cor:torsion} rules out exactly this mechanism on locally
bipartite families of bounded degree. The other lift $I+A$, with eigenvalues $6$, $2$ and $-2$, has full rank and ten
invariant factors divisible by~$2$: the torsion counted in \cref{cor:torsion} depends on the lift, and its conclusion
holds for every bounded lift. A single graph of girth~$4$ is not a locally bipartite family; the Clebsch graph shows
that local bipartiteness cannot be dropped from Theorem~A and \cref{cor:torsion}, even for bounded rational codes.
\end{remark}

The nilpotent identity $(I+W)^2=0$, which the Clebsch check matrix satisfies modulo~$2$, gives no gain once the girth is
at least~$5$.

\begin{lemma}[Index-two nilpotence at girth five]\label{lem:ana-matching}
\notformalized
Let $G$ be a graph of girth at least $5$, and let $W\in F^{V\times V}$, over any field $F$, have zero diagonal and
$W_{vu}=0$ whenever $uv\notin E(G)$. If $W^2=I$, then there is a perfect matching $P$ of $G$ such that $W_{vu}\ne0$
exactly when $uv\in P$, and $W_{vu}W_{uv}=1$ for $uv\in P$; in particular $\rank(I+W)=\rank(I-W)=N/2$. Over a field of
characteristic $2$, $(I+W)^2=I+W^2$, so every matrix $I+W$ that fits $G$ and satisfies $(I+W)^2=0$ has rank exactly
$N/2$.
\end{lemma}

\begin{informalproof}
For $v\ne w$, $(W^2)_{vw}$ is a sum over the common neighbours of $v$ and $w$, of which there is at most one, as $G$
has no cycle of length $3$ or $4$. Since $(W^2)_{vv}=\sum_{u\sim v}W_{vu}W_{uv}=1$, some neighbour $m(v)$ of $v$ has
$W_{v\,m(v)}W_{m(v)\,v}\ne0$. If $W_{vu'}\ne0$ for a neighbour $u'\ne m(v)$, then $m(v)$ and $u'$ are non-adjacent with
unique common neighbour $v$, and $(W^2)_{m(v)\,u'}=W_{m(v)\,v}W_{vu'}\ne0$, contradicting $W^2=I$; so row $v$ of $W$ is
supported on $\{m(v)\}$. Likewise, if $W_{u'v}\ne0$ for a neighbour $u'\ne m(v)$, then
$(W^2)_{u'\,m(v)}=W_{u'v}W_{v\,m(v)}\ne0$; so column $v$ is supported on $\{m(v)\}$. As $W_{m(v)\,v}\ne0$, row $m(v)$ is
supported on $\{v\}$, so $m(m(v))=v$, the pairs $\{v,m(v)\}$ form a perfect matching, and
$W_{v\,m(v)}W_{m(v)\,v}=(W^2)_{vv}=1$. Ordering the vertices in matched pairs, $I\pm W$ is block diagonal with blocks
$\bigl(\begin{smallmatrix}1&\pm a\\\pm a^{-1}&1\end{smallmatrix}\bigr)$ of rank $1$.
\end{informalproof}

In odd characteristic there is a different constraint, which compares a code with its sign twin.

\begin{proposition}[Sign twins in odd characteristic]\label{prop:oddchar}
\notformalized
Let $F$ be a field of characteristic different from $2$ and let $W\in F^{V\times V}$. Then
\[
  \rank(I+W)+\rank(I-W)\;=\;\rank\begin{pmatrix}I&W\\W&I\end{pmatrix}\;\ge\;N .
\]
Hence $\dim\ker(I+W)+\dim\ker(I-W)\le N$, and at most one of $\ker(I+W)$ and $\ker(I-W)$ has dimension larger than
$N/2$. If $W$ has zero diagonal and is supported on the edges of a graph $G$, the block matrix fits the bipartite double
cover $G\times K_2$, the graph on $V\times\{0,1\}$ in which $(v,i)\sim(u,j)$ if and only if $uv\in E(G)$ and $i\ne j$.
\end{proposition}

\begin{informalproof}
Let $Q=\bigl(\begin{smallmatrix}I&I\\I&-I\end{smallmatrix}\bigr)$; then $Q^2=2I$, so $Q$ is invertible. A direct
computation gives $Q\bigl(\begin{smallmatrix}I&W\\W&I\end{smallmatrix}\bigr)Q^{-1}
=\bigl(\begin{smallmatrix}I+W&0\\0&I-W\end{smallmatrix}\bigr)$, which proves the equality. The block matrix contains
$I$ as a principal $N\times N$ submatrix, so its rank is at least $N$. The last claim holds because the off-diagonal
entries of the block matrix are the entries $W_{vu}$ at positions $((v,0),(u,1))$ and $((v,1),(u,0))$.
\end{informalproof}

\begin{remark}[Gauge invariance]\label{rem:ana-gauge}
Over a field of characteristic different from~$2$, if $G$ is bipartite with a proper colouring $s\colon V\to\{\pm1\}$, the diagonal matrix $S=\mathrm{diag}(s)$ satisfies
$SWS=-W$ for $W$ supported on the edges, so $\ker(I+W)$ and $\ker(I-W)$ have the same dimension, at most $N/2$ by
\cref{prop:oddchar}. On a locally bipartite sequence the same conclusion up to $o(N_n)$, for all matrices over a field
of characteristic different from $2$ and uniformly in the number of symbols per vertex, would follow from approximate
gauge invariance: $\dim\ker(I+W_n)-\dim\ker(I-W_n)=o(N_n)$. In the bipartite limit object of the proof of \cref{thm:A}
the coin-flip colouring does conjugate $I+W$ into $I-W$, although that proof uses the colouring differently, through the
independent set $\beta^{-1}(+1)$. In characteristic $2$,
$I+W=I-W$ and \cref{prop:oddchar} says nothing; on the Clebsch graph $I+A=I-A$ has a kernel of dimension $10>8$ over
$\F_2$.
\end{remark}

\paragraph{Why it matters.} For constructions, \cref{cor:torsion} is the finite-field part of the design rule stated
after \cref{thm:A}; in particular, a gain cannot be inherited by reduction modulo~$p$ from a code over $\Q$ with bounded
coefficients, as the Clebsch code is (\cref{rem:ana-clebsch}). In characteristic~$2$, a gain also cannot come from the identity
$(I+W)^2=0$ once the girth is at least~$5$ (\cref{lem:ana-matching}). For the impossibility side,
\cref{prop:oddchar} shows that in odd characteristic a gain on one sign twin is paid for on the other. In
characteristic~$2$ only \cref{cor:torsion} applies.

\subsection{Limits of local and spectral arguments}\label{sec:ana-limits}

Theorem~S settles all scalar linear codes on triangle-free graphs of maximum degree~$3$, and Theorem~A settles bounded
complex linear codes on locally bipartite graphs of every bounded degree by a spectral argument. The undirected crux needs finite
fields (indeed all codes), at degree at least~$4$ for scalar linear codes and already at degree~$3$ for vector-linear
and nonlinear ones. This subsection
records three limits: inside tree-like regions, the local necessary conditions of a smallest counterexample are
compatible with rate above $\frac12+o(1)$ at degree at least~$4$ (\cref{prop:L}); real eigenvalue codes, the mechanism
behind several known codes that beat $\frac12$ at bounded degree, cannot reach rate $\frac12$ at girth~$7$ or more
(\cref{prop:E}); and girth~$5$ does not replace low degree (\cref{prop:HS}).

A class $\mathcal K$ of pairs $(G,M)$, where $M$ fits $G$ over some field, is \emph{hereditary} if
$(G[W],M[W,W])\in\mathcal K$ whenever $(G,M)\in\mathcal K$ and $W\subseteq V(G)$ is nonempty; for example, triangle-free
graphs of maximum degree at most $4$ with all matrices that fit them, over all fields. For $C=\ker M$ and
$U\subseteq V$ let $r(U)=\dim C|_U$ be the dimension of the restriction of $C$ to the coordinates in $U$ (the rank
function of the matroid of $C$), and let $C^\perp=\{y\in F^V:\sum_vy_vx_v=0\text{ for all }x\in C\}$, the row space of
$M$. For a hereditary class, the floor bound $\dim\ker M\le\lfloor N/2\rfloor$ for all
members implies $\dim\ker M\le\nu(G)$, the matching number, by the argument of \cref{thm:S-barrier}: for every vertex
set $U$, $\dim\ker M\le\card U+\sum_H\lfloor\card{H}/2\rfloor$ over the components $H$ of $G-U$, and the Tutte--Berge
formula turns the minimum over $U$ into $\nu(G)$. The floor bound is therefore the statement to prove.

\begin{lemma}[Smallest counterexamples]\label{lem:ana-mincex}
\notformalized
Let $\mathcal K$ be hereditary, and let $(G,M)\in\mathcal K$ satisfy $\dim\ker M>\lfloor N/2\rfloor$ with $N$ as small as
possible. Then
\begin{enumerate}[label=(\alph*)]
\item $N$ is odd and $\dim\ker M=(N+1)/2$;
\item $r(U)\ge\lfloor\card U/2\rfloor+1$ for every nonempty $U\subseteq V$;
\item for every nonempty $X\subseteq V$, the vectors of $C^\perp$ supported in $X$, in particular the rows of $M$
  supported in $X$, span a space of dimension at most $\lceil\card X/2\rceil-1$.
\end{enumerate}
\end{lemma}

\begin{informalproof}
A $1\times1$ matrix that fits a graph is $(1)$, so $N\ge2$. For $U\subseteq V$, the codewords vanishing on $U$ form a
subspace of dimension $\dim C-r(U)$, and their restrictions to $V\setminus U$ lie in $\ker M[V\setminus U,V\setminus U]$
and determine them; so for $U\ne V$ the pair $(G-U,M[V\setminus U,V\setminus U])\in\mathcal K$ has a kernel of dimension
at least $\dim C-r(U)$. (a) If $N$ were even, deleting one vertex ($r\le1$) would leave a kernel of dimension at least
$N/2>\lfloor(N-1)/2\rfloor$ on fewer vertices; if $\dim C\ge(N+3)/2$, deleting one vertex would leave a kernel of
dimension at least $(N+1)/2>\lfloor(N-1)/2\rfloor$. (b) By minimality for $U\ne V$, and trivially for $U=V$,
$\dim C-r(U)\le\lfloor(N-\card U)/2\rfloor$, so $r(U)\ge(N+1)/2-\lfloor(N-\card U)/2\rfloor=\lfloor\card U/2\rfloor+1$,
as $N$ is odd. (c) The vectors of $C^\perp$ supported in $X$ form the annihilator of the
restriction $C|_X$ in $F^X$, of dimension $\card X-r(X)\le\card X-\lfloor\card X/2\rfloor-1$.
\end{informalproof}

One way to reach a contradiction in a smallest counterexample is to exhibit a set $X$ with $\lceil\card X/2\rceil$
independent vectors of $C^\perp$ supported in $X$, for instance rows of $M$. The proof of Theorem~S does this with a
maximal chain grown from a single vertex along recovery neighbourhoods; such a chain closes only through cycles or the
boundary of the vertex set, never inside a tree-like region. The next proposition shows that inside a tree-like region of degree
$d\ge4$ no such set exists, however the region is explored, and that the local conditions there are met by a code of
rate close to $(d-2)/(d-1)$.

\begin{proposition}[Proposition L: tree codes are locally consistent]\label{prop:L}
\notformalized
Let $d\ge3$ and $r\ge1$, let $B$ be the ball of radius $r$ around the root of the $d$-regular tree, $B^\circ\subseteq B$
the set of vertices at distance less than $r$ from the root, and $F$ a field. For each $v\in B^\circ$ let $z_v\in F^B$
have $z_v(v)=1$, $z_v(u)\ne0$ for every neighbour $u$ of $v$, and $z_v(w)=0$ otherwise (over $\F_2$ these are the rows
of $I+A$ indexed by $B^\circ$). Let $C_B=\{x\in F^B:\sum_wz_v(w)x_w=0\text{ for all }v\in B^\circ\}$ and
$r(U)=\dim C_B|_U$.
\begin{enumerate}[label=(\roman*)]
\item $\dim C_B=d(d-1)^{r-1}$, the number of vertices at distance $r$, and $\dim C_B/\card B\to(d-2)/(d-1)$ as
  $r\to\infty$: this is $\frac12$ for $d=3$ and $\frac23$ for $d=4$.
\item $r(U)\ge\card U-\lfloor\card U/(d-1)\rfloor$ for every $U\subseteq B$.
\item If $d\ge4$, then $r(U)\ge\lfloor\card U/2\rfloor+1$ for every nonempty $U\subseteq B$, and for every nonempty
  $X\subseteq B$ the vectors of $C_B^\perp$ supported in $X$, in particular the vectors $z_v$ supported in $X$, span a
  space of dimension at most $\lceil\card X/2\rceil-1$.
\end{enumerate}
\end{proposition}

\begin{informalproof}
Orient the tree away from the root; for $v\in B^\circ$ let $\mathrm{ch}(v)$ be its set of children, of size $d$ at the
root and $d-1$ elsewhere. (i) Order $B^\circ$ by distance from the root. A child of $v$ lies in the support of $z_w$, for
$w\in B^\circ$, only if $w=v$, or $w$ is that child or one of its children, which come after $v$; so each $z_v$ is
nonzero at a coordinate where all earlier vectors vanish, the vectors $z_v$ are linearly independent, and
$\dim C_B=\card B-\card{B^\circ}$ is the number of vertices at distance $r$. Since $\card B=1+d((d-1)^r-1)/(d-2)$, the
ratio tends to $(d-2)/(d-1)$.

(ii) The annihilator of $C_B|_U$ in $F^U$ is $C_B^\perp\cap F^U$, where $C_B^\perp$ is spanned by the vectors $z_v$, so
$r(U)=\card U-\dim(C_B^\perp\cap F^U)$. Let $D_U=\{v\in B^\circ:\mathrm{ch}(v)\subseteq U\}$. If
$y=\sum_v\lambda_vz_v$ is supported in $U$ and $\lambda\ne0$, let $v$ be a vertex of maximal depth with $\lambda_v\ne0$.
For a child $c$ of $v$, the only vectors $z_w$ with $z_w(c)\ne0$ are $z_v$, $z_c$ and the $z_w$ for the children $w$ of
$c$, and $\lambda_c=\lambda_w=0$ because $c$ and $w$ are deeper than $v$; so $y_c=\lambda_vz_v(c)\ne0$, $c\in U$, and
$v\in D_U$. Since the vectors $z_v$ are independent by (i), each element of $C_B^\perp\cap F^U$ has a unique coefficient
vector $\lambda$, and $\lambda\mapsto\lambda|_{D_U}$ is injective on these coefficient vectors; so
$\dim(C_B^\perp\cap F^U)\le\card{D_U}$. The children sets of distinct vertices are disjoint subsets of size at least
$d-1$, so $\card{D_U}\le\lfloor\card U/(d-1)\rfloor$.

(iii) For $d\ge4$, (ii) gives $r(U)\ge\card U-\lfloor\card U/3\rfloor$, which is at least $\lfloor\card U/2\rfloor+1$:
check $\card U\le5$ directly, and for $\card U\ge6$ use $\card U-\lfloor\card U/3\rfloor\ge2\card U/3\ge\card U/2+1$.
The second claim follows as in \cref{lem:ana-mincex}(c).
\end{informalproof}

\paragraph{Why it matters.} In a $d$-regular graph of girth at least $2r+2$, every ball of radius $r$ is a copy of $B$.
By \cref{prop:L}, inside such a ball the constraints that a smallest counterexample satisfies, namely (b) and (c) of
\cref{lem:ana-mincex} for subsets of the ball and the recoverability of every interior vertex from its neighbours, are
all met by a code of rate close to $(d-2)/(d-1)$, which exceeds $\frac12$ when $d\ge4$. So an argument that uses only
these local necessary conditions inside tree-like regions cannot give even the asymptotic bound $\frac12+o(1)$ at degree
at least~$4$. At degree~$3$ the ball code $C_B$ has rate $\frac12+1/\card B$, so such conditions are compatible with the
asymptotic bound there; the exact floor bound needs global information even at degree~$3$, as the closing of the maximal
chain in the proof of Theorem~S shows. For such arguments, large girth is the hardest case, not the easiest. And the
degree-$4$ statement over $\F_2$ on $4$-regular graphs of large girth is the scalar-linear, characteristic-$2$ case of the
undirected crux at degree~$4$, for which no characteristic-$p$ analogue of Theorem~A is known.

\begin{remark}[Computational evidence at degree four]\label{rem:ana-deg4}
Exhaustive exact searches with SAT solvers, not independently verified and not used in any proof, found no
counterexample to the floor bound for triangle-free graphs of maximum degree $4$: over $\F_2$ none of the
$193{,}900$ connected $4$-regular triangle-free graphs on $17$ vertices has a linear storage code of dimension $9$; over
$\F_3$ none of the connected $4$-regular triangle-free graphs on $13$ and $15$ vertices, and none of the non-bipartite
connected triangle-free graphs with degrees between $2$ and $4$ on $9$ vertices, or on $11$ vertices with a vertex of
degree $4$, has a code of dimension larger than its matching number.
\end{remark}

Several known codes that beat $\frac12$ on graphs of bounded degree, for instance on the Clebsch graph over $\Q$ and on
the Hoffman--Singleton graph, come from eigenvalues. For a graph $G$ with adjacency matrix $A$ and a real number
$\theta\ne0$, the matrix $I-A/\theta$ fits $G$, has coefficients of modulus $1/\abs\theta$, and its kernel, the \emph{eigenvalue code} of $\theta$, is the
$\theta$-eigenspace of $A$, of dimension $\mathrm{mult}(\theta)$ (zero if $\theta$ is not an eigenvalue). The next
proposition, a linear-programming bound in the spirit of Nozaki's bounds for regular graphs~\cite{Nozaki2015}, shows
where this mechanism stops. Let
\[
  p_0=1,\quad p_1(x)=\frac{x}{\sqrt d},\quad p_2(x)=\frac{x^2-d}{\sqrt{d(d-1)}},\quad
  p_3(x)=\frac{x(x^2-2d+1)}{\sqrt d\,(d-1)},\qquad K_n(\theta)=\sum_{k=0}^np_k(\theta)^2 ;
\]
these are the first orthonormal polynomials of the Kesten--McKay law, the spectral measure of the $d$-regular tree at a
vertex~\cite{Kesten1959,McKay1981}.

\begin{proposition}[Eigenvalue codes and girth]\label{prop:E}
\notformalized
Let $d\ge3$ and let $G$ be a $d$-regular graph on $N$ vertices with girth at least $g$.
\begin{enumerate}[label=(\roman*)]
\item For $n\in\{1,2,3\}$ with $2n+1\le g$ and every real $\theta$, $\mathrm{mult}(\theta)\le N/K_n(\theta)$.
\item $K_1(\theta)<2$ if and only if $\theta^2<d$; and $K_2(\theta)<2$ if and only if $1<\theta^2<d$.
\item $K_3(\theta)>2$ for every real $\theta$. Hence, if $g\ge7$, every eigenvalue has multiplicity less than $N/2$.
\item If $g\ge7$, no eigenvalue code of $G$ has dimension $N/2$ or more. If $g\ge5$ and an eigenvalue code has dimension
  larger than $N/2$, then $\theta$ is an integer with $2\le\abs\theta<\sqrt d$, so $d\ge5$.
\end{enumerate}
Equality holds in (i) for the Petersen graph ($d=3$, $g=5$, $n=2$, $\theta=1$, multiplicity $5$) and for the
Hoffman--Singleton graph ($d=7$, $g=5$, $n=2$, $\theta=2$, multiplicity $28$).
\end{proposition}

\begin{informalproof}
\emph{Orthonormality.} Let $A_T$ be the adjacency operator of the $d$-regular tree, $o$ its root, $S_k$ the sphere of
radius $k$ around $o$, and $e_k=\card{S_k}^{-1/2}\mathbf 1_{S_k}$. Counting neighbours in adjacent spheres gives
$A_Te_0=\sqrt d\,e_1$, $A_Te_1=\sqrt d\,e_0+\sqrt{d-1}\,e_2$ and $A_Te_k=\sqrt{d-1}\,(e_{k-1}+e_{k+1})$ for $k\ge2$. A
direct computation from these relations gives $p_k(A_T)e_0=e_k$ for $k\le3$; for instance
$(A_T^2-d)e_0=\sqrt{d(d-1)}\,e_2$. Hence $\langle p_j(A_T)e_0,p_k(A_T)e_0\rangle=\delta_{jk}$ for $j,k\le3$.

\emph{Moments.} For $k<g$ and every vertex $v$ of $G$, $(A^k)_{vv}=\langle A_T^ke_0,e_0\rangle$. Indeed, let
$\pi\colon T\to G$ be a covering of the component of $v$ by the $d$-regular tree with $\pi(o)=v$. Walks from $v$ lift
uniquely to walks from $o$ of the same length, and a lift of a closed walk of length $k<g$ is closed: otherwise the path
in $T$ between its distinct endpoints would project to a closed non-backtracking walk of positive length at most $k$,
which contains a cycle of length at most $k<g$. Conversely, closed walks from $o$ project to closed walks from $v$.

(i) Let $q=\sum_{k\le n}c_kp_k$ with $c_k=p_k(\theta)/K_n(\theta)$, so that $q(\theta)=1$. The symmetric matrix $q(A)^2$
has eigenvalue $1$ on the $\theta$-eigenspace and nonnegative eigenvalues elsewhere, so
$\mathrm{mult}(\theta)\le\operatorname{tr}q(A)^2$. Since $q^2$ has degree $2n<g$, the moments give
$\operatorname{tr}q(A)^2=N\lVert q(A_T)e_0\rVert^2=N\sum_kc_k^2=N/K_n(\theta)$.

(ii) $K_1(\theta)=1+\theta^2/d$. With $s=\theta^2$, $d(d-1)\bigl(K_2(\theta)-2\bigr)=(d-1)s+(s-d)^2-d(d-1)=(s-1)(s-d)$.

(iii) With $s=\theta^2$, $d(d-1)^2\bigl(K_3(\theta)-2\bigr)=f(s)=s^3-(3d-1)s^2+(3d^2-4d+2)s+d(d-1)$, and
$f(d+y)=y^3+y^2-2(d-1)y+d(d-1)^2$. The derivative $3y^2+2y-2(d-1)$ has roots $y_\pm=(-1\pm\sqrt{6d-5})/3$, a local
maximum at $y_-<0$ and a local minimum at $y_+>0$. On $[-d,0]$ the minimum of $f(d+y)$ is therefore attained at an
endpoint, where the values are $f(0)=d(d-1)>0$ and $f(d)=d(d-1)^2>0$. On $[0,\infty)$ the minimum is
$f(d+y_+)=d(d-1)^2-y_+^2(2y_++1)$, using $3y_+^2+2y_+=2(d-1)$. From $y_+<\sqrt{6d}/3$ we get
$y_+^2(2y_++1)<\frac{2d}3\bigl(\frac{2\sqrt{6d}}3+1\bigr)$, which is less than $d(d-1)^2$ because
$(d-1)^2>\frac49\sqrt{6d}+\frac23$ for $d\ge3$ (at $d=3$, $4>2.56$, and the left side grows faster). So $f(s)>0$ for
all $s\ge0$. The second claim follows from (i) with $n=3$.

(iv) The first claim is (iii). For the second: $A$ has integer entries, so the roots of the minimal polynomial of
$\theta$ over $\Q$ are eigenvalues of the same multiplicity; if this multiplicity exceeds $N/2$, the minimal polynomial
has degree $1$, and $\theta$, an algebraic integer, is an integer. By (i) with $n=2$ and (ii), $1<\theta^2<d$, so
$\abs\theta\ge2$ and $d>4$.

\emph{Equality cases.} $K_2(1)=1+\frac13+\frac46=2$ for $d=3$, and $K_2(2)=1+\frac47+\frac9{42}=\frac{25}{14}$ for
$d=7$, with $50\cdot\frac{14}{25}=28$; the multiplicities are those of the Petersen graph (spectrum $3$, $1^5$,
$(-2)^4$) and of the Hoffman--Singleton graph (computed in the proof of \cref{prop:HS}).
\end{informalproof}

\paragraph{Why it matters.} \Cref{prop:E} bounds the girth up to which the spectral mechanism can work: real eigenvalue codes,
which give the Clebsch graph a code of dimension $10$ over $\Q$ (eigenvalue $1$) and the Hoffman--Singleton graph one of
dimension $28$ (\cref{prop:HS}), never reach rate $\frac12$ at girth $7$ or more, at any degree, and beat it at girth $5$
or $6$ only when $d\ge5$; whether they beat it at girth~$6$ at all is open. It is a finite-girth counterpart of
Theorem~A for eigenvalue codes, with explicit thresholds.
It does not see $p$-torsion: over $\F_p$ the rank of $A-\theta I$ can be smaller than over $\Q$, and over $\F_2$ a code
need not come from an eigenvalue at all, and by \cref{cor:torsion} the finite-field gains that matter here are torsion
effects.

\begin{proposition}[Proposition HS: Hoffman--Singleton beats one half outside characteristic two]\label{prop:HS}
\notformalized
Let $A$ be the adjacency matrix of the Hoffman--Singleton graph ($50$ vertices, $7$-regular, girth $5$). Over every field
$F$ of characteristic different from $2$, the matrix $I-A/2$ fits the graph and has rank at most $22$, with equality if
the characteristic is $0$. So $\ker(I-A/2)$ is a linear storage code over $F$ of dimension at least $28>25$, with
recovery coefficients $\frac12$, and disjoint unions of copies of the graph form a $7$-regular family of girth $5$ that
carries such codes of rate at least $0.56$.
\end{proposition}

\begin{informalproof}
By~\cite{HoffmanSingleton1960} the Hoffman--Singleton graph is strongly regular with parameters $(50,7,0,1)$:
adjacent vertices have no common neighbour and non-adjacent vertices exactly one. Hence $A^2=7I+(J-I-A)$, where $J$ is
the all-ones matrix, that is, $(A-2I)(A+3I)=J$. The all-ones vector is an eigenvector with eigenvalue $7$, and $J$
vanishes on its orthogonal complement, which is invariant under $A$; so the other eigenvalues are $2$ and $-3$, with
multiplicities $f+g=49$ and $7+2f-3g=\operatorname{tr}A=0$, and $f=28$. Thus
$\rank_\Q(2I-A)=22$. Minors of an integer matrix that vanish over $\Z$ vanish modulo $p$, so
$\rank_{\F_p}(2I-A)\le22$, and $I-A/2=(2I-A)/2$ has the same rank whenever $2$ is invertible. The rank over a field
depends only on its characteristic.
\end{informalproof}

\begin{remark}[Computations on the Hoffman--Singleton graph]\label{rem:ana-hs}
By exact computation (not used in any proof), the rank of $2I-A$ is $22$ over $\F_3$ and $\F_7$ but $21$ over $\F_5$,
an example of the $p$-torsion that \cref{prop:E} does not see. Over $\F_2$ the matrix $I+A$ has rank $29$, giving only
dimension $21$, and we do not know whether the Hoffman--Singleton graph carries an $\F_2$-linear storage code of
dimension larger than $25$.
\end{remark}

\paragraph{Why it matters.} Theorem~S needs maximum degree~$3$ (and no triangles), and Theorem~A needs local
bipartiteness. One might hope
that large girth alone gives the bound $\frac12$ at every bounded degree, over every field. \Cref{prop:HS} shows that
girth~$5$ does not suffice: at degree~$7$, over $\Q$ and over every field of odd characteristic, there are linear codes
of rate $0.56$ with bounded coefficients. By \cref{prop:E}, eigenvalue codes need girth at most~$6$ to beat $\frac12$;
whether, for each $d$, some girth forces rate at most $\frac12+o(1)$ for all linear codes over a fixed finite field is
open, and in characteristic~$2$ even the Hoffman--Singleton graph is undecided.

Taken together, on side information of bounded degree $d\ge4$ that is locally bipartite, an $\F_p$-linear code with a
fixed number of symbols per vertex that beats $\frac12$ by a constant must have $p$-torsion of linear rank in every
bounded integer lift of its check matrix (\cref{cor:torsion}); in characteristic~$2$ it gains nothing from the identity
$(I+W)^2=0$ once the girth is at least $5$ (\cref{lem:ana-matching}); it does not come from a real eigenvalue code once
the girth is at least $7$ (\cref{prop:E}); and the local necessary conditions of a smallest counterexample cannot refute
its existence inside tree-like regions (\cref{prop:L}). These are design
constraints for constructions, and they say where a characteristic-$p$ proof of the crux would have to work.

\subsection{Targets, Mycielskians and Question G}\label{sec:ana-targets}

By Theorem~H (\cref{thm:H}), a homomorphism from side information to a good graph $T$ certifies the rate bound $\frac12$
for every code on the side information, linear or not. If the homomorphism is defined only after the deletion of a
vertex set $U$, conditioning on the symbols at $U$ weakens the bound by at most $\card U/2$ symbols: for a valid
polymatroid~$h$, the function $S\mapsto h(S\cup U)-h(U)$ on subsets of $V\setminus U$ is a valid polymatroid for $G-U$,
so $h(V)\le h(U)+\frac12\sum_{v\notin U}h(\{v\})\le\frac12\sum_vh(\{v\})+\frac12\sum_{u\in U}h(\{u\})$. A single good target
that receives the side information of the crux after the deletion of $o(N)$ vertices would therefore prove the crux for
all codes. Two facts constrain the targets. Good graphs are triangle-free: the triangle carries the code
$\{(a,b,a+b)\}$ of rate $\frac23$, and a graph that contains a triangle receives it (\cref{thm:H}). And homomorphisms do not decrease
the fractional chromatic number $\chi_f$, the least total weight of a family of weighted independent sets that covers
every vertex with weight at least $1$, because preimages of independent sets are independent. So side information $H$
with a deleted set $U$ needs a target with $\chi_f(T)\ge\chi_f(H-U)\ge(\card{V(H)}-\card U)/\alpha(H)$, where $\alpha$
is the independence number.

The classical way to raise $\chi_f$ while keeping a graph triangle-free is Mycielski's construction. For a graph $G$ on
$V$, the \emph{Mycielskian} $\mathrm M(G)$ has vertex set $V\cup V'\cup\{z\}$, where $V'=\{v':v\in V\}$, and edges $uv$,
$uv'$ and $u'v$ for $uv\in E(G)$, and $v'z$ for $v\in V$. Thus $\mathrm M(K_2)$ is the $5$-cycle and $\mathrm M^2(K_2)$
is the Gr\"otzsch graph. The Mycielskian of a triangle-free graph is triangle-free~\cite{Mycielski1955}, and
$\chi_f(\mathrm M(G))=\chi_f(G)+1/\chi_f(G)$~\cite{LPU1995}, so $\chi_f(\mathrm M^k(K_2))\to\infty$.

\begin{lemma}[Homomorphisms into a Mycielskian]\label{lem:ana-myc}
\notformalized
For graphs $H$ and $G$, $H$ admits a homomorphism to $\mathrm M(G)$ if and only if $H$ has an independent set $Z$ such that
the set $N_H(Z)$ of vertices outside $Z$ with a neighbour in $Z$ is independent and $H-Z$ admits a homomorphism to $G$.
\end{lemma}

The `if' direction is the case $t=1$ of~\cite[Lemma~2.3]{GHW2025}; we include the short proof of both directions.

\begin{informalproof}
Let $\varphi\colon H\to\mathrm M(G)$, and put $Z=\varphi^{-1}(z)$ and $Y=\varphi^{-1}(V')$. Both are independent,
because $z$ has no loop and $V'$ is independent in $\mathrm M(G)$, and $N_H(Z)\subseteq Y$ because $z$ is adjacent only
to $V'$. Define $\psi$ on $H-Z$ by $\psi(w)=\varphi(w)$ if $\varphi(w)\in V$ and $\psi(w)=v$ if $\varphi(w)=v'$. An edge
of $H-Z$ has at most one end in $Y$; if it has none, $\varphi$ maps it to an edge of $G$; if it has one, $\varphi$ maps it
to an edge $uv'$ with $uv\in E(G)$, and $\psi$ maps it to $uv$. So $\psi$ is a homomorphism to $G$. Conversely, given
$Z$ and $\psi\colon H-Z\to G$, map $Z$ to $z$, each $w\in N_H(Z)$ to $\psi(w)'$, and every other vertex $w$ to $\psi(w)$.
There are no edges inside $Z$ or inside $N_H(Z)$, and none between $Z$ and the vertices outside $Z\cup N_H(Z)$; edges
between $Z$ and $N_H(Z)$ go to edges $zv'$; and every other edge $ww''$ goes to $\psi(w)\psi(w'')$, $\psi(w)'\psi(w'')$ or
$\psi(w)\psi(w'')'$, which are edges of $\mathrm M(G)$ because $\psi(w)\psi(w'')\in E(G)$.
\end{informalproof}

\begin{corollary}[Odd cycle transversals]\label{cor:ana-oct}
\notformalized
If $H$ is triangle-free and $H-X$ is bipartite for a set $X$ of $k$ vertices, then $H$ admits a homomorphism to
$\mathrm M^k(K_2)$.
\end{corollary}

\begin{informalproof}
Induction on $k$; a bipartite graph maps to $K_2$. For $k\ge1$ pick $x\in X$. The set $Z=\{x\}$ is independent, $N_H(x)$
is independent because $H$ is triangle-free, and $H-x$ is triangle-free and becomes bipartite after the deletion of
$X\setminus\{x\}$, so $H-x\to\mathrm M^{k-1}(K_2)$ by induction. \Cref{lem:ana-myc} gives $H\to\mathrm M^k(K_2)$. (A
similar iteration is used in~\cite{GHW2025}.)
\end{informalproof}

\begin{proposition}[Mycielskians leave the good graphs]\label{prop:M}
\notformalized
\begin{enumerate}[label=(\roman*)]
\item The Clebsch graph admits a homomorphism to $\mathrm M^6(K_2)$. Hence $\mathrm M^6(K_2)$ is not good.
\item $K_2$ and $\mathrm M(K_2)$, the $5$-cycle, are good. Hence there is $j\in\{1,\dots,5\}$ such that $\mathrm M^j(K_2)$
  is good and $\mathrm M^{j+1}(K_2)$ is not: the Mycielski construction does not preserve goodness.
\item No good graph receives every triangle-free graph of maximum degree $5$. More generally, if $T_1,T_2,\dots$ are
  graphs and every finite triangle-free graph admits a homomorphism to some $T_k$, then some $T_k$ is not good; by
  \cref{cor:ana-oct} this applies to the graphs $\mathrm M^k(K_2)$, which are universal for triangle-free graphs by
  a theorem of Cropper, Gyárfás and Lehel~\cite{CGL2006}.
\end{enumerate}
\end{proposition}

\begin{informalproof}
(i) The Clebsch graph $\Gamma$ is triangle-free and $5$-regular on $16$ vertices. For an edge $uv$, the neighbourhoods
$N(u)$ and $N(v)$ are disjoint independent sets (as $\Gamma$ is triangle-free), so $\Gamma[N(u)\cup N(v)]$ is bipartite
on $10$ vertices and the remaining set $X$ of $6$ vertices has $\Gamma-X$ bipartite. By \cref{cor:ana-oct},
$\Gamma\to\mathrm M^6(K_2)$. The Clebsch graph is not good: by \cref{cor:H4} it has guessing number $10>8$, while a good
graph on $16$ vertices has guessing number at most $8$ by \cref{thm:entropic-rate}. By \cref{thm:H}, a graph that
receives a graph that is not good is not good. (ii) The $5$-cycle is good (\cref{prop:C5}), and $K_2$ maps to it, so
$K_2$ is good by \cref{thm:H}. With (i), the least $k$ for which $\mathrm M^k(K_2)$ is not good satisfies $2\le k\le6$;
take $j=k-1$. (iii) The Clebsch graph is triangle-free of maximum degree $5$, and a graph that receives it is not good.
\end{informalproof}

\begin{remark}[Which iterate fails; computational]\label{rem:ana-myc}
The Gr\"otzsch graph $\mathrm M^2(K_2)$ is good by an exact certificate, a nonnegative rational combination of
elemental Shannon inequalities found by computer and checked in exact arithmetic; the certificate is not published here,
not formalized and not used in any proof. With it,
$j\in\{2,3,4,5\}$ in \cref{prop:M}(ii). A SAT computation (not used in any proof) indicates that the Clebsch graph does
not map to $\mathrm M^5(K_2)$, so the argument of (i) cannot locate the first iterate that is not good; which iterate it
is remains open. At the level of Shannon inequalities the construction may fail at the first step: numerically, by a
floating-point linear program without an exact certificate, the Mycielskian of the circular clique $K_{8/3}$ (which is
good by \cref{prop:targets}) has Shannon bound about $7/13>\frac12$.
\end{remark}

\paragraph{Why it matters.} Mycielski's construction is the standard way to build triangle-free graphs of large
fractional chromatic number, and its iterates are universal for triangle-free graphs~\cite{CGL2006}; \cref{cor:ana-oct} refines this, bounding
the number of iterations by the size of an odd cycle transversal.
\Cref{prop:M} shows that this universality forces them out of the good graphs. Targets of growing $\chi_f$ for the
homomorphism route to the crux must therefore be built differently, and they must avoid receiving the Clebsch graph and
the other graphs known not to be good.

\begin{remark}[Good graphs with $\chi_f>3$; computational]\label{rem:ana-circulants}
Write $C_n(S)$ for the Cayley graph of $\Z_n$ with connection set $\pm S$; for vertex-transitive graphs
$\chi_f=N/\alpha$. Exact rational certificates, nonnegative combinations of elemental Shannon inequalities found by
computer and checked in exact arithmetic, show that the triangle-free circulants
$C_{13}(1,5)$ ($4$-regular, $\chi_f=13/4$), $C_{16}(1,3,8)$ ($5$-regular, $\chi_f=16/5$) and $C_{20}(1,6,10)$
($5$-regular, $\alpha=6$, $\chi_f=10/3$) are good. The certificate for $C_{20}(1,6,10)$ consists of $777$ inequalities;
$10/3$ is the largest value of $\chi_f$ among good graphs known to us. The certificates are not published here, and none
of this is formalized or used in any proof. These examples show that $\chi_f\le3$ is no barrier, and that $\chi_f$ does not
determine goodness: $C_{16}(1,3,8)$ has the same number of vertices, degree, independence number and fractional
chromatic number as the Clebsch graph, which is not good. Nearby circulants are not good: for instance
$C_{20}(2,5,6)$, with $\chi_f=10/3$, carries the $\F_2$-linear code $\ker(I+A)$ of dimension $12>10$.
\end{remark}

\begin{remark}[Cayley graphs of elementary abelian $2$-groups; partly computational]\label{rem:ana-z2}
Proof sketch only: every non-bipartite Cayley graph $\Gamma=\mathrm{Cay}(\Z_2^k,S)$ is not good. The graph $\Gamma$ is
bipartite if and only if some linear functional equals $1$ on $S$, and by linear duality over $\F_2$ this fails exactly
when some subset of $S$ with an odd number of elements sums to $0$. A minimal such subset $S'$ consists of $m-1$ linearly independent
elements and their sum, with $m$ odd (and $m\ge5$ if $\Gamma$ is triangle-free), so $\mathrm{Cay}(\langle
S'\rangle,S')$ is the folded $m$-cube, a subgraph of $\Gamma$. In the group algebra the matrix $I+A$ of a folded cube
satisfies $(I+A)^2=0$ over $\F_2$, so its kernel has dimension at least half; exact computation shows that the dimension
is $(2^{m-1}+2^{(m-1)/2})/2$ for $m\le11$, and we have not proved this formula in general. Where it holds, the folded
$m$-cube is not good, and neither is $\Gamma$, by \cref{thm:H}. The argument is complete for $k\le10$, where $m\le11$;
it would exclude, as targets, all non-bipartite triangle-free Cayley graphs of elementary abelian $2$-groups, a family that contains
the Clebsch graph and the folded cubes.
\end{remark}

These examples lead to the following question.

\begin{question}[Question G: the fractional chromatic number of good graphs]\label{prob:G}
Is $\chi_f$ unbounded on good graphs? That is, is $\sup\{\chi_f(T):T\text{ a good graph}\}$ infinite?
\end{question}

Good graphs are triangle-free, and triangle-free graphs have unbounded $\chi_f$; the question is whether goodness bounds
it. The analogous question for digraphs, with the fractional dichromatic number (fractional covers by vertex sets
inducing acyclic subgraphs), is also open.

\paragraph{Why it matters.} A negative answer, $\chi_f(T)\le C$ for every good graph $T$, would rule out the homomorphism
route for every family of side information whose underlying graphs keep fractional chromatic number above $C$ after the
deletion of $o(N)$ vertices, for instance families whose independence ratio stays below $1/C$. A positive answer is
necessary for the route but not sufficient, since the targets must also receive the side information. With the methods
available to us, certifying a single target by the Shannon linear program is limited to about $20$ vertices, while a
triangle-free graph with large $\chi_f$ must be large, so Question~G needs a structural argument in either direction.

\section{What a linear-time transposer would need}\label{sec:construction}

Our aim is to decide~(T), and we would prefer a construction. This section collects what the results above require of
any linear-time transposer, and of any oblivious one in particular. It contains no new theorems: every claim is one of
the results above or the first two papers, a citation, or explicitly labelled as a heuristic.

\subsection{Necessary conditions}

A machine witnessing~(T) must have all of the following properties, in addition to those of the first two
papers~\cite{FirstPaper2026,SecondPaper2026} ($\Omega(\log d)$ passes, folding macroscopic motion, and a coding gain of order
$\log d$ over routing on the block network of a slightly modified machine; within the abstract local-coding models of
the second paper, a fixed fractional saving is also excluded when the interaction graph has treewidth $o(d)$).
\begin{enumerate}[label=(C\arabic*),leftmargin=3em]
\item \emph{Curved motion.} In every macroscopic limit, the work heads travel a total distance of at least
half a block unit, a linear number of cells, along motion that is not countably-piecewise affine (Theorem~Cq,
\cref{thm:Cq}). In every clock, that is, every reparametrization of time by a homeomorphism whose inverse preserves null
sets, and for every countable set of affine maps, at least one block unit of travel has first returns that are not given
by these maps (Theorem~Cr, \cref{thm:Cr}); so neither a common change of clock nor fixed changes of scale remove the
curvature (\cref{cor:mot-clock}). The curved motion must be head-dependent, and
at least two work-head limits must be non-monotone and must not all have the same level sets (\cref{cor:mot-design}). Curved monotone expanders of the Bourgain--Yehudayoff type~\cite{BourgainYehudayoff2013} remain
possible.
\item \emph{Coding on bounded-depth side information, or depth-robust dependencies.} For an oblivious machine, Theorem~D
(\cref{thm:D}) leaves two ways out. Either some genie of $\eps L$ bits cuts the time--space dependency DAG to bounded
depth, and then the side information of the cones admits index codes of broadcast length $\eps L$, defeating the crux;
or no such genie exists, and the walks of the work heads are depth-robust. Home-based heads always admit such a genie
(Theorem~B, \cref{thm:B}), so for them only the first way is open (Corollary~B, \cref{cor:B}), and by the trade-off
(\cref{prop:tradeoff}) the depth is at least about $\beta/\eps$.
\item \emph{Codes that beat the pairing, of a particular kind.} If, as in the hardest instances we know, the relevant side
information is sparse, symmetric and locally tree-like, a construction needs storage codes of rate at least
$\tfrac12+\delta$ on it, for a fixed $\delta>0$ (a heuristic reading of \cref{prob:undirected-crux}). The classes of codes
that our results exclude are listed in \cref{sec:con-fail}.
Over a finite field, with a fixed number of symbols per vertex, a gain must come from $p$-torsion (\cref{cor:torsion}).
The binary code of rate $5/8$ on the Clebsch graph (\cref{cor:H4}) is not a torsion effect: it is the reduction of a rational
eigenvalue code with coefficients~$1$ (\cref{rem:ana-clebsch}), a mechanism that \cref{cor:torsion} rules out on locally
bipartite families of bounded degree. In odd
characteristic at most one of the two signed versions $I\pm W$ of a code can gain (\cref{prop:oddchar}).
\item \emph{If the dependencies are depth-robust: the deep-nesting template.} The only candidate we know for depth-robust walk pairs is
the deep-nesting pair (\cref{def:dnp}). Proof sketches indicate that per-head band rules fail already on one deeply nested walk
(\cref{rem:dep-N,rem:dep-Nprime}), although one such walk alone is easy to cut (\cref{lem:dep-visit}); and with one level
of nesting, labels exact on one head pay a price per hop on the cover independent of the depth (Theorem~E$_1$,
\cref{thm:E1}). Whether the pairs are depth-robust as $m$ grows is open (\cref{conj:Aprime}).
\end{enumerate}

\subsection{An architecture compatible with these conditions}

We describe, as a design and not as a result, how the first paper's geometric recursion~\cite[\S13]{FirstPaper2026} would
have to look in the light of (C1)--(C4). The recursion runs $\Theta(\log d)$ levels of geometrically decreasing volume,
each with $O(1)$ passes of an encoder and a decoder. Condition~(C1) forces the heads of each level to follow curved
rulers, for example near-identity Möbius maps realizing a monotone expander, and the second paper excludes the affine
variants~\cite{SecondPaper2026}. The encoder of each level must then save a constant fraction of the data by coding across
the expander's relations. By (C3), the code it uses must be nonlinear, or linear over a finite field,
with a $p$-torsion gain if it has a fixed number of symbols per vertex, and
a scalar linear code needs degree at least~$4$; its side-information graph must not map to any good target. Codes that beat $\tfrac12$ exist on single graphs with short
odd cycles, such as the Clebsch graph (girth~$4$) and, outside characteristic~$2$, the Hoffman--Singleton graph
(girth~$5$, \cref{prop:HS}). We do not know any family of bounded degree and growing odd girth that carries one; such a
family would answer \cref{prob:undirected-crux} in the negative. For an oblivious design, Corollary~B applies if its
work heads are all home-based, and the design then stands or falls with the crux; otherwise, if the crux holds, its
dependency structure must be depth-robust, and the deep-nesting pair is the template.

\subsection{Mechanisms that fail}\label{sec:con-fail}

The following classes of constructions fail for reasons proved here, in addition to those listed in the first two papers.
\begin{itemize}[leftmargin=2em]
\item \emph{Home-based motion of all work heads with legs of boundedly many constant-rate runs}, obliviously or
adaptively, with any number of legs at any scales (Theorem~C).
\item \emph{One expanding tape, affine returns, a common change of clock, or fixed changes of scale}: no single work tape
can carry the expansion, returns given by countably many affine maps in a common clock cannot carry it either, and
neither a common reparametrization of time nor a fixed reparametrization of each tape adds anything
(\cref{cor:mot-one-tape,thm:Cr,cor:mot-clock}).
\item \emph{Stroke trees for all work heads}, the recursive schedules of \cref{def:mot-stroke}: their limits are affine
wherever the heads move (\cref{prop:mot-strokes}).
\item \emph{Linear codes with bounded real or complex coefficients}, scalar or with a bounded number of symbols per
vertex, on locally bipartite side information of bounded degree (Theorem~A and \cref{prop:ana-A-ext}),
\emph{scalar linear codes of degree~$3$} on triangle-free side information (Theorem~S), and \emph{all codes} on
side information that is $C_5$-, $K_{14/5}$- or Petersen-colourable (\cref{cor:H2}).
\end{itemize}

\section{Roadmap and open problems}\label{sec:open}

\subsection{Where the decision of (T) stands}

For oblivious machines, Theorem~D reduces $\neg$(T) to two statements, and this paper narrows both; the last row of
the table records the unconditional results on motion.
\begin{center}
\begin{tabular}{@{}p{0.27\textwidth}p{0.33\textwidth}p{0.33\textwidth}@{}}
\toprule
 & \textbf{proved here} & \textbf{open} \\
\midrule
Depth reduction (TCP) &
one walk (by planarity) and all couplings of home-based walks (Theorem~B); packing certificates cannot show
depth-robustness (\cref{prop:packing}); on the deep-nesting pair with one level of nesting, labels exact on one head pay
a fixed price (Theorems~E, E$_1$) &
TCP$(2)$ at fixed $\eps$; its sharpest instance is the deep-nesting pair (\cref{conj:Aprime}), where per-head band
rules appear to fail (sketch only, \cref{rem:dep-N,rem:dep-Nprime}) \\
\addlinespace
The crux (index coding on bounded-depth side information) &
undirected form, rate $\tfrac12$: bounded complex linear codes on locally bipartite graphs (Theorem~A); scalar linear codes
over every field at degree~$3$ (Theorem~S); all codes on $C_5$-, $K_{14/5}$- and Petersen-colourable graphs (\cref{cor:H2}) &
on locally bipartite graphs, scalar linear codes over finite fields at degree at least~$4$ (vector-linear ones
already at degree~$3$)
and nonlinear codes at degree at least~$3$; the directed, depth-bounded form \\
\addlinespace
Motion &
machines whose work heads are all home-based with legs of boundedly many constant-rate runs fail unconditionally
(Theorem~C); half a block unit of travel must be curved (Theorem~Cq) &
whether curved motion of the Bourgain--Yehudayoff type can be coupled to a coding scheme \\
\bottomrule
\end{tabular}
\end{center}
For home-based oblivious machines only the crux remains (Corollary~B), and only in a weak form: a rate $\rho(D)$ with
$\rho(D)\log D\to\infty$ suffices. Theorem~C also covers adaptive machines; for them the decomposition
needs further hypotheses, which we do not treat here.

\subsection{Open problems}

\begin{enumerate}[label=\arabic*.,leftmargin=2em]
\item \emph{Decide (T).} Construct a machine transposing $d\times d$ matrices with $\clog d$-bit entries in time
$O(d^2\log d)$, or prove that none exists. The first paper's frozen checker accepts either answer~\cite{FirstPaper2026}.

\item \emph{Depth reduction for interacting walks} (\cref{prob:tcp}). Does every coupling of $k$ walks admit an
$\eps T$-edge cut that leaves depth $D(k,\eps)$? A positive answer for every~$k$, with the crux, refutes oblivious~(T)
(Theorem~D). The first open case is $k=2$, and its sharpest instance is the deep-nesting pair. There, the missing step is
a robust bound, logarithmic in the depth, on how many levels a label of small cost can serve on both heads at once
(\cref{conj:Aprime}, stated for labels on a cyclic cover and for a suitable odd shift); the same statement for labels on the walk DAG would
make these pairs depth-robust as $m\to\infty$, give
$\TCP(2)$ a negative answer, and point constructions toward deep nesting.

\item \emph{The undirected crux} (\cref{prob:undirected-crux}). Do storage codes on locally bipartite graphs of bounded
degree have rate at most $\tfrac12+o(1)$? Open for scalar linear codes over finite fields at degree at least~$4$ (and for vector-linear ones already at
degree~$3$), where a gain would need $p$-torsion (\cref{cor:torsion}), and for nonlinear codes at every degree at
least~$3$. Arguments that use only the
local necessary conditions of a smallest counterexample inside tree-like regions cannot give rate $\tfrac12+o(1)$ at
degree at least~$4$ (\cref{prop:L}), and eigenvalue codes cannot reach rate~$\tfrac12$ at girth at least~$7$
(\cref{prop:E}).

\item \emph{Targets} (\cref{prob:G}). Is the fractional chromatic number unbounded on good graphs? Theorem~H
proves rate~$\tfrac12$ on every graph that maps to a good target, so a family of good targets with growing fractional
chromatic number would extend \cref{cor:H2} to side information that needs more colours; a negative answer, $\chi_f\le C$ on all good graphs, would
close this route for side information whose fractional chromatic number stays above~$C$ after the deletion of $o(N)$
vertices. Mycielski's construction does not preserve goodness (\cref{prop:M}).

\item \emph{Cubic graphs of large girth.} Do all storage codes on cubic graphs of large girth have rate at most~$\tfrac12$?
By \cref{cor:H3} this would follow from a positive answer to Nešetřil's Pentagon Problem~\cite{Nesetril1999}, or
from a proof that such graphs have circular chromatic number at most~$14/5$, which is open. For scalar linear codes, Theorem~S answers it.

\item \emph{The crux in full.} Prove an index-coding lower bound of any rate $\rho(D)\gg\beta/D$ for the side-information
systems of bounded-depth cones of oblivious schedules (\cref{def:crux}), linear or nonlinear; by \cref{prop:tradeoff},
no depth reduction can use a smaller rate. For home-based heads, where depth reduction is
proved, the rate $\rho(D)\gg1/\log D$ suffices (Corollary~B).

\item \emph{Formalization.} Formalize the results of
\cref{sec:decomposition,sec:motion,sec:depth,sec:ana-A,sec:ana-fields,sec:ana-limits,sec:ana-targets} that are not
formalized.
\end{enumerate}


\appendix
\section{Index of Lean declarations}\label{app:index}

Every declaration of this paper's library cited in the text, grouped by the Lean module that contains it, relative to the namespace
\texttt{MotionDepthCodes}. All of them depend only on the axioms \texttt{propext}, \texttt{Classical.choice} and
\texttt{Quot.sound} (\cref{app:verify}). The declarations of the first paper's library cited in \cref{sec:motion} are
verified there~\cite{FirstPaper2026}.

\paragraph{\texttt{MotionDepthCodes.CircularClique}}
{\small\raggedright \texttt{circularClique}, \texttt{circularCliqueHom}, \texttt{circularClique\_adj\_of\_lt}, \texttt{cycleToCircularClique}, \texttt{rateBound\_circularClique\_14\_5}, \texttt{rateBound\_circularClique\_of\_le}, \texttt{rateBound\_cycleGraph\_odd}.\par}

\paragraph{\texttt{MotionDepthCodes.Clebsch}}
{\small\raggedright \texttt{clebschCode}, \texttt{clebschCode\_card}, \texttt{clebschCode\_isStorageCode}, \texttt{clebschGraph}, \texttt{clebschGraph\_adj\_iff}, \texttt{exists\_clebsch\_storageCode}.\par}

\paragraph{\texttt{MotionDepthCodes.ClebschCertificate}}
{\small\raggedright \texttt{clebschCode\_eq}, \texttt{clebsch\_card\_le}, \texttt{clebsch\_guessingNumber}, \texttt{clebsch\_rateBound\_iff}, \texttt{finrank\_ker\_clebsch}, \texttt{mem\_ker\_clebsch\_iff}, \texttt{rank\_clebsch}, \texttt{rateBound\_clebsch}, \texttt{rateBound\_clebschGraph}.\par}

\paragraph{\texttt{MotionDepthCodes.DepthLabels}}
{\small\raggedright \texttt{card\_descents\_eq}, \texttt{card\_violated\_eq}, \texttt{congr\_of\_exact}, \texttt{cost\_ge\_of\_negative\_cycles}, \texttt{cycleDefect}, \texttt{cycle\_errors\_congr}, \texttt{le\_of\_modEq\_neg}, \texttt{lemma\_Zq}.\par}

\paragraph{\texttt{MotionDepthCodes.Pentagon}}
{\small\raggedright \texttt{nbhd\_cycleGraph\_five}, \texttt{pentagonNbhd}, \texttt{rateBound\_cycleGraph\_five}, \texttt{rateBound\_pentagon}.\par}

\paragraph{\texttt{MotionDepthCodes.Petersen}}
{\small\raggedright \texttt{petersenGraph}, \texttt{petersenGraph\_adj\_iff}, \texttt{petersenKneser\_card}, \texttt{petersenKneser\_injective}, \texttt{rateBound\_petersenGraph}.\par}

\paragraph{\texttt{MotionDepthCodes.Polymatroid}}
{\small\raggedright \texttt{IsGood}, \texttt{IsGood.of\_hom}, \texttt{IsPolymatroid}, \texttt{IsPolymatroid.eq\_of\_closure}, \texttt{IsPolymatroid.insert\_eq\_of\_subset}, \texttt{IsPolymatroid.pushRank}, \texttt{IsValid}, \texttt{IsValid.pushRank}, \texttt{RateBound}, \texttt{RateBound.mono}, \texttt{RateBound.of\_digraphHom}, \texttt{RateBound.of\_graphHom}, \texttt{RateBound.of\_hom}, \texttt{pushRank}, \texttt{pushRank\_univ}, \texttt{sum\_pushRank\_singleton\_le}.\par}

\paragraph{\texttt{MotionDepthCodes.RateHalf}}
{\small\raggedright \texttt{IsGood.card\_sq\_le}, \texttt{IsGood.codeEntropy\_le}, \texttt{card\_sq\_le\_of\_rateBound\_half}.\par}

\paragraph{\texttt{MotionDepthCodes.StorageCode}}
{\small\raggedright \texttt{IsRecoverable}, \texttt{IsStorageCode}, \texttt{codeEntropy}, \texttt{codeEntropy\_isPolymatroid}, \texttt{codeEntropy\_isValid}, \texttt{codeEntropy\_univ\_le\_of\_rateBound}, \texttt{codeEntropy\_univ\_uniform}, \texttt{log\_card\_le\_of\_rateBound}.\par}

\paragraph{\texttt{MotionDepthCodes.Subcubic}}
{\small\raggedright \texttt{IsLinearStorageCode}, \texttt{IsLinearStorageCode.card\_sq\_le}, \texttt{IsLinearStorageCode.exists\_coeffs}, \texttt{IsLinearStorageCode.exists\_pattern}, \texttt{IsLinearStorageCode.finrank\_le\_half}, \texttt{RecTriangleFree}, \texttt{Star}, \texttt{ThmSHyp}, \texttt{exists\_linearStorageCode\_of\_matching}, \texttt{exists\_pattern\_finrank\_ker\_ge\_of\_matching}, \texttt{finrank\_ker\_le\_barrier}, \texttt{finrank\_ker\_le\_half\_of\_degree\_le\_three}, \texttt{finrank\_ker\_le\_heavy}, \texttt{finrank\_subCode\_biUnion\_le}, \texttt{finrank\_subCode\_le\_half}, \texttt{finrank\_subCode\_union\_le}, \texttt{isLinearStorageCode\_of\_le\_ker}, \texttt{isStorageCode\_iff\_isLinearStorageCode}, \texttt{rank\_ge\_half\_of\_degree\_le\_three}, \texttt{star\_of\_degree\_le\_three}, \texttt{thmSHyp\_of\_cliqueFree}.\par}

\paragraph{\texttt{MotionDepthCodes.Symmetry}}
{\small\raggedright \texttt{IsAut}, \texttt{IsPolymatroid.symmetrize}, \texttt{IsValid.symmetrize}, \texttt{RateBound.of\_invariant}, \texttt{sum\_symmetrize\_singleton}, \texttt{symmetrize}, \texttt{symmetrize\_invariant}, \texttt{symmetrize\_univ}.\par}

\paragraph{Also named in the text}
{\small\raggedright \texttt{emod\_of\_bounds}, \texttt{nbhd}, \texttt{outNbhd}.\par}

\section{Building and verifying the formalization}\label{app:verify}

The Lean project is the directory \leanfile{lean/} next to this paper. It uses Lean and Mathlib, both at version v4.34.1
(\leanfile{lean-toolchain}, \leanfile{lakefile.toml}, \leanfile{lake-manifest.json}), and imports the library of the second
paper~\cite{SecondPaper2026}, which imports that of the first~\cite{FirstPaper2026}, through relative paths. Keep the
repository's directory layout, so that the sibling directories of the earlier papers stay in place.

\paragraph{Build.} With \texttt{elan} installed, run in \leanfile{lean/}:
\begin{center}
\texttt{lake exe cache get \&\& lake build}
\end{center}
The first command downloads the compiled Mathlib; the second compiles the library and the parts of the earlier libraries
that it uses. Before the first build, fetch the second paper's pinned external dependency as its README describes
(\texttt{python3 verification/fetch\_openai\_metric.py}, run in that paper's directory).

\paragraph{Axioms.} The file \leanfile{PaperAxioms.lean}, generated from the \LaTeX{} sources of this paper by
\leanfile{tools/lean\_index.py}, runs \texttt{\#print axioms} on every declaration of this library named here (\cref{app:index}):
\begin{center}
\texttt{lake env lean PaperAxioms.lean}
\end{center}
Every declaration depends at most on \texttt{propext}, \texttt{Classical.choice} and \texttt{Quot.sound}. No Lean file of this
library uses \texttt{sorry}, \texttt{admit} or \texttt{native\_decide}, or declares an \texttt{axiom}. The heartbeat limit is
raised once, for the Clebsch certificate, which is a long but routine computation.

\paragraph{Certificates.} The certificates of \cref{sec:targets,sec:clebsch} are Lean proofs generated by
\leanfile{tools/certificates.py} (which needs \texttt{numpy} and \texttt{scipy}). For example,
\begin{center}
\texttt{python3 tools/certificates.py clebsch 5 8 clebsch clebschNbhd --all-pairs}
\end{center}
prints the automorphisms and the proof used in \leanfile{ClebschCertificate.lean}. Different solver versions may produce
different certificates; any certificate printed is checked by Lean, and the files in the library are the record.

\paragraph{Where things are.} \leanfile{Polymatroid.lean} (polymatroids, closure, Theorem~H), \leanfile{StorageCode.lean}
(entropic polymatroids of codes), \leanfile{Symmetry.lean} (symmetrization), \leanfile{Pentagon.lean},
\leanfile{Petersen.lean} and \leanfile{CircularClique.lean} (good targets; the Bondy--Hell homomorphisms and odd cycles), \leanfile{RateHalf.lean} (Corollary~H2),
\leanfile{Clebsch.lean} and \leanfile{ClebschCertificate.lean} (Corollary~H4), \leanfile{Subcubic.lean} (Theorem~S, linear storage codes and the matching code),
\leanfile{DepthLabels.lean} (Lemma~Zq, the exact cases of the chain and transfer identities, and Theorem~E), and \leanfile{FinSums.lean} (sums over small index sets used by the
certificates).


\bibliographystyle{plainurl}
\bibliography{references}

\end{document}
